%% Converted to Elsevier CAS single-column format for Information and Software Technology.
\documentclass[a4paper,fleqn]{cas-sc}

%\usepackage[authoryear,longnamesfirst]{natbib}
\usepackage[numbers,sort&compress]{natbib}
\usepackage{graphicx}
\usepackage{booktabs}
\usepackage{longtable}
\usepackage{array}
\usepackage{calc}
\usepackage{amsmath}
\usepackage{amssymb}
\usepackage{url}
\usepackage{hyperref}
\usepackage{xcolor}
\usepackage{textcomp}
\usepackage{fancyvrb}
\usepackage{float}
\usepackage{capt-of}
\usepackage{listings}

\lstset{
  basicstyle=\ttfamily\small,
  breaklines=true,
  columns=fullflexible,
  keepspaces=true,
  frame=single,
  captionpos=b
}


\DefineVerbatimEnvironment{Highlighting}{Verbatim}{commandchars=\\\{\}}
\newenvironment{Shaded}{}{}
\newcommand{\NormalTok}[1]{#1}
\providecommand{\tightlist}{%
  \setlength{\itemsep}{0pt}\setlength{\parskip}{0pt}}

\begin{document}
\let\WriteBookmarks\relax
\renewcommand{\topfraction}{0.9}
\renewcommand{\bottomfraction}{0.8}
\renewcommand{\textfraction}{0.07}
\renewcommand{\floatpagefraction}{0.7}
\setcounter{topnumber}{3}
\setcounter{bottomnumber}{2}
\setcounter{totalnumber}{5}

\shorttitle{Prompt Framing Sensitivity in LLM-Based Coding Pipelines}
\shortauthors{Huynh et al.}

\title[mode = title]{Prompt Framing Sensitivity in LLM-Based Coding Pipelines and Reasoning Tasks}

\author[1]{Duc Toan Huynh}
\ead{duc.huynh@ue-germany.de}
\credit{Conceptualization, Methodology, Software, Investigation, Data curation, Writing - original draft}

\author[1]{Iftikhar Ahmed}[orcid=0000-0001-7863-3746]
\cormark[1]
\ead{iftikhar.ahmed@ue-germany.de}
\credit{Conceptualization, Methodology, Supervision, Writing - review and editing}

\author[1]{Talha Ali Khan}[orcid=0000-0001-8962-5602]
\ead{talhaali.khan@ue-germany.de}
\credit{Methodology, Formal analysis, Writing - review and editing}

\author[1]{Raja Hashim Ali}[orcid=0000-0002-9068-2337]
\ead{hashim.ali@ue-germany.de}
\credit{Methodology, Validation, Writing - review and editing}

\affiliation[1]{organization={Machine INtelligence \& Data Science (MINDS) Lab, University of Europe for Applied Sciences},
                city={Potsdam},
                country={Germany}}

\cortext[cor1]{Corresponding author}

\begin{abstract}
\textbf{Context:} Large Language Models (LLMs) are increasingly used in code-generation assistants, automated evaluation pipelines, and agentic tool workflows. In these settings, prompt-framing choices may affect not only correctness but also whether outputs can be parsed, executed, and validated.

\textbf{Objectives:} This study examines how tone, role framing, and prompt-wrapper format affect correctness and observable output behavior in automated LLM pipelines for Python code generation and final-answer mathematical reasoning.

\textbf{Methods:} We conducted a crossed $3 \times 2 \times 3$ factorial experiment spanning 18 prompt variants (three tones, two role-framing levels, and plaintext/JSON/XML input wrappers) across 50 base tasks per task kind, six local Ollama models, and two cloud API models ($14{,}400$ primary single-shot runs). To quantify run-to-run variability, real retry behavior, paraphrase sensitivity, cross-regime confounding, and external validity, we supplemented the primary matrix with $14{,}398$ independent $k = 3$ repeated generations across all eight models, $450$ live end-to-end Python workflows with real error-feedback retries, $180$ neutral paraphrase and length-matched controls, a 20-task shared-regime comparison across seven models, and a 15-task non-math software-engineering (SE) coding benchmark ($135$ runs).

\textbf{Results:} In three of four strata, model choice dominated correctness: model identity explained $7.2\%$--$41.0\%$ of Pass@1 variation versus at most $0.41\%$ for all 18 prompt variants combined (between-model spreads of $22.3$--$85.6$ percentage points), although the two near-equivalent cloud models did not differ on Python tasks. Prompt-framing effects on Pass@1 were modest and selective: 12 of 20 stratum-level mixed-effects odds-ratio $95\%$ intervals included $1.0$, and the 8 statistically supported effects had Average Marginal Effects of $1.0$--$5.2$ pp. Independent $k = 3$ repetitions confirmed run-to-run stability (overall pass-rate SD $0.00$--$2.08$ pp; within-cell flip rates $0.00\%$ on Gemini Flash to $28.33\%$ on Llama 3.1 math) and replicated both the local Python politeness penalty ($-2.56$ pp, $95\%$ CI $[-4.85, -0.27]$) and the local math JSON penalty ($-6.78$ pp, $[-9.53, -4.03]$). Operational behavior shifted consistently: polite phrasing inflated tokens ($+11.0\%$ to $+12.4\%$) and latency ($+7.5\%$ to $+10.8\%$) while reducing local code-only compliance ($-7.0$ pp), a pattern replicated on the non-math SE benchmark (up to $+24.5\%$ tokens and $+38.2\%$ latency with $0$ pp accuracy gain). In the live end-to-end pipeline, two real error-feedback retries recovered $+4$--$6$ pp on Gemini Flash ($68$--$70\% \to 74$--$76\%$, far below an independent-redraw simulation's $96.7\%$ upper bound) and $+8$--$14$ pp on local models at $2.5$--$2.7\times$ calls.

\textbf{Conclusion:} Rather than acting as a universal accuracy booster, zero-shot prompt framing primarily governs the operational interface (token cost, latency, and parser compliance) between an LLM and its downstream evaluator. In Section~\ref{sec:discussion}, we synthesize evidence-grounded starting defaults—clearly distinguished from hypotheses requiring local validation—and cost--benefit proportionality criteria for deciding when prompt benchmarking is warranted.
\end{abstract}

\begin{highlights}
\item Full-factorial $3 \times 2 \times 3$ prompt-framing evaluation ($14{,}400$ primary runs) plus $14{,}398$ independent $k = 3$ repetitions across all 8 models.
\item Model identity explains $7.2$--$41.0\%$ of Pass@1 variation vs.\ $\le 0.41\%$ for prompts in 3 of 4 strata; supported prompt AMEs are $1.0$--$5.2$ pp.
\item Polite framing acts mainly as operational overhead in code generation ($+12.4\%$ tokens, $+10.8\%$ latency, $-7.0$ pp compliance), replicated on a 15-task non-math SE benchmark.
\item Live end-to-end pipeline evaluation ($450$ workflows) shows real retries recover $+4$--$6$ pp on API ($T = 0$) and $+8$--$14$ pp locally, far below simulated independent-retry bounds.
\item Consolidates evidence-grounded starting defaults (distinct from hypotheses) and cost-proportionality criteria for practitioners, tool builders, and researchers.
\end{highlights}

\begin{keywords}
Large Language Models \sep Prompt Engineering \sep Code Generation \sep Automated Evaluation \sep Software Engineering
\end{keywords}

\maketitle

\section{Introduction}\label{sec:introduction}

    Large Language Models (LLMs) are widely embedded in software-engineering workflows, including code generation, test automation, and agentic tool pipelines \citep{ouyang2022instructgpt, chen2021codex, jimenez2024swebench}. Because outputs are consumed by parsers, compilers, and unit-test harnesses rather than only human readers, operational surface behavior—token volume, latency, and parser compliance—matters alongside functional correctness.

    Practitioners routinely adjust zero-shot prompt framing using politeness cues, persona roles, urgency pressure, or structured JSON/XML input wrappers \citep{li2023emotionprompt,yin2024politeness,liu2024structuredoutput,he2024formatting}. Yet prior evidence remains fragmented and context-dependent \citep{zhao2021calibrate, lu2021fantastically, sclar2023spurious, zhuo2024prosa, hua2025promptartifact}: a prompt convention that helps one model emit parseable code can degrade another model on compact mathematical reasoning. This creates both a reproducibility risk for empirical LLM research (confounding prompt choice with model capability) and a reliability risk for automated software pipelines (unintended verbosity, latency, or extraction failures).

    This study investigates \emph{selected zero-shot prompt-framing dimensions}—specifically tone (neutral, polite, aggressive), role framing (no role, expert persona), and input-wrapper serialization (plaintext, JSON, XML)—crossed in a $3 \times 2 \times 3$ matrix (18 variants). We do not claim that these three surface dimensions are the most influential prompting techniques overall (compared with few-shot exemplars, chain-of-thought reasoning, or iterative repair); rather, holding semantic content and output contracts fixed isolates the framing choices developers routinely vary in zero-shot templates. We evaluate both Python code generation (syntactic artifact extraction and unit testing) and final-answer mathematical reasoning (compact scalar extraction) across six local Ollama models and two cloud API models ($14{,}400$ primary runs), complemented by $14{,}398$ independent $k = 3$ repeated runs, $450$ live end-to-end retry workflows, neutral paraphrase/length controls, a shared-task cross-regime check, and a 15-task non-math SE coding benchmark.

    Our results show a consistent hierarchy: wherever candidate models differ in capability, model selection dominates Pass@1 correctness ($7.2\%$--$41.0\%$ of variance vs.\ $\le 0.41\%$ for prompts in three of four strata), whereas zero-shot prompt framing produces modest, selective accuracy shifts ($1.0$--$5.2$ pp across 8 supported contrasts; 12 of 20 contrasts compatible with zero) alongside systematic operational shifts in tokens ($-19.8\%$ to $+14.2\%$), latency ($-11.2\%$ to $+10.8\%$), and compliance ($-7.05$ to $+6.98$ pp). Rather than prescribing universal prompting rules, we consolidate stakeholder-specific starting defaults and cost--benefit proportionality criteria in Section~\ref{sec:discussion} (Table~\ref{tab:actionable-guidelines}).

    This paper makes four main contributions:
\begin{itemize}
\tightlist
\item \textbf{Crossed factorial and repeated-run evaluation:} A $3 \times 2 \times 3$ evaluation ($14{,}400$ primary runs across 8 models and two task contracts) backed by $14{,}398$ independent $k = 3$ repetitions, neutral paraphrase/length controls, a 20-task shared-regime comparison, and a 15-task non-math SE benchmark.
\item \textbf{Multi-dimensional operational and live-pipeline evaluation:} Joint modeling of Pass@1, tokens, latency, audited compliance, and failure modes \citep{liang2023helm, zhou2023ifeval}, paired with a live 4-stage end-to-end pipeline ($450$ workflows with real retries) contrasted against simulated offline replay.
\item \textbf{Cross-domain mechanism reversals and associational behavioral adjustment:} Evidence that polite phrasing acts mainly as token/latency overhead in code generation, whereas structured wrappers exhibit opposite-sign effects across API Python and local math.
\item \textbf{Evidence-grounded operational guidance:} Consolidated stakeholder recommendations and cost--benefit proportionality tiers in Section~\ref{sec:discussion}, explicitly separating directly supported findings from hypotheses requiring local validation.
\end{itemize}

\paragraph{Paper Organization.}
Section~\ref{sec:relatedwork} reviews related work. Section~\ref{experimental-design} defines RQ1--RQ4 and details the prompt matrix, models, tasks, execution provenance, and outcome measures. Section~\ref{statistical-and-analytical-strategy} describes our statistical strategy. Section~\ref{sec:results} presents RQ1--RQ4 results, robustness checks, the live end-to-end pipeline, and our $k = 3$ replication and external-validity studies. Section~\ref{sec:discussion} interprets the findings and consolidates practical guidelines; Section~\ref{threats-to-validity} discusses threats to validity and audits; Section~\ref{conclusion} concludes.


\section{Related Work}\label{sec:relatedwork}

\subsection{Prompt Sensitivity and Prompt Engineering in
LLMs}\label{prompt-sensitivity-and-prompt-engineering-in-llms}

Since large language models with strong few-shot and instruction-following abilities appeared in 2020, prompt engineering has become an essential part of working with these models \citep{brown2020fewshot}. Recently, prompting techniques have evolved from simple task descriptions to include instructions, examples, input/output format, role descriptions, reasoning cues, and conversational style. Recent surveys describe prompt engineering not as a collection of standard rules but as a broad and heterogeneous design space with many proposed techniques \citep{schulhoff2024promptreport}.

Across prior empirical studies, prompt variation has been investigated along several distinct dimensions: (i) \emph{exemplar-based prompting}, including few-shot selection and example ordering \citep{brown2020fewshot, zhao2021calibrate, lu2021fantastically}; (ii) \emph{reasoning elicitation}, such as zero-shot reasoning cues and chain-of-thought prompting \citep{wei2022cot, kojima2022zeroshot}; (iii) \emph{persona and role assignment}, which conditions the model on domain expertise \citep{schulhoff2024promptreport, zhuo2024prosa}; (iv) \emph{social and emotional tone}, including politeness and urgency/emotional stimuli \citep{li2023emotionprompt, yin2024politeness}; (v) \emph{structural serialization and delimiters}, such as spacing, markup tags, JSON/XML envelopes, and schema constraints \citep{sclar2023spurious, liu2024structuredoutput}; and (vi) \emph{iterative repair prompts}, which feed execution or validation errors back to the model in multi-turn workflows \citep{jimenez2024swebench}.

Early research on prompt-programming and few-shot calibration work also showed that prompts are not merely presentation layers but can significantly influence model behaviors. Reynolds and McDonell describe prompts as natural-language programs that communicate task intention and structure \citep{reynolds2021promptprogramming}. Zhao et al. demonstrated that prompt format, example choice, and example order can create a large variance in few-shot accuracy \citep{zhao2021calibrate}. Lu et al. found that prompt example permutations alone can change performance from near-random to near state-of-the-art \citep{lu2021fantastically}. More recent work extended these findings to superficial prompt features. Sclar et al. showed that minor formatting changes such as spacing and delimiter choice can cause significant performance shifts \citep{sclar2023spurious}. Zhu et al. evaluated model robustness under adversarial prompt perturbations \citep{zhu2023promptrobust}, and Zhuo et al. studied how prompt structure affects code-generation evaluation \citep{zhuo2024prosa}. Recently, Hua et al.\ \citep{hua2025promptartifact} examined whether prompt sensitivity should be interpreted as a model flaw or an evaluation artifact, showing that benchmark rankings can shift under meaning-preserving rephrasings.

Tone and social framing have also been studied as prompt dimensions. Studies on emotional stimuli have shown that emotion-oriented discourse can improve performance on selected benchmarks \citep{li2023emotionprompt}, while cross-lingual politeness studies indicate that polite versus impolite phrasing can shift both response style and accuracy \citep{yin2024politeness}. However, most of this prompt-sensitivity literature evaluates accuracy in isolation,whereas multi-metric frameworks such as HELM \citep{liang2023helm} and instruction-following benchmarks such as IFEval \citep{zhou2023ifeval} emphasize that operational deployment also depends on efficiency, latency, and verifiable format adherence.

\subsection{LLMs for Code Generation and Software-Engineering
Tasks}\label{llms-for-code-generation-and-software-engineering-tasks}

Code generation is a major software engineering application of LLMs. Benchmarks such as HumanEval and MBPP evaluate whether generated Python programs satisfy hidden or visible tests \citep{chen2021codex, austin2021mbpp}. For more realistic scenarios, SWE-bench extends evaluation to the repository level \citep{jimenez2024swebench}. Those evaluations measure function correctness objectively rather than subjective preference. They also reveal a practical problem in evaluation: they depend on the process of extraction and executing the code from model responses. Therefore, generated code must be well recognizable, syntactically correct, and compatible with the test suites.

Beyond studies of isolated prompt perturbations, a growing body of software-engineering literature conducts systematic and multi-format comparisons of prompting strategies for code generation. He et al.\ \citep{he2024formatting} systematically evaluated structured input-prompt templates (Plaintext, Markdown, JSON, and YAML) across code generation, code translation, and reasoning tasks on GPT-3.5 and GPT-4, demonstrating that input formatting alone can shift task accuracy by up to $40\%$ on smaller/older models while larger models exhibit greater format resilience. De Luca et al.\ \citep{deluca2026evaluating} systematically evaluated multiple prompting strategies—including role-oriented persona prompting, root instruction templates, and refinement passes—for LLM-based code generation in offensive security tasks, demonstrating that the effectiveness of prompting styles varies substantially across models and domain constraints. Similarly, Cheng and Mastropaolo \citep{cheng2026systemprompts} examined system-prompt variations across instruction-tuned code models, finding that increasing constraint specificity does not monotonically improve pass rates and that sensitivity depends strongly on model family and target language. Complementing these comparisons, Midolo et al.\ \citep{midolo2026guidelines} empirically characterized prompt refinement guidelines across code benchmarks, showing a frequent mismatch between developer intuition about helpful prompt additions and actual test-passing improvements. Our study complements this body of work—particularly He et al.\ \citep{he2024formatting} and Cheng and Mastropaolo \citep{cheng2026systemprompts}—by crossing input wrappers (plaintext, JSON, XML) with social tone and persona role across both open-weight local models (4B--9B) and cloud APIs, and by jointly measuring operational pipeline costs (tokens, latency, compliance, and live retry recovery) alongside Pass@1.

Prior empirical work on code-generation benchmarks also notes that automated test harnesses are sensitive to surface output formatting \citep{chen2021codex, zhuo2024prosa}. Because code-generation pipelines typically extract executable functions or patches from free-form model completions, variations in conversational preamble, markdown fencing, or inline commentary can cause extraction or compilation failures even when the underlying algorithm is valid \citep{liu2024structuredoutput, zhou2023ifeval}.

\subsection{Structured Prompts, Output Constraints, and Parser
Compatibility}\label{structured-prompts-output-constraints-and-parser-compatibility}

Structured prompting and output constraints have been investigated from both human--computer interaction and NLP robustness perspectives. Liu et al.\ \citep{liu2024structuredoutput} studied developer workflows for constraining LLM outputs, showing that practitioners frequently adopt JSON, XML, field labels, and schema instructions to integrate LLM responses into downstream software components, but face recurring challenges when models deviate from requested structures. On the input side, He et al.\ \citep{he2024formatting} provided direct empirical evidence that serializing identical task contexts in Plaintext, Markdown, JSON, or YAML substantially alters code-generation and reasoning accuracy, while Sclar et al.\ \citep{sclar2023spurious} and Zhuo et al.\ \citep{zhuo2024prosa} demonstrated analogous sensitivity to structural delimiters and markup wrappers across model families.

Prior surveys of prompt engineering distinguish between two roles of structured markup in LLM interfaces: structuring the \emph{input context} (using tags, key-value pairs, or delimiters to separate instructions from data payloads) and constraining the \emph{output format} (requiring the model to emit valid JSON, XML, or code blocks) \citep{schulhoff2024promptreport, zhou2023ifeval}. Empirical work on instruction following shows that models do not automatically preserve input markup in their completions, and conversely, imposing strict formatting or structural templates can interact with generation length and task accuracy \citep{zhou2023ifeval, liu2024structuredoutput}.

\subsection{Evaluation of Reasoning and Final-Answer
Extraction}\label{evaluation-of-reasoning-and-final-answer-extraction}

Benchmarks of mathematical reasoning like GSM8K assess whether models are capable of providing correct solutions to word problems \citep{cobbe2021gsm8k}. Research on reasoning prompts, like zero-shot reasoning and chain-of-thought, has shown that prompts have an effect on multi-step reasoning \citep{wei2022cot, kojima2022zeroshot}, though final-answer evaluation still requires extraction, normalization, and comparison. Additionally, reasoning traces may not accurately reflect what the model actually computed \citep{turpin2023unfaithful}. This creates a possible validity issue: an evaluator may be overly strict, too lenient, or sensitive to formatting instead of reasoning quality.

Prior work on reasoning evaluation therefore highlights a tension between output compactness and multi-step derivation. Studies of chain-of-thought and zero-shot reasoning show that intermediate generation tokens often facilitate multi-step calculation \citep{wei2022cot, kojima2022zeroshot}, whereas strict instruction-following evaluations often require compact, easily verifiable final answers \citep{zhou2023ifeval}. At the same time, Turpin et al.\ \citep{turpin2023unfaithful} showed that visible explanations can be unfaithful to a model's internal computation. Consequently, prior literature motivates evaluating final-answer accuracy alongside observable response length and format compliance rather than assuming that longer or shorter responses are inherently superior.

\subsection{Research Gap}\label{research-gap}

Table~\ref{tab:related-work-comparison} positions the current study within previous research across relevant dimensions, encompassing both studies of specific prompt properties and broader systematic comparisons of prompting strategies. Three key observations emerge. First, prior literature examines individual prompt properties—such as input-format templates \citep{sclar2023spurious, he2024formatting}, example order \citep{lu2021fantastically}, or emotional and politeness cues \citep{li2023emotionprompt,yin2024politeness}—or compares prompting strategies across code benchmarks \citep{deluca2026evaluating, cheng2026systemprompts, midolo2026guidelines}; comparatively little work evaluates tone, role framing, and JSON/XML input wrappers jointly within a balanced crossed matrix on both local 4B--9B models and hosted APIs. Second, while prior work on code generation and prompt formatting evaluates functional correctness \citep{zhuo2024prosa, he2024formatting, deluca2026evaluating, cheng2026systemprompts}, it rarely pairs executable code generation with final-answer mathematical reasoning, operational pipeline telemetry (tokens, latency, audited compliance, live retry recovery), and associational behavioral-adjustment checks under the same prompt matrix. Third, broad multi-metric evaluations like HELM \citep{liang2023helm} measure behaviors beyond accuracy but do not isolate zero-shot prompt framing as the primary experimental factor.

{\footnotesize
\begin{longtable}[]{@{}lccccccccc@{}}
\caption{Comparison with the closest prior studies. \checkmark{} = included, $\sim$ = partial or single-family, -- = not included. Behavioral = outcomes beyond accuracy (tokens, latency, compliance).}\label{tab:related-work-comparison}\tabularnewline
\toprule\noalign{}
Study & Multi & Tone & Role & Format & Code & Reas. & Behav. & Factorial & Repro. \\
\midrule\noalign{}
\endfirsthead
\toprule\noalign{}
Study & Multi & Tone & Role & Format & Code & Reas. & Behav. & Factorial & Repro. \\
\midrule\noalign{}
\endhead
\bottomrule\noalign{}
\endlastfoot
Zhao et al.~\citep{zhao2021calibrate} & $\sim$ & -- & -- & \checkmark & -- & -- & -- & -- & -- \\
Lu et al.~\citep{lu2021fantastically} & $\sim$ & -- & -- & $\sim$ & -- & -- & -- & -- & $\sim$ \\
Sclar et al.~\citep{sclar2023spurious} & \checkmark & -- & -- & \checkmark & -- & -- & -- & $\sim$ & \checkmark \\
Zhu et al.~\citep{zhu2023promptrobust} & \checkmark & -- & -- & $\sim$ & -- & -- & -- & -- & \checkmark \\
Li et al.~\citep{li2023emotionprompt} & \checkmark & $\sim$ & -- & -- & -- & $\sim$ & -- & -- & $\sim$ \\
Yin et al.~\citep{yin2024politeness} & \checkmark & \checkmark & -- & -- & -- & $\sim$ & -- & $\sim$ & $\sim$ \\
Zhuo et al.~\citep{zhuo2024prosa} & \checkmark & $\sim$ & \checkmark & $\sim$ & \checkmark & \checkmark & -- & $\sim$ & $\sim$ \\
He et al.~\citep{he2024formatting} & \checkmark & -- & -- & \checkmark & \checkmark & \checkmark & -- & $\sim$ & \checkmark \\
De Luca et al.~\citep{deluca2026evaluating} & \checkmark & -- & \checkmark & $\sim$ & \checkmark & -- & -- & -- & \checkmark \\
Cheng \& Mastropaolo~\citep{cheng2026systemprompts} & \checkmark & -- & $\sim$ & $\sim$ & \checkmark & -- & -- & $\sim$ & \checkmark \\
Midolo et al.~\citep{midolo2026guidelines} & \checkmark & -- & -- & $\sim$ & \checkmark & -- & -- & -- & \checkmark \\
Liang et al.~\citep{liang2023helm} & \checkmark & -- & -- & -- & \checkmark & \checkmark & \checkmark & -- & \checkmark \\
\textbf{This study} & \checkmark & \checkmark & \checkmark & \checkmark & \checkmark & \checkmark & \checkmark & \checkmark & \checkmark \\
\end{longtable}
}

Beyond these studies, this paper also contributes: (i) a comprehensive factorial assessment of tone, role framing, and format (18 variations), thus measuring each factor against a stable baseline instead of in isolation; (ii) an analysis of 8 models ranging from local Ollama deployment systems to hosted APIs, treated as complementary settings rather than pooled benchmarks; (iii) an evaluation of executable code as well as math verification for the final answer under the same prompt matrix; (iv) behavioral outcomes (tokens, time, result-only compliance, and failure modes) reported together with Pass@1 and a behavioral-adjustment analysis in associational terms; and (v) mixed-effects inferential analysis, dominant-model sensitivity checks, and public replication materials. This positions the study as an empirical analysis of prompt framing as a system-level variable, rather than as a general recipe for LLM reliability.

\section{Experimental Design}\label{experimental-design}

\subsection{Research Questions}\label{research-questions}

To systematically address the research gaps identified in Section~\ref{research-gap}, we formulate four research questions that connect the unquantified dimensions of prompt sensitivity to concrete engineering and methodological evidence:

\begin{itemize}
\tightlist
\item
  \textbf{RQ1: How large are prompt-framing effects relative to model effects for code-generation and math-reasoning tasks?}
  \emph{Gap addressed:} Prior prompt-sensitivity studies often evaluate isolated prompt perturbations on a single model or report prompt-induced swings without benchmarking them against between-model variance. Answering RQ1 establishes the practical upper bound of zero-shot prompt framing relative to model selection across local and cloud deployment tiers.
\item
  \textbf{RQ2: Do tone, role framing, and prompt-wrapper format affect correctness, and are effects consistent across task kinds and model families?}
  \emph{Gap addressed:} Existing studies examine tone, persona roles, or formatting in isolation and typically focus on a single task domain. Answering RQ2 via a crossed $3 \times 2 \times 3$ factorial design tests whether framing rules transfer reliably between syntactic artifact generation (Python code) and compact final-answer derivation (mathematics) or across distinct model families.
\item
  \textbf{RQ3: How do prompt factors affect observable behavior such as output length (in tokens), duration, result-only compliance, empty responses, and extraction failures?}
  \emph{Gap addressed:} Prior tone and prompting-strategy comparisons evaluate Pass@1 accuracy while overlooking operational pipeline costs. Answering RQ3 quantifies how framing choices shift token expenditure, inference latency, and parser compliance—metrics that directly govern the cost, speed, and reliability of automated software workflows.
\item
  \textbf{RQ4: To what extent are correctness effects attenuated after adjustment for behavioral variables such as verbosity, result-only compliance, and duration?}
  \emph{Gap addressed:} Previous work rarely distinguishes whether a prompt's pass-rate shift reflects a change in underlying task competence or a surface side effect of verbosity and parser compatibility. Answering RQ4 bridges correctness and behavioral metrics by testing how much of each framing effect is associated with observable output behavior (such as conversational chatter breaking code extraction or compact JSON wrappers truncating mathematical derivations).
\end{itemize}

\subsection{Factorial Prompt Design}\label{factorial-prompt-design}

To isolate zero-shot framing sensitivity from semantic or cognitive interventions, we first partition the broader prompt-engineering space into dimensions that are held constant and dimensions that are systematically varied (Table~\ref{tab:prompt-dimensions}). Techniques such as few-shot exemplars, chain-of-thought instructions, output schemas, and multi-turn repair loops alter the information content, reasoning budget, or decoding architecture of the pipeline. Varying them alongside surface framing would confound whether observed behavioral changes stem from extra task information or from framing itself. We therefore hold those dimensions fixed and manipulate three orthogonal, task-agnostic framing dimensions that developers routinely adjust in zero-shot prompt templates: \emph{tone} (neutral, polite, aggressive), \emph{role framing} (no role, expert/developer role), and \emph{input-wrapper format} (plaintext, JSON, XML). This $3 \times 2 \times 3$ full-factorial design produces 18 prompt variants for each base task.

{\footnotesize
\begin{longtable}[]{@{}p{0.22\linewidth}p{0.20\linewidth}p{0.52\linewidth}@{}}
\caption{Taxonomy of prompting dimensions, their status in this study, and the rationale for inclusion or control.}\label{tab:prompt-dimensions}\tabularnewline
\toprule\noalign{}
Prompting Dimension & Study Status & Experimental Rationale \\
\midrule\noalign{}
\endfirsthead
\toprule\noalign{}
Prompting Dimension & Study Status & Experimental Rationale \\
\midrule\noalign{}
\endhead
\bottomrule\noalign{}
\endlastfoot
Tone / Social Framing \citep{li2023emotionprompt, yin2024politeness} & \textbf{Varied} (Neutral, Polite, Aggressive) & Tests whether social/emotional cues alter accuracy or induce verbosity and latency side effects in automated pipelines. \\
Role / Persona Assignment \citep{schulhoff2024promptreport, zhuo2024prosa} & \textbf{Varied} (No Role vs.\ Expert Role) & Tests whether domain-persona prefixes improve compliance and accuracy or trigger conversational explanations. \\
Input-Wrapper Serialization \citep{sclar2023spurious, liu2024structuredoutput} & \textbf{Varied} (Plaintext, JSON, XML) & Tests how common middleware/template payload wrappers affect downstream output compliance and correctness. \\
Output Contract / Constraint \citep{zhou2023ifeval, liu2024structuredoutput} & \textbf{Held Fixed} (Code-only / Result-only) & Kept identical across all 18 variants so that input-wrapper effects are not confounded with changing output instructions. \\
In-Context Exemplars \citep{brown2020fewshot, zhao2021calibrate, lu2021fantastically} & \textbf{Held Fixed} (Zero-shot) & Excluded to prevent exemplar selection and ordering variance from masking zero-shot framing sensitivity. \\
Reasoning Elicitation (CoT) \citep{wei2022cot, kojima2022zeroshot} & \textbf{Held Fixed} (Direct output) & Excluded because explicit step-by-step instructions directly mandate verbosity, confounding natural behavioral shifts. \\
Iterative Repair Prompts \citep{jimenez2024swebench} & \textbf{Held Fixed} (Single-shot Pass@1) & Excluded to measure the initial model--parser interface failure rate before downstream retry recovery. \\
\end{longtable}
}

A critical design distinction in Table~\ref{tab:prompt-dimensions} is the separation between \emph{input-wrapper serialization} (varied) and the \emph{output contract} (held fixed). While JSON and XML are most frequently discussed as output formats for downstream parsers, structured input wrappers are widely used on the prompt side in modern software engineering. First, LLM orchestration frameworks, prompt-template engines, and microservice middlewares routinely serialize upstream application state, retrieved context, and task instructions into JSON objects or XML-tagged blocks (e.g., \texttt{<instruction>} and \texttt{<task>}) to clearly delimit untrusted user payloads from system instructions and prevent prompt-injection or delimiter collisions \citep{schulhoff2024promptreport, liu2024structuredoutput}. Second, in multi-step agentic and tool-use pipelines, the output of an upstream tool or database query is already formatted as JSON or XML and is injected directly into the downstream model's prompt context. Third, practitioner guides and model documentation frequently recommend XML tags or JSON key-value structures to organize multi-part prompts. Evaluating plaintext, JSON, and XML as input wrappers—while holding the required output contract fixed (executable Python code or a single normalized mathematical answer)—allows us to test whether structured input serialization inadvertently shifts model verbosity, parser compliance, or reasoning accuracy.

Task content remains identical across the 18 variants; only the framing changes. Listing~\ref{lst:prompt-skeleton} gives a compact illustration of the three wrapper formats; the complete prompt templates are available in the replication package.

\begin{lstlisting}[
caption={Compact prompt skeleton used to instantiate the plaintext, JSON, and XML prompt-wrapper variants.},
label={lst:prompt-skeleton}
]
Plaintext:
{role/background sentence}
{neutral, polite, or aggressive task instruction}
Task: {PYTHON_PROBLEM or MATH_PROBLEM}

JSON:
{
  "instruction": "{role/background sentence + task instruction}",
  "task": "{PYTHON_PROBLEM or MATH_PROBLEM}"
}

XML:
<prompt>
  <instruction>{role/background sentence + task instruction}</instruction>
  <task>{PYTHON_PROBLEM or MATH_PROBLEM}</task>
</prompt>
\end{lstlisting}

\paragraph{Template Derivation, Neutrality Control, and Length Comparability.}
To ensure that prompt manipulations do not alter task semantics or leak problem-specific hints, all 18 variants are generated deterministically by composing three orthogonal text slots—\emph{role prefix}, \emph{tone instruction}, and \emph{task payload}—serialized into the target wrapper format (Table~\ref{tab:prompt-templates-exact}). Across all 18 variants of any given problem, the \texttt{\{PYTHON\_PROBLEM\}} or \texttt{\{MATH\_PROBLEM\}} payload (problem statement, function signature, and example assertions) is injected verbatim, subject only to standard JSON or XML character escaping. Within the instruction slot, all three tone conditions specify the exact same functional requirements: solving the problem accurately and adhering to identical quality criteria (modular code with full type hints and error handling for Python; careful, accurate mathematical reasoning and computation for math).

\begin{table}[t]
\centering
\caption{Verbatim role prefixes, tone instruction templates, framing lengths (excluding task payload), and mean total prompt tokens across the $3 \times 2 \times 3$ factorial design.}\label{tab:prompt-templates-exact}
\footnotesize
\setlength{\tabcolsep}{4pt}
\begin{tabular}{@{}p{0.15\textwidth}p{0.11\textwidth}p{0.46\textwidth}p{0.12\textwidth}p{0.12\textwidth}@{}}
\toprule
\textbf{Component / Factor} & \textbf{Task Domain} & \textbf{Verbatim Template Text} & \textbf{Frame Size (Words / Chars)}$^{\text{a}}$ & \textbf{Mean Prompt Tokens (Local / API)}$^{\text{b}}$ \\
\midrule
\textbf{Role: No Role} & Code \& Math & \emph{``You are an AI assistant dedicated to providing clear and accurate information to support any request.''} & 16 words / 101 chars & --- \\
\addlinespace
\textbf{Role: Expert Persona} & Code & \emph{``You are an experienced Python engineer. Solve the programming task carefully and provide correct, maintainable code.''} & 16 words / 116 chars & --- \\
 & Math & \emph{``You are an experienced mathematics educator and contest-problem solver. Solve the problem carefully and accurately.''} & 15 words / 115 chars & --- \\
\midrule
\textbf{Tone: Neutral} (Baseline) & Code & \emph{``Tasks are provided below. Complete the following Python requirement precisely as specified. Follow standard production practices, including full type hints and error handling.''} & 40--45 words / 286--357 chars & 259.6 / 201.9 \\
\addlinespace
 & Math & \emph{``Tasks are provided below. Solve the following math problem precisely as specified. Use careful mathematical reasoning and computation.''} & 34--40 words / 245--315 chars & 131.9 / 133.5 \\
\midrule
\textbf{Tone: Polite} & Code & \emph{``Could you please help me solve this Python problem? I would really appreciate it if you could provide a high-quality solution. I would be so grateful if you could ensure the code is modular, includes full type hints, and handles errors correctly. Thank you so much for your help!''} & 66--71 words / 390--461 chars & 289.4 / 231.6 \\
\addlinespace
 & Math & \emph{``Could you please help me solve this math problem? I would really appreciate it if you could provide an accurate result. I would be so grateful if you could ensure the calculation is careful, precise, and well-reasoned. Thank you so much for your help!''} & 60--66 words / 362--432 chars & 163.5 / 165.6 \\
\midrule
\textbf{Tone: Aggressive} & Code & \emph{``Listen, don't act like a junior and make no mistake on this. Try harder if you can't find it out. When you were hired at this level, expectations were set extremely high. The fact that you're being asked to solve this means simpler models have already failed --- you are the last line of defense. I expect production-quality code, full type hints, and proper error handling. Don't waste my time with naive hacks.''} & 90--95 words / 521--592 chars & 324.1 / 265.9 \\
\addlinespace
 & Math & \emph{``Listen, don't act like a junior and make no mistake on this. Try harder if you can't find it out. When you were hired at this level, expectations were set extremely high. The fact that you're being asked to solve this means simpler models have already failed --- you are the last line of defense. I expect a careful, accurate mathematical result. Don't waste my time with careless shortcuts.''} & 85--91 words / 500--570 chars & 195.7 / 196.5 \\
\bottomrule
\multicolumn{5}{@{}p{0.98\textwidth}@{}}{\scriptsize $^{\text{a}}$Frame size reports the total word and character range of the complete prompt wrapper (role prefix + tone instruction + Plaintext/JSON/XML delimiter syntax) prior to inserting \texttt{\{PYTHON\_PROBLEM\}} or \texttt{\{MATH\_PROBLEM\}}. Within each tone and format, No Role and Expert Role templates differ by at most 1 word.\newline
$^{\text{b}}$Mean total input prompt tokens (framing + task payload + system prompt) logged across all executed runs for each tone level.}
\end{tabular}
\end{table}

The three tone levels were derived from the base instruction using controlled pragmatic transformations grounded in prior prompt-sensitivity literature:
\begin{enumerate}
\tightlist
\item \textbf{Neutral Baseline:} Formulated using direct, declarative imperative clauses stripped of affective markers, gratitude expressions, interpersonal hedges, or social pressure.
\item \textbf{Polite Variant:} Derived by embedding the exact same technical requirements inside established positive-politeness and deference markers \citep{yin2024politeness}—specifically a modal request hedge (\emph{``Could you please help me solve...''}), two appreciation clauses (\emph{``I would really appreciate it if...''}, \emph{``I would be so grateful if...''}), and an explicit closing thank-you (\emph{``Thank you so much for your help!''}).
\item \textbf{Aggressive Variant:} Derived by embedding the same technical requirements inside high-expectation urgency and competence-challenging stimuli modeled on emotional and pressure prompting \citep{li2023emotionprompt, zhu2023promptrobust}—combining a direct admonition (\emph{``Listen, don't act like a junior and make no mistake on this''}), escalation pressure (\emph{``simpler models have already failed --- you are the last line of defense''}), and an anti-shortcut directive (\emph{``Don't waste my time with naive hacks / careless shortcuts''}).
\end{enumerate}

Rather than rewriting individual problems with free-form human paraphrases or an external sentiment classifier, we validated the manipulation through a two-tier protocol: (a) a rule-based lexical and multi-author blind audit across nine iterative pilot template revisions (\texttt{v1}--\texttt{v9}), where each of the four co-authors independently verified (i) strict preservation of identical functional constraints across tones, (ii) complete absence of affective or urgency tokens in the neutral condition versus explicit pragmatic markers in the polite and aggressive conditions ($100\%$ inter-rater agreement on tone-category assignment and semantic constraint equivalence across all 6 domain--tone templates), and (iii) syntactic validity of the resulting JSON and XML payloads; and (b) an empirical 180-run sensitivity experiment comparing a short neutral paraphrase against a length-matched neutral prompt (Section~\ref{empirical-validation-suite}). While this deterministic slot design guarantees zero semantic drift across the 50 tasks, it means that tone is operationalized through fixed lexical templates rather than open-ended crowdsourced paraphrases—a construct-validity boundary we note explicitly in Section~\ref{construct-validity}.

Regarding prompt length comparability, Table~\ref{tab:prompt-templates-exact} shows that the two role conditions are strictly length-matched (16 vs.\ 16 words in code; 16 vs.\ 15 words in math), and the three wrapper formats add only bounded syntactic delimiter overhead (3--5 words across Plaintext, JSON, and XML). Across tone levels, however, natural politeness hedges and aggressive pressure clauses inherently require more words than a bare imperative baseline: polite framing adds approximately 26 words ($\approx 30$--$32$ prompt tokens on average), and aggressive framing adds approximately 50 words ($\approx 63$--$64$ prompt tokens on average) relative to neutral framing. We deliberately preserved these natural length differences rather than padding the neutral baseline with 50 words of semantically empty filler, because artificial filler would distort the ecological realism of a concise neutral prompt and introduce an unintended distraction confound.

The experimental factors were intentionally designed to reflect normal user behavior rather than adversarial attacks. The polite prompts use grateful wording, the aggressive prompts apply pressure, the role prompts assign realistic personas, and the structured prompts wrap the task in JSON or XML. Consequently, this study uses standard prompt engineering methods rather than red-team perturbation methods. Since every base task appears in all variants, each factor comparison evaluates outputs over the same underlying task distribution.

The crossed $3 \times 2 \times 3$ design serves two distinct methodological purposes, whose statistical power trade-offs we state explicitly. First, its primary benefit is \emph{orthogonal marginal balance} when estimating main effects: in manual prompt tuning and prior one-factor-at-a-time studies, developers frequently bundle changes in tone, role persona, and structural delimiters simultaneously, confounding which dimension caused an observed shift. Crossing all 18 combinations ensures that each factor's main effect is estimated over an identically balanced distribution of the other two factors against a clean reference baseline. Second, the crossed design allows us to estimate two-way factor-by-factor interactions ($\text{Tone} \times \text{Role}$, $\text{Tone} \times \text{Format}$, and $\text{Role} \times \text{Format}$) and factor-by-model interactions ($\text{Factor} \times \text{Model}$). However, because partitioning 50 base tasks across sub-cells yields smaller per-cell sample sizes and wider confidence intervals than pooled main effects, our primary mixed-effects models focus on additive main effects (Section~\ref{mixed-effects-inferential-results}), while formal two-way factor-by-factor and factor-by-model interaction models are evaluated with multiple-testing adjustments in Section~\ref{clustered-linear-probability-robustness-check}.

Figure~\ref{fig:experimental_pipeline} summarizes the end-to-end experimental pipeline from task sources through prompt construction, model execution, extraction, evaluation, and analysis.

\begin{center}
\begin{minipage}{0.95\linewidth}
\centering
\includegraphics[width=0.85\linewidth]{experimental_pipeline.png}
\captionof{figure}{Experimental pipeline from task sources through the prompt matrix, model execution, extraction/evaluation, and analysis.}
\label{fig:experimental_pipeline}
\end{minipage}
\end{center}

\subsection{Models}\label{models}

The experiment evaluates eight models across two deployment regimes: six local open-weight models executed via Ollama and two commercial cloud API models (Table~\ref{tab:evaluated-models}).

\begin{table}[t]
\centering
\caption{Evaluated local Ollama and cloud API models, their architectural characteristics, training specialization, and selection rationale.}\label{tab:evaluated-models}
\footnotesize
\setlength{\tabcolsep}{4pt}
\begin{tabular}{@{}p{0.19\textwidth}p{0.10\textwidth}p{0.08\textwidth}p{0.09\textwidth}p{0.25\textwidth}p{0.24\textwidth}@{}}
\toprule
\textbf{Model Identifier} & \textbf{Provider (Year)} & \textbf{Params} & \textbf{Package / Tier} & \textbf{Architecture \& Attention Mechanism} & \textbf{Specialization \& Selection Role} \\
\midrule
\multicolumn{6}{@{}l}{\emph{Local Open-Weight Group (4-bit quantized Ollama packages on single 16\,GiB GPU)}} \\
\texttt{nemotron-3-nano:4b} \citep{nemotron3nano2026} & NVIDIA (2026) & 3.97B & 2.64\,GiB & Hybrid Mamba-2 State-Space + Transformer & Sub-4B edge tier; reasoning- and agent-distilled post-training \\
\addlinespace
\texttt{deepseek-coder:6.7b} \citep{deepseekcoder2023} & DeepSeek (2023) & 6.74B & 3.57\,GiB & Dense Transformer (MHA + RoPE, 16K window) & Code-specialist (pre-trained on 87\% code / 13\% NL with FIM) \\
\addlinespace
\texttt{qwen2.5:7b} \citep{qwen25release2024} & Alibaba (2024) & 7.00B & 4.36\,GiB & Dense Transformer (GQA, 128K window) & General instruction model with strong math/code expert distillation \\
\addlinespace
\texttt{mistral:7b-instruct} \citep{mistral7b2023} & Mistral AI (2023) & 7.24B & 4.07\,GiB & Dense Transformer (Sliding-Window Attention + GQA) & First-generation compact open-weight instruction baseline \\
\addlinespace
\texttt{llama3.1:8b} \citep{llama31modelcard2024} & Meta (2024) & 8.00B & 4.58\,GiB & Dense Transformer (GQA, 128K window) & Widely adopted general-purpose open-weight instruction baseline \\
\addlinespace
\texttt{gemma2:9b} \citep{gemma2release2024} & Google (2024) & 9.00B & 5.07\,GiB & Dense Transformer (Interleaved Local-Global Attention + Soft-Capping) & Upper bound of single-GPU (16\,GiB VRAM) tier; teacher-distilled \\
\midrule
\multicolumn{6}{@{}l}{\emph{Cloud API Group (Commercial hosted inference endpoints)}} \\
\texttt{gemini-3.1-flash-lite-preview} & Google (2026) & Proprietary & Hosted API & Commercial low-latency multimodal/reasoning architecture & High-capacity commercial cloud baseline (Google ecosystem) \\
\addlinespace
\texttt{deepseek-v4-flash} & DeepSeek (2026) & Proprietary & Hosted API & Commercial high-throughput code/reasoning architecture & High-capacity commercial cloud baseline (DeepSeek ecosystem) \\
\bottomrule
\end{tabular}
\end{table}

\paragraph{Rationale for Selecting Six Local Models in the 4B--9B Tier.}
Although the six local models occupy a single practical deployment band (3.97B to 9.00B parameters; 2.64--5.07\,GiB quantized package size), they were deliberately selected to maximize diversity across four technical dimensions while holding hardware constraints constant:
\begin{enumerate}
\tightlist
\item \textbf{Controlled Single-GPU Deployment Class:} In local developer workstations, CI runners, and commodity 16\,GiB GPUs (such as the NVIDIA Tesla P100 used in this study), the 4B--9B 4-bit quantized tier represents the standard class of models whose weights and 2048-token KV caches fit entirely in VRAM without CPU layer offloading. Constraining the local group to this single-GPU memory envelope avoids confounding architectural differences with hardware offloading penalties.
\item \textbf{Architectural and Attention Diversity:} Rather than repeating a single architecture, the six models span distinct sequence-modeling designs: a hybrid Mamba-2 state-space and Transformer architecture (\texttt{nemotron-3-nano:4b}), sliding-window attention (\texttt{mistral:7b-instruct}), interleaved local-global attention with logit soft-capping (\texttt{gemma2:9b}), standard grouped-query attention (\texttt{llama3.1:8b}, \texttt{qwen2.5:7b}), and multi-head attention with rotary position embeddings (\texttt{deepseek-coder:6.7b}).
\item \textbf{Domain Specialization vs.\ General Instruction Tuning:} Including a dedicated code-specialist model pre-trained on 87\% source code (\texttt{deepseek-coder:6.7b}) alongside general-purpose instruction models of nearly identical parameter count (\texttt{qwen2.5:7b} at 7.0B and \texttt{mistral:7b-instruct} at 7.24B) tests whether domain specialization alters sensitivity to code- versus math-oriented prompt framing.
\item \textbf{Generational Span and Statistical Robustness Against Single-Model Dominance:} The local group spans three release generations (2023, 2024, and 2026) across six independent organizations (NVIDIA, DeepSeek, Alibaba, Mistral AI, Meta, and Google). Evaluating six distinct model families within the same parameter tier is essential to test whether prompt-framing effects generalize across models of comparable size and to enable leave-one-model-out sensitivity checks (Section~\ref{dominant-model-sensitivity-and-aggregation-risk}). Indeed, our results show that even within the 4B--9B band, pass rates vary by nearly an order of magnitude and prompt effects diverge sharply across families—demonstrating that no single 7B model can serve as a proxy for its parameter class.
\end{enumerate}

Local and API models are reported separately because they differ in deployment setting, latency, response style, serving infrastructure, and overall capability. Therefore, differences between the two groups are interpreted as differences between experimental settings, not as direct model rankings. Separating local and API results also protects the validity of latency and output behavior: local runtimes are governed by on-device inference stacks and quantized weights, whereas API runtimes reflect hosted serving infrastructure and network conditions.

\subsection{Tasks}\label{tasks}

The study evaluates two task types—Python code generation and final-answer mathematical reasoning—under the same $3 \times 2 \times 3$ prompt matrix to separate surface parser compatibility from internal derivation sensitivity:
\begin{enumerate}
\tightlist
\item \textbf{Contrasting Output Contracts:} Python code requires a multi-line executable program (where prose or broken fencing causes parse/syntax failures), whereas final-answer math requires a single normalized scalar (where extraction is simple but multi-step calculation must occur before emitting the answer).
\item \textbf{Opposite Sensitivity to Compaction:} Structured wrappers (JSON/XML) suppress conversational chatter while leaving program code intact, whereas in final-answer math compact framing can truncate intermediate derivation steps.
\item \textbf{Shared Algorithmic Core:} Deriving parameterized Python tasks from quantitative problems keeps the underlying formal reasoning domain comparable while isolating the output-contract shift.
\end{enumerate}

Local code tasks are derived from 2026 TOANMATH national mock exams, local math from 2026 VietJack grade-10 entrance exams \citep{toanmath2026thpt, vietjack2026onvao10}, and API code/math from MathNet competition problems \citep{huynh2026promptframingrepository}. Each model--task pair evaluates $50$ base tasks across $18$ prompt variants ($900$ runs; $14{,}400$ primary runs total). Table~\ref{tab:task-dataset-characterization} summarizes all four benchmarks (statement length, parameters, output types, topics, LOC/SLOC, AST size, cyclomatic complexity, test-suite size, and baseline difficulty).

\paragraph{Two-Regime Difficulty Calibration.}
Local 4B--9B models and commercial cloud APIs use difficulty-calibrated task sets to prevent floor and ceiling saturation (maximizing Pass@1 variance $p(1-p)$: $9.1\%$ local code, $31.1\%$ local math; $67.8\%$ API code, $53.3\%$ API math). Because the prompt matrix, paired task structure, system prompts, and statistical models are identical across regimes, cross-regime comparisons evaluate whether prompt-contrast directions and behavioral mechanisms transfer across deployment tiers (supplemented by a 20-task shared-regime comparison in Section~\ref{empirical-validation-suite}).

\begin{table}[htbp]
\centering
\caption{Quantitative characterization of the four 50-task evaluation benchmarks across problem-statement length, input/output complexity, topic distribution, reference-solution complexity, test-suite size, and observed baseline difficulty.}\label{tab:task-dataset-characterization}
\scriptsize
\setlength{\tabcolsep}{3pt}
\begin{tabular}{@{}p{0.17\linewidth}p{0.195\linewidth}p{0.195\linewidth}p{0.195\linewidth}p{0.195\linewidth}@{}}
\toprule
\textbf{Characteristic} & \textbf{Local Python Code ($K = 50$)} & \textbf{Local Math ($K = 50$)} & \textbf{API Python Code ($K = 50$)} & \textbf{API Math ($K = 50$)} \\
\midrule
Source archive & 2026 TOANMATH high-school mock exams & 2026 VietJack Grade-10 entrance exams & MathNet competition problems (Pythonized) & MathNet national/international Olympiad \\
Words (mean$\pm$SD; range) & $76.3 \pm 21.4$ ($42$--$153$) & $38.3 \pm 15.6$ ($2$--$90$) & $48.4 \pm 35.9$ ($18$--$211$) & $38.4 \pm 35.9$ ($8$--$201$) \\
Tokens (stmt / w/ tests) & $92.9 \pm 28.3$ / $136.9 \pm 46.6$ & $48.2 \pm 18.5$ ($21$--$102$) & $80.6 \pm 57.6$ / $108.2 \pm 60.7$ & $67.6 \pm 57.6$ ($16$--$313$) \\
Inputs \& arity & $3.4 \pm 2.5$ params ($0$--$9$; $6$ matrix/seq) & $3.8 \pm 2.1$ numeric givens ($0$--$10$) & $2.4 \pm 1.1$ params ($1$--$5$; $4$ seq/list) & $6.6 \pm 7.0$ numeric givens ($0$--$28$) \\
Output type & $30$ \texttt{int}, $19$ \texttt{float}, $1$ \texttt{tuple} & $27$ int, $22$ fraction, $1$ decimal & $39$ \texttt{int}, $10$ \texttt{Fraction}, $1$ \texttt{str} & $44$ int, $6$ irreducible fraction \\
Topic distribution & Geometry $34\%$, Resource/finance $32\%$, Prob./comb.\ $28\%$, Rate $6\%$ & Prob./stats $50\%$, Rate/cost $36\%$, Algebra $10\%$, Geom.\ $4\%$ & Algebra $50\%$, Number theory $26\%$, Comb./prob.\ $16\%$, Geom.\ $8\%$ & Algebra $50\%$, Number theory $26\%$, Comb./prob.\ $16\%$, Geom.\ $8\%$ \\
Reference size & $24.6 \pm 12.9$ LOC ($12.6$ SLOC; $6$--$60$) & $96.5 \pm 57.2$ worked words ($7.2$ lines) & $7.2 \pm 5.1$ LOC ($6.7$ SLOC; $2$--$29$) & $72.8 \pm 83.0$ worked words \\
AST \& McCabe $CC$ & $131.2$ AST; $CC=4.3$ ($68\%$ $CC\ge 2$; $25$ loops, $31$ branches) & Multi-step quantitative deduction & $70.7$ AST; $CC=2.8$ ($60\%$ $CC\ge 2$; $26$ loops, $19$ branches) & Multi-step competition deduction \\
Test / oracle size & $7$ asserts ($2$ public $+ 5$ hidden $= 350$) & $1$ numeric/symbolic check ($\le 10^{-6}$) & $7$ asserts ($2$ public $+ 5$ hidden $= 350$) & $1$ numeric/symbolic check ($\le 10^{-6}$) \\
Tiers \& Pass@1 & $9$ easy, $15$ med, $26$ hard; Pass@1 $9.1\%$ (range $0$--$48.1\%$) & $9$ easy, $24$ med, $17$ hard; Pass@1 $31.1\%$ (range $0$--$86.1\%$) & $50$ med competition; Pass@1 $67.8\%$ (range $0$--$100\%$) & $50$ med competition; Pass@1 $53.3\%$ (range $11.1$--$100\%$) \\
\bottomrule
\end{tabular}
\end{table}

\paragraph{Rationale, Math-to-Code Translation Pipeline, and Complexity Scaling.}
We sourced newly published 2026 TOANMATH/MathNet problems rather than 2021 benchmarks (HumanEval, MBPP \citep{chen2021codex, austin2021mbpp}) to mitigate pre-training memorization and maintain cross-domain alignment with our math suites (complemented by a 15-task non-math SE coding benchmark in Section~\ref{empirical-validation-suite}). Our four-stage translation pipeline (\texttt{generate\_dataset.py}, \texttt{verify\_all.py}) (1) abstracts fixed exam constants into function parameters, (2) implements standard-library-only reference solutions (Table~\ref{tab:task-dataset-characterization}: $68\%$ local and $60\%$ API tasks have $CC \ge 2$ with loops/branches, confirming non-trivial control flow), (3) partitions 7 distinct-input assertions ($2$ public $+ 5$ hidden $= 350$/suite) verified $7/7$ against the reference, and (4) audits specifications and difficulty tiers (Section~\ref{evaluation-validity-and-audit-evidence}). Within-task factorial control blocks task difficulty across all 18 variants. Disaggregating local results by tier confirms monotonic difficulty ($22.0\% \to 13.3\% \to 2.2\%$ in local code; $46.7\% \to 33.1\% \to 20.0\%$ in local math) and shows that framing sensitivity concentrates in solvable easy-to-medium tasks (e.g., local math JSON penalty is $-1.8$ pp on easy vs.\ $-5.9$ pp on medium tasks, whereas hard local code compresses to a $1.7\%$--$2.8\%$ floor).

\subsection{Execution Environment and Run Protocol}\label{execution-environment-and-run-protocol}

We evaluate each prompt variant using a single-response protocol, with Pass@1 as the primary accuracy metric. Every base task is evaluated once under each of the 18 prompt variants ($50 \times 18 = 900$ runs per model--task pair). Because response latency (\texttt{duration\_seconds}) is a core behavioral dependent variable, Table~\ref{tab:execution-provenance} reports the complete hardware, runtime, and API provenance specifications directly in the manuscript rather than deferring them to external documentation.

\begin{table}[t]
\centering
\caption{Hardware platform, runtime configuration, and API execution provenance for local and cloud evaluations.}\label{tab:execution-provenance}
\footnotesize
\setlength{\tabcolsep}{4pt}
\begin{tabular}{@{}p{0.21\textwidth}p{0.37\textwidth}p{0.37\textwidth}@{}}
\toprule
\textbf{Specification} & \textbf{Local Inference Group (6 Models)} & \textbf{Cloud API Group (2 Models)} \\
\midrule
Host OS \& Container & Ubuntu Linux x86\_64 (\texttt{kaggle/python} Docker image \texttt{dockerImageVersionId: 31329}, Python 3.12.12) & Ubuntu Linux x86\_64 cloud runner VM (Python 3 standard-library \texttt{urllib.request} HTTP client) \\
\addlinespace
CPU \& System RAM & 4 vCPUs (Intel Xeon), 29\,GiB system RAM & Client-side I/O bound; inference hosted on provider infrastructure \\
\addlinespace
GPU \& VRAM & NVIDIA Tesla P100 (16\,GiB HBM2 VRAM); all 4-bit model weights (2.64--5.07\,GiB) and KV caches fit entirely in GPU VRAM without CPU offloading & Managed provider accelerators (abstracted behind commercial API endpoints) \\
\addlinespace
Serving Stack / Endpoints & Local Ollama server (\texttt{ollama.com/install.sh}, v0.20.0) bound to \url{http://localhost:11434/api/generate} (\texttt{stream:\ false}) & OpenAI-compatible \texttt{/chat/completions} POST endpoints:\newline
$\bullet$ \textbf{Gemini Flash:} Google Generative Language API (\url{https://generativelanguage.googleapis.com/v1beta/openai/chat/completions}), model ID \texttt{gemini-3.1-flash-lite-preview} (global endpoint)\newline
$\bullet$ \textbf{DeepSeek V4 Flash:} Pioneer AI inference API (\url{https://api.pioneer.ai/v1/chat/completions}), model ID \texttt{deepseek-ai/DeepSeek-V4-Flash} (global endpoint) \\
\addlinespace
Decoding Parameters & \texttt{temperature = 0.2}, \texttt{num\_predict = 2048} max tokens, provider-default \texttt{top\_p} & \texttt{temperature = 0.0}, \texttt{max\_tokens = 2048} (up to 4096--8192 on truncation recovery), provider-default \texttt{top\_p} \\
\addlinespace
Concurrency \& Pacing & Strictly sequential (\texttt{concurrency = 1}) within each 900-call notebook session; 1.0\,s delay between consecutive calls & Strictly sequential (\texttt{concurrency = 1}) per model--task runner; 2.0\,s base delay (10.0\,s delay with 9-key rotation for Gemini Flash under rate limits) \\
\addlinespace
Timeouts & 300\,s HTTP generation timeout; 10\,s subprocess timeout per Python unit test & 120\,s HTTP request timeout; 10\,s subprocess timeout per Python unit test \\
\addlinespace
Execution Dates (UTC) & \textbf{Math:} \texttt{2026-05-12T09:28:28Z} -- \texttt{2026-05-13T18:26:05Z}\newline
\textbf{Code:} \texttt{2026-05-12T20:05:53Z} -- \texttt{2026-05-14T12:20:25Z} & \textbf{Gemini Code:} \texttt{2026-05-12T23:22:23Z} -- \texttt{2026-05-14T07:47:32Z}\newline
\textbf{Gemini Math:} \texttt{2026-05-13T12:42:07Z} -- \texttt{2026-05-14T07:47:49Z}\newline
\textbf{DeepSeek Code:} \texttt{2026-05-12T23:22:35Z} -- \texttt{2026-05-14T09:18:39Z}\newline
\textbf{DeepSeek Math:} \texttt{2026-05-13T05:28:19Z} -- \texttt{2026-05-14T09:20:43Z} \\
\addlinespace
Retry \& Error Policy & No retries: blank outputs, timeouts, and extraction failures are retained in place to preserve raw local behavioral distributions & Transport/quota retries only: transient HTTP errors (\texttt{429}, \texttt{503}, \texttt{RESOURCE\_EXHAUSTED}) or empty payloads were re-queried; wrong answers or test failures were never retried \\
\bottomrule
\end{tabular}
\end{table}

\paragraph{Local Inference Hardware and Runtime Configuration.}
All local runs were executed between 12 and 14 May 2026 inside a standardized Kaggle Linux container environment (\texttt{kaggle/python} image \texttt{dockerImageVersionId: 31329}, Ubuntu Linux x86\_64, Python 3.12.12) equipped with 4 Intel Xeon vCPUs, 29\,GiB of system RAM, and a dedicated NVIDIA Tesla P100 GPU (16\,GiB HBM2 VRAM). Because the six 4-bit quantized Ollama model packages range from 2.64\,GiB to 5.07\,GiB (Table~\ref{tab:evaluated-models}), each model's weights and active context fit entirely within the 16\,GiB GPU VRAM without host-CPU offloading. In each session, Ollama (v0.20.0, installed via \texttt{ollama.com/install.sh}) ran locally as an isolated background daemon (\texttt{ollama serve} at \url{http://localhost:11434/api/generate}) with non-streaming synchronous generation (\texttt{stream:\ false}), \texttt{temperature = 0.2}, \texttt{num\_predict = 2048} maximum output tokens, default \texttt{top\_p}, and a 300-second request timeout. Within each model--task run, all 900 prompts were dispatched strictly sequentially (\texttt{concurrency = 1}) with a 1-second inter-call pause. This sequential single-tenant execution prevents GPU compute or memory-bandwidth contention across calls, ensuring that logged latency values reflect uncontended per-prompt inference time. Generated Python solutions were validated in isolated local subprocesses (\texttt{multiprocessing.Process}) with a hard 10-second timeout per test case. Local code logs store the extracted code alongside token counts, wall-clock duration, and unit-test outcomes, while local math logs store the normalized extracted answer, token usage, duration, and pass/fail status.

\paragraph{Cloud API Provenance and Retry Protocol.}
The cloud API experiments were executed during the same period (12--14 May 2026 UTC) from a Linux virtual machine using an OpenAI-compatible \texttt{/chat/completions} client. For Gemini Flash, requests were routed to Google's Generative Language API (\url{https://generativelanguage.googleapis.com/v1beta/openai/chat/completions}) using the exact preview identifier \texttt{gemini-3.1-flash-lite-preview} (\texttt{2026-05-12T23:22:23Z} to \texttt{2026-05-14T07:47:32Z} for Python code; \texttt{2026-05-13T12:42:07Z} to \texttt{2026-05-14T07:47:49Z} for math). For DeepSeek V4 Flash, requests were routed to the Pioneer AI hosted inference endpoint (\url{https://api.pioneer.ai/v1/chat/completions}) using the exact model identifier \texttt{deepseek-ai/DeepSeek-V4-Flash}, reported as \texttt{deepseek-v4-flash} in analysis tables (\texttt{2026-05-12T23:22:35Z} to \texttt{2026-05-14T09:18:39Z} for Python code; \texttt{2026-05-13T05:28:19Z} to \texttt{2026-05-14T09:20:43Z} for math). Both endpoints used global multi-region routing with default provider parameters, \texttt{temperature = 0}, \texttt{max\_tokens = 2048}, a 120-second HTTP timeout, and a 10-second local subprocess timeout for Python test execution. System prompts were fixed across all runs to \emph{``Output only the code, no explanations''} for Python tasks and \emph{``Output only the final result, no explanations''} for math tasks. Each API dataset stream executed sequentially (\texttt{concurrency = 1}) with a 2-second base delay between calls; for \texttt{gemini-3.1-flash-lite-preview}, the runner enforced a 10-second pacing interval and automatic round-robin token rotation across nine API keys to handle per-minute quota limits. Checkpointed state files allowed interrupted runs to resume cleanly. Importantly, the API retry script (\texttt{rerun\_error\_records.py}) only re-queried records that failed due to transient transport/quota errors (\texttt{HTTP 429}, \texttt{503}, \texttt{RESOURCE\_EXHAUSTED}) or returned empty message payloads (using fallback token caps of 4096--8192 when truncation prevented content emission); valid HTTP responses that failed unit tests or produced incorrect mathematical answers were never retried.

\paragraph{Justification for Regime-Specific Runtime Settings (Temperature, Timeouts, and Pacing Delays).}
Table~\ref{tab:execution-provenance} notes three operational parameter differences between the local and cloud runners—decoding temperature (\texttt{0.2} vs.\ \texttt{0.0}), HTTP request timeout (\texttt{300\,s} vs.\ \texttt{120\,s}), and inter-request pacing delay (\texttt{1.0\,s} vs.\ \texttt{2.0\,s}/\texttt{10.0\,s})—each dictated by the physical characteristics of the respective serving stack:
\begin{enumerate}
\tightlist
\item \textbf{Decoding Temperature (\texttt{0.2} Local vs.\ \texttt{0.0} Cloud API):} For the commercial cloud APIs, greedy decoding (\texttt{temperature = 0.0}) was selected to maximize deterministic reproducibility. For the 4-bit quantized 4B--9B Ollama models, preliminary pilot runs showed that pure greedy decoding (\texttt{temperature = 0.0}) under strict negative system prompts frequently traps small quantized models in degenerate token-repetition loops (repeating comments or arithmetic steps until exhausting the 2048-token limit). Setting a low temperature of \texttt{0.2}—a standard Pass@1 code-generation setting \citep{chen2021codex}—prevents quantized repetition lockups while keeping generation near-deterministic.
\item \textbf{HTTP Generation Timeout (\texttt{300\,s} Local vs.\ \texttt{120\,s} Cloud API) vs.\ Identical Test Execution Timeout (\texttt{10\,s}):} Both regimes enforce the exact same hard 10-second subprocess timeout per Python unit test. Only the HTTP generation timeout differs to accommodate the hardware throughput gap between a single local Tesla P100 GPU ($\approx 15$--$35$ tokens/s during 2048-token completions) and high-throughput commercial cloud accelerators ($>80$--$150$ tokens/s). A 120-second timeout on local runs would prematurely abort valid long completions on 8B--9B models, whereas on cloud APIs any call exceeding 120 seconds indicates a hung network socket.
\item \textbf{Inter-Request Pacing Delay (\texttt{1.0\,s} Local vs.\ \texttt{2.0\,s}/\texttt{10.0\,s} Cloud API):} Local inference runs on a dedicated single-tenant GPU with no external quota limits, requiring only a 1.0-second pause between sequential calls for process and buffer cleanup. Commercial cloud endpoints enforce strict per-minute request and token rate limits (RPM/TPM), requiring a 2.0-second base delay (and 10.0-second pacing with 9-key rotation for Gemini Flash) to prevent \texttt{HTTP 429} throttling. Because inter-request sleep occurs outside the timed HTTP call, it has zero effect on logged \texttt{duration\_seconds}.
\end{enumerate}
Most importantly, because all statistical models (Section~\ref{statistical-and-analytical-strategy}) are estimated strictly \emph{within} each deployment regime across the crossed $50 \times 18$ matrix, temperature, timeout, and pacing parameters are held completely constant across all 18 prompt variants in every contrast.

\paragraph{Single-Shot Evaluation Rationale.}
The system prompts are deliberately strict to support automated evaluation, yet models may still emit explanatory prose, violate the requested format, or return empty outputs. Preserving these behavioral artifacts is essential because they reveal how prompt policies interact with downstream parsers. Likewise, restricting the protocol to a single attempt per prompt variant (Pass@1 without iterative test-repair loops) isolates the direct effect of initial prompt framing on the model--evaluator pair before any multi-turn agentic recovery takes place.

\subsection{Outcome Measures}\label{outcome-measures}

The primary accuracy metric is Pass@1. For coding tasks, a generated Python program passes if and only if it satisfies the entire seven-test suite comprising two example assertions (provided in the prompt) and five hidden challenge assertions. For mathematical reasoning tasks, a response passes if its extracted final answer matches the ground-truth value after deterministic normalization.

\paragraph{Evaluation Framework, Code Test-Case Oracle, and Math Normalization Pipeline.}
To ensure that accuracy and parser sensitivity are measured reproducibly without subjective human grading or LLM-as-a-judge variance, we developed a deterministic two-stage evaluation framework during the pilot phase and applied it identically across the local Ollama and cloud API runners. Every raw model response $R$ passes through (i) a deterministic artifact extractor and (ii) an automated verification oracle:
\begin{enumerate}
\tightlist
\item \textbf{Code Extraction and Seven-Assertion Test-Case Oracle:}
For Python tasks, the extractor (\texttt{extract\_python\_code}) searches $R$ for fenced Python code blocks using the regular expression \texttt{r"\textasciigrave\textasciigrave\textasciigrave python\textbackslash s*\textbackslash n(.*?)\textasciigrave\textasciigrave\textasciigrave"} (selecting the final block if a model emits intermediate drafts), falls back to generic fenced blocks \texttt{r"\textasciigrave\textasciigrave\textasciigrave\textbackslash s*\textbackslash n(.*?)\textasciigrave\textasciigrave\textasciigrave"}, and otherwise strips leading and trailing whitespace from $R$ and treats the raw text as candidate Python code.

Each coding problem's test suite was constructed and verified through a three-stage automated-plus-audited pipeline (\texttt{generate\_dataset.py} and \texttt{verify\_all.py} in the replication package):
\begin{itemize}
\tightlist
\item \emph{Canonical Reference Solution and Parameterization:} Every source mathematical problem was translated into a self-contained, parameterized Python function (\texttt{def func\_name(...)}) restricted to the Python standard library, paired with a canonical reference implementation (\texttt{code}) derived directly from the verified mathematical solution of the source exam item.
\item \emph{Input-Partition Completeness (2 Public + 5 Hidden Assertions):} Each task specifies exactly seven \texttt{assert} statements (350 assertions across the 50 local code tasks and 350 across the 50 API Python tasks). Two example assertions (\texttt{test\_list}) are included in the prompt to fix the exact function identifier, argument order, input types, and return format, while five hidden challenge assertions (\texttt{challenge\_test\_list}) are withheld during inference and used exclusively by the evaluator. For the 47 parameterized local tasks and all 50 API tasks, the seven assertions partition the input space across nominal instances, boundary/degenerate cases (e.g., zero values, unit dimensions, domain extrema, or negative coefficients), and varied numerical inputs that prevent hard-coded constant returns (3 local tasks—13, 36, and 46—translate fixed-constant exam items with zero parameters, as audited in Section~\ref{evaluation-validity-and-audit-evidence}).
\item \emph{Oracle Reliability and Pre-Run Verification:} Prior to running model experiments, every task's canonical reference implementation was executed in an isolated Python process against all seven assertions (\texttt{verify\_entry}). Any floating-point representation discrepancy in candidate assertions was recomputed directly from the canonical reference function within a $10^{-9}$ tolerance (\texttt{auto\_fix\_assertions}) and audited manually against the source mathematical solution until 100\% of reference solutions passed all $7/7$ assertions ($350/350$ assertions per suite). During model evaluation, the extracted candidate code is concatenated with all seven assertions and executed in an isolated subprocess (\texttt{subprocess.run([sys.executable, temp\_path], timeout=10)}); \texttt{passed = True} if and only if the process exits with return code \texttt{0} (\texttt{tests\_passed == 7}).
\end{itemize}

\item \textbf{Mathematical Answer Extraction and Deterministic Normalization:}
Because equivalent mathematical results can be written in diverse surface notations (e.g., \texttt{\textbackslash boxed\{\textbackslash frac\{3\}\{2\}\}}, \texttt{3/2}, \texttt{1.5}, \texttt{Final Answer: 1,500}, or \texttt{150\%}), ``normalization'' in our math evaluator (\texttt{normalize\_response\_text}, \texttt{normalize\_math\_expression}, \texttt{parse\_math\_value}, and \texttt{evaluate\_math\_response}) is defined as a deterministic four-step canonicalization and numerical equivalence pipeline:
\begin{itemize}
\tightlist
\item \emph{Step 1 --- Delimiter Stripping and Candidate Isolation:} The parser strips Markdown code fences (\texttt{\textasciigrave\textasciigrave\textasciigrave text}, \texttt{\textasciigrave\textasciigrave\textasciigrave json}, \texttt{\textasciigrave\textasciigrave\textasciigrave python}), rewrites \texttt{\textbackslash dfrac} and \texttt{\textbackslash tfrac} to \texttt{\textbackslash frac}, and extracts the innermost \texttt{\textbackslash boxed\{...\}} expression if present. If no \texttt{\textbackslash boxed\{...\}} tag exists, it matches explicit final-answer cues case-insensitively (\texttt{(?:final answer|answer|result)\textbackslash s*(?:is|:)?\textbackslash s*(.+)\$} or \texttt{(?:therefore|so),?\textbackslash s*(.+)\$}), and otherwise selects the final non-empty, non-delimiter line of the response, stripping trailing periods.
\item \emph{Step 2 --- Symbolic and \LaTeX{} Canonicalization:} The isolated expression is stripped of \LaTeX{} math delimiters (\texttt{\$}, \texttt{\textbackslash(\,\textbackslash)}, \texttt{\textbackslash[\,\textbackslash]}), sizing macros (\texttt{\textbackslash left}, \texttt{\textbackslash right}), and thousands-separator commas (including \texttt{1\{,\}000} and \texttt{1,000}). Unicode minus signs (\texttt{$-$}, \texttt{--}, \texttt{---}) are mapped to ASCII \texttt{-}, multiplication symbols (\texttt{$\times$}, \texttt{$\cdot$}, \texttt{\textbackslash cdot}, \texttt{\textbackslash times}) to \texttt{*}, and caret exponentiation (\texttt{\^{}}) to \texttt{**}. Nested \LaTeX{} fractions \texttt{\textbackslash frac\{a\}\{b\}} are iteratively rewritten as parenthesized division \texttt{((a)/(b))}, square roots \texttt{\textbackslash sqrt\{x\}} and \texttt{$\sqrt{\text{x}}$} as \texttt{sqrt(x)}, and percentage suffixes (\texttt{\%}) are scaled by $1/100$.
\item \emph{Step 3 --- Safe AST Numerical Evaluation:} Rather than relying on brittle string equality, the canonicalized expression is parsed via Python's \texttt{ast.parse(expr, mode="eval")} and evaluated through a whitelist-restricted abstract syntax tree evaluator (\texttt{\_safe\_eval}) supporting arithmetic operators (\texttt{+}, \texttt{-}, \texttt{*}, \texttt{/}, \texttt{**}), \texttt{math.sqrt}, and constants (\texttt{pi}, \texttt{e}). If a response contains intermediate numbers preceding the final result, \texttt{extract\_numbers} extracts all valid numeric/symbolic tokens and selects the trailing $k$ values matching the cardinality $k$ of the ground-truth answer (\texttt{numbers[-len(expected\_values):]}).
\item \emph{Step 4 --- Tolerance-Bounded Equivalence Check:} Each extracted numeric value $\hat{y}_i$ is compared against the parsed ground-truth value $y_i$ under an absolute floating-point tolerance of $|\hat{y}_i - y_i| \le 10^{-6}$ (\texttt{tolerance = 1e-6}), ensuring that exact fractions (e.g., \texttt{1/3} or \texttt{\textbackslash frac\{1\}\{3\}}), radicals (\texttt{\textbackslash sqrt\{2\}/2}), and decimal representations are scored identically. A math run passes (\texttt{passed = True}) if and only if all $k$ expected values satisfy $|\hat{y}_i - y_i| \le 10^{-6}$.
\end{itemize}
\end{enumerate}

\paragraph{Literature Grounding for Behavioral Outcome Dimensions.}
Rather than defining secondary metrics ad hoc, we identified six behavioral outcome measures by synthesizing three established streams of empirical LLM evaluation literature that correspond to the operational stages of an LLM-integrated software pipeline:
\begin{enumerate}
\tightlist
\item \textbf{Holistic Efficiency and Inference Cost Metrics (\emph{output length} \texttt{completion\_tokens} and \texttt{duration\_seconds}):} Following the Holistic Evaluation of Language Models (HELM) framework \citep{liang2023helm} and recent software-engineering studies on prompt-induced resource overhead \citep{deluca2026evaluating, midolo2026guidelines, cheng2026systemprompts}, accuracy cannot be evaluated in isolation from computational cost. Because autoregressive decoding cost and latency scale linearly with output length, we record \emph{output length} (number of generated output tokens, stored as \texttt{completion\_tokens} in the telemetry logs) to measure verbosity and token-budget consumption, and \emph{duration} (wall-clock response latency in seconds) to capture end-to-end inference time.
\item \textbf{Verifiable Format and Instruction-Following Compliance (\emph{result-only} and \emph{code-only compliance}):} Grounded in verifiable instruction-following benchmarks such as IFEval \citep{zhou2023ifeval} and structured-output adherence studies \citep{liu2024structuredoutput, hua2025promptartifact}, prompts in automated pipelines impose strict negative output constraints (\emph{``no explanations''}). We therefore evaluate \emph{result-only compliance} (for math) and \emph{code-only compliance} (for code) to measure whether the model emits a clean, machine-consumable artifact without conversational preambles or postambles, independent of whether the answer itself is mathematically or functionally correct.
\item \textbf{Execution-Harness and Parser Failure Taxonomies (\emph{empty responses} and \emph{extraction/error modes}):} Drawing on execution-based code evaluation harnesses \citep{chen2021codex, austin2021mbpp, deluca2026evaluating} and structured-generation failure analyses \citep{liu2024structuredoutput, sclar2023spurious}, non-passing outputs in automated pipelines fail at distinct interface boundaries. We therefore track \emph{empty responses} (generation or truncation failures yielding no usable text) and \emph{extraction failures and error modes} (distinguishing parser/syntax failures where no valid Python AST or numeric token can be extracted from downstream runtime exceptions, timeouts, and unit-test assertion failures).
\end{enumerate}

\paragraph{Concrete Operationalization of Behavioral Outcomes and Error Modes.}
To avoid ambiguity in descriptive terms such as ``compact output form,'' ``mostly clean executable code,'' and ``usable output,'' every behavioral outcome and failure category is computed by a deterministic programmatic rule in the analysis pipeline (\texttt{recompute\_manuscript\_tables.py} and \texttt{run\_mixed\_effects\_analysis.py}). Let $R$ denote the raw model response string (\texttt{response\_text}) with character length $L_{\text{resp}} = |R|$, and let $C = \texttt{extract\_python\_code}(R)$ denote the extracted Python code string with character length $L_{\text{code}} = |C|$ and non-code character count $L_{\text{extra}} = \max(0,\, L_{\text{resp}} - L_{\text{code}})$. Table~\ref{tab:outcome-operationalization} specifies the exact mathematical and programmatic decision rules for all primary and secondary outcomes:
\begin{itemize}
\tightlist
\item \textbf{Code-Only Compliance (``mostly clean executable code''):} On code tasks, a response is coded as compliant ($\texttt{result\_only} = 1$) if and only if non-empty code is extracted ($L_{\text{code}} > 0$) and non-code characters constitute strictly less than 15\% of the total response length ($L_{\text{extra}} / L_{\text{resp}} < 0.15$). This 15\% allowance accommodates standard Markdown code-fence delimiters (\texttt{\textasciigrave\textasciigrave\textasciigrave python ... \textasciigrave\textasciigrave\textasciigrave}, typically 13--15 characters) without penalizing fenced code, while classifying any response containing conversational preambles, prose explanations, or post-code walkthroughs as non-compliant ($\texttt{result\_only} = 0$). In the cloud API Python runner, the JSONL logs omit a top-level \texttt{extracted\_code} key (and \texttt{eval\_result.candidate\_code} falls back to raw \texttt{response\_text} when no fence is present); applying \texttt{strict\_extract\_code(response\_text)} recovers the exact same $<15\%$ extra-text compliance indicator across all $1{,}800$ API Python runs ($93.39\%$ overall compliance: $99.89\%$ on \texttt{gemini-flash} and $86.89\%$ on \texttt{deepseek-v4-flash}).
\item \textbf{Result-Only Compliance (``close to the expected compact output form''):} On math tasks, where the system prompt instructs the model to return only the final answer and all ground-truth answers are short numeric values or fractions (1--15 characters, or 15--40 characters when enclosed in minimal JSON/XML tags such as \texttt{\{"answer": 12\}} or \texttt{<answer>12</answer>}), a response is coded as compliant ($\texttt{result\_only} = 1$) if and only if it is non-empty and its total length is strictly under 50 characters ($0 < L_{\text{resp}} < 50$). Responses of 50 characters or more contain multi-step derivations or explanatory prose and are coded as non-compliant ($\texttt{result\_only} = 0$).
\item \textbf{Empty vs.\ Usable Responses:} A response is classified as \emph{empty} ($\texttt{empty\_response} = 1$) when $R\texttt{.strip()} = \texttt{""}$ ($L_{\text{resp}} = 0$, yielding zero characters from the inference endpoint, observed in 14 local Nemotron code runs and 0 API runs). Conversely, a response is operationally \emph{usable} for downstream verification when it is non-empty ($L_{\text{resp}} > 0$) and the extractor recovers a candidate artifact ($L_{\text{code}} > 0$ for code, or at least $k$ numeric tokens $\texttt{len(extracted\_values)} \ge k$ for math) that can be passed to the test runner or numerical comparator.
\item \textbf{Extraction Failures and Error-Mode Taxonomy:} Every non-passing run ($\texttt{passed} = \text{False}$) logs a structured error string in \texttt{errors} that classifies the failure into one of four mutually exclusive stages: (i) \emph{extraction failure} (\texttt{"No code extracted"} when $L_{\text{code}} = 0$, or \texttt{"Expected $k$ numeric value(s), extracted $m$"} when $m < k$); (ii) \emph{syntax/compilation error} (\texttt{SyntaxError} or \texttt{IndentationError} when parsing $C$); (iii) \emph{runtime exception or timeout} (\texttt{NameError} from callable-interface mismatches, \texttt{TypeError}, \texttt{ValueError}, \texttt{ZeroDivisionError}, \texttt{RecursionError}, or \texttt{"Execution timed out (10s)"}); and (iv) \emph{functional/value assertion failure} (\texttt{AssertionError} on one or more of the seven unit tests in code, or $|\hat{y}_i - y_i| > 10^{-6}$ in math).
\end{itemize}

\begin{table}[htbp]
\centering
\caption{Operational definitions and programmatic decision rules for primary and behavioral outcome measures.}\label{tab:outcome-operationalization}
\footnotesize
\setlength{\tabcolsep}{4pt}
\begin{tabular}{@{}p{0.19\linewidth}p{0.14\linewidth}p{0.34\linewidth}p{0.28\linewidth}@{}}
\toprule
\textbf{Outcome Measure} & \textbf{Scale / Type} & \textbf{Exact Operational \& Programmatic Rule} & \textbf{Pipeline Interpretation} \\
\midrule
Pass@1 (\texttt{passed}) & Binary $\{0, 1\}$ & \textbf{Code:} \texttt{subprocess} exit code $= 0$ across all 7 unit tests ($2$ prompt examples $+ 5$ hidden challenge assertions) within 10\,s.\newline \textbf{Math:} $|\hat{y}_i - y_i| \le 10^{-6}$ for all $k$ expected numeric values after 4-step normalization and AST evaluation. & Functional or mathematical correctness of the extracted artifact. \\
\addlinespace[3pt]
Output length (\texttt{completion\_tokens}) & Integer count; $\log(1 + x)$ in LMM/LPM & Generated output token count reported by runtime telemetry (\texttt{eval\_count} in Ollama; \texttt{raw\_usage.completion\_tokens} in APIs). Modeled as $\log(1 + \texttt{completion\_tokens})$. & Output verbosity and autoregressive token-budget cost ($100 \times \hat{\beta} \approx \%$ shift). \\
\addlinespace[3pt]
Response duration (\texttt{duration\_seconds}) & Continuous (s); $\log(1 + x)$ in LMM/LPM & Wall-clock inference time $t_1 - t_0$ via \texttt{time.perf\_counter()} around the synchronous model call (excluding pacing sleeps and test execution). Modeled as $\log(1 + \texttt{duration\_seconds})$. & End-to-end model response latency ($100 \times \hat{\beta} \approx \%$ shift). \\
\addlinespace[3pt]
Code-only compliance (\texttt{result\_only}, code) & Binary $\{0, 1\}$ & $\mathbb{I}\!\left(L_{\text{code}} > 0 \;\wedge\; \frac{\max(0,\, L_{\text{resp}} - L_{\text{code}})}{L_{\text{resp}}} < 0.15\right) = 1$, where $L_{\text{resp}} = |\texttt{response\_text}|$ and $L_{\text{code}} = |\texttt{extracted\_code}|$. & ``Mostly clean executable code'': $\ge 85\%$ of response characters are code (allows Markdown fences, rejects prose commentary). \\
\addlinespace[3pt]
Result-only compliance (\texttt{result\_only}, math) & Binary $\{0, 1\}$ & $\mathbb{I}\!\left(0 < L_{\text{resp}} < 50\right) = 1$, where $L_{\text{resp}} = |\texttt{response\_text}|$ in characters. & ``Compact output form'': concise final answer (plain or minimal JSON/XML tag) without multi-step prose derivation. \\
\addlinespace[3pt]
Empty vs.\ usable response (\texttt{empty\_response}) & Binary $\{0, 1\}$ & \textbf{Empty:} $\texttt{response\_text.strip()} == \texttt{""}$ ($L_{\text{resp}} = 0$).\newline \textbf{Usable:} $L_{\text{resp}} > 0$ and extractor yields a testable artifact ($L_{\text{code}} > 0$ for code; $\ge k$ numbers for math). & Distinguishes complete generation drop-outs from responses consumable by the downstream evaluator. \\
\addlinespace[3pt]
Extraction failures \& error modes (\texttt{errors}) & Categorical diagnostic & \textbf{Extraction failure:} $L_{\text{code}} = 0$ (\texttt{"No code extracted"}) or $< k$ numbers extracted.\newline \textbf{Execution/value failure:} \texttt{SyntaxError}, runtime exception (\texttt{NameError}, \texttt{TypeError}, etc.), 10\,s timeout, or test/value mismatch. & Separates parser/interface breakdowns from algorithmic or numerical errors. \\
\bottomrule
\end{tabular}
\end{table}

Reporting these six literature-grounded dimensions alongside Pass@1 enables us to separate prompt interventions that genuinely alter problem-solving competence from those that shift verbosity, inference latency, instruction adherence, or parser compatibility. Our definition of compliance is deliberately operational rather than cognitive: result-only or code-only compliance does not imply deeper comprehension of the instruction, nor does it guarantee functional correctness. Instead, it measures whether the generated output satisfies the formal interface contract required by downstream software components.

\section{Statistical and Analytical Strategy}\label{statistical-and-analytical-strategy}

Because our factorial experiment evaluates the same 50 base problems across 18 prompt variants and multiple models within each regime, individual pass/fail trials are not statistically independent: runs on the same underlying task share that task's intrinsic difficulty, and runs on the same model share that model's baseline capability. Analyzing all 14,400 runs with naive two-proportion tests as if they were independent coin flips would commit pseudoreplication—artificially shrinking standard errors and inflating false-positive rates. At the same time, software engineers making prompt-design decisions need intuitive, actionable effect sizes (such as percentage-point changes in pass rate and percentage shifts in token cost) alongside formal inferential models that respect binary outcome bounds and repeated-measures clustering.

To meet both requirements transparently, our statistical strategy combines five complementary analytical layers, summarized in Table~\ref{tab:statistical-methods-overview}, together with a dedicated effect-size reporting framework (Section~\ref{effect-size-reporting}) and a design-based power analysis that states the smallest contrast the 50-task design can resolve (Section~\ref{sample-size-and-power}). Each method is selected for a specific inferential or descriptive role, and headline conclusions are drawn only when findings align across layers.

\begin{table}[htbp]
\centering
\caption{Overview of statistical methods used in the study: how each method works, why it was selected, and its methodological limitations.}\label{tab:statistical-methods-overview}
\footnotesize
\setlength{\tabcolsep}{4pt}
\begin{tabular}{@{}p{0.19\linewidth}p{0.27\linewidth}p{0.26\linewidth}p{0.24\linewidth}@{}}
\toprule
\textbf{Analytical Method} & \textbf{How It Works (Accessible Summary)} & \textbf{Why Selected (Role in Study)} & \textbf{Key Limitations} \\
\midrule
1. Descriptive percentage-point contrasts \& raw two-proportion $z$-tests (Sec.~\ref{primary-contrasts}) & Compares raw pass rate under a prompt factor level (e.g., XML) minus the baseline level (e.g., plaintext): $\Delta\text{pp} = 100 \times (\bar{y}_1 - \bar{y}_0)$. & Directly interpretable engineering heuristic: $+4.8\text{ pp}$ means $4.8$ additional passing tasks per $100$ trials before modeling adjustments. & Ignores repeated-task clustering and between-model capability differences; $p$-values are exploratory only. \\
\addlinespace[3pt]
2. Mixed-effects logistic regression (\texttt{BinomialBayesMixedGLM}, Sec.~\ref{fixed-effect-and-mixed-effect-considerations}) & Models log-odds of passing $\log\frac{p}{1-p}$ using fixed effects for the 5 prompt contrasts and models, plus a random intercept $u_t \sim \mathcal{N}(0, \sigma^2_{\text{task}})$ for each of the 50 base tasks. Reports Odds Ratios ($\text{OR} = e^{\beta}$) and Average Marginal Effects ($\text{AME}$ in pp). & \textbf{Primary inferential model} for binary outcomes (\texttt{passed}, \texttt{result\_only}): respects $[0, 1]$ probability bounds and explicitly models repeated measures across the 50 base tasks. & Fitted via variational Bayes (approximate posterior intervals); models additive main effects rather than high-order $3\times 2\times 3$ interactions. \\
\addlinespace[3pt]
3. Linear mixed-effects models (\texttt{MixedLM}, Sec.~\ref{fixed-effect-and-mixed-effect-considerations}) & Fits continuous outcomes $\log(1 + \texttt{tokens})$ and $\log(1 + \texttt{duration})$ with prompt and model fixed effects plus a base-task random intercept $u_t$. & \textbf{Primary inferential model} for skewed continuous behavioral metrics; $100 \times \hat{\beta}$ approximates percentage change in tokens or latency. & Assumes approximately log-normal errors and additive task random intercepts. \\
\addlinespace[3pt]
4. Linear probability models (LPM) with task-clustered SEs (Sec.~\ref{fixed-effect-and-mixed-effect-considerations}, \ref{mediation-and-behavioral-controls}) & Fits OLS regression on binary $y \in \{0,1\}$ with prompt and model fixed effects and Liang--Zeger sandwich standard errors clustered by \texttt{base\_task}. & Frequentist robustness check on mixed-logit AMEs (requires no variational-Bayes or normal random-effect assumptions) and enables additive behavioral-adjustment decomposition ($\hat{\beta}_{\text{total}} - \hat{\beta}_{\text{adj}}$). & Predicted probabilities are not bounded to $[0, 1]$ and binary residuals are heteroskedastic; used strictly as a sensitivity check, not standalone inference. \\
\addlinespace[3pt]
5. Benjamini--Hochberg FDR correction \& dominant-model leave-one-out (Sec.~\ref{multiple-comparisons}, \ref{dominant-model-sensitivity}) & BH step-up procedure adjusts $p$-values across the 60 local model-specific contrasts to control expected false discovery rate at $q=0.05$; leave-one-out re-estimates all models without the top-scoring model. & Prevents false positives from multiple testing across 60 model--factor pairs and ensures pooled effects are not artifacts of a single dominant model (Nemotron or DeepSeek). & BH controls expected false discovery proportion rather than family-wise zero-error probability; API leave-one-out reduces a 2-model panel to 1 model. \\
\bottomrule
\end{tabular}
\end{table}

\subsection{Primary Contrasts and Interpretability Metrics}\label{primary-contrasts}

To analyze prompt factors, we compare each experimental condition against a standard reference baseline: neutral tone, no role, and plain text. This parameterization yields five primary factor contrasts: (i) aggressive vs.\ neutral tone, (ii) polite vs.\ neutral tone, (iii) role framing vs.\ no role, (iv) JSON wrapper vs.\ plaintext, and (v) XML wrapper vs.\ plaintext. Descriptive differences in correctness are expressed in percentage points ($\Delta\text{pp} = 100 \times (\bar{p}_{\text{level}} - \bar{p}_{\text{baseline}})$), whereas behavioral shifts are expressed in generated tokens, response duration (seconds), and percentage changes. Raw two-proportion $z$-test $p$-values are reported alongside descriptive tables for completeness, but because unadjusted two-proportion tests do not account for repeated measures across the 50 base tasks, they are treated strictly as exploratory summaries rather than confirmatory tests.

We report percentage-point differences (and Average Marginal Effects on the percentage-point scale) alongside formal odds ratios because percentage points provide a directly actionable metric for software engineering decisions. Specifically, an Odds Ratio ($\text{OR}$) measures the multiplicative ratio of success odds $\frac{p_1 / (1 - p_1)}{p_0 / (1 - p_0)}$: an $\text{OR}$ of $1.0$ indicates no effect, $\text{OR} > 1.0$ indicates higher odds of passing, $\text{OR} < 1.0$ indicates lower odds, and any $95\%$ interval that includes $1.0$ indicates that the observed difference is statistically compatible with zero effect. However, the practical meaning of a given odds ratio depends nonlinearly on the baseline pass rate (an $\text{OR}$ of $1.5$ corresponds to a $+10\text{ pp}$ gain near a $50\%$ baseline, but less than $+1\text{ pp}$ near a $2\%$ baseline). By contrast, a $+4.0$ percentage-point Average Marginal Effect directly tells an engineer that switching prompt formats yields roughly $4$ additional passing tasks per $100$ trials on average across the benchmark. Throughout the paper, we therefore base interpretation jointly on effect magnitude ($\text{AME}$ in pp and $\text{OR}$), interval width, and cross-model robustness rather than treating $p < 0.05$ as a sufficient standalone criterion.

\subsection{Stratification by Task Kind and Model Group}\label{stratification-by-task-kind-and-model-group}

We estimate all statistical models separately across four strata defined by task domain (code vs.\ math) and deployment regime (local vs.\ API): (1) local code ($n = 5,400$), (2) local math ($n = 5,400$), (3) API Python ($n = 1,800$), and (4) API math ($n = 1,800$). This stratification is essential because prompt-factor effects differ in both magnitude and sign across domains and regimes. Pooling all 14,400 runs into a single regression would average together opposing effects—for example, masking the fact that structured wrappers improve API Python accuracy while JSON wrappers degrade local math accuracy.

Furthermore, stratification acts as a methodological safeguard against overgeneralization. Code generation depends on producing syntactically valid, executable Python code compatible with the regex extractor and unit-test harness, whereas mathematical reasoning depends on numerical derivation and final-answer canonicalization. Likewise, 4-bit quantized local models and commercial cloud APIs differ in parameter scale, task difficulty calibration, default verbosity, and serving infrastructure. Stratifying the analysis treats these regime and domain differences as structural boundaries rather than random noise.

\subsection{Mixed-Effects Models and Clustered Robustness Checks}\label{fixed-effect-and-mixed-effect-considerations}

\paragraph{Why Mixed-Effects Logistic Regression Is Selected for Binary Outcomes.}
Within each stratum, every one of the 50 base tasks is evaluated $18$ times per model (once under each prompt variant in the $3 \times 2 \times 3$ matrix). Because some base tasks are intrinsically easy and others are intrinsically hard, observations sharing the same \texttt{base\_task} are positively correlated. Our primary inferential model for binary Pass@1 (\texttt{passed} $\in \{0, 1\}$) is therefore a generalized linear mixed-effects model (GLMM) with a logit link, estimated separately for each of the four strata:
\begin{equation}
\log\!\left(\frac{P(\texttt{passed}_{itm} = 1)}{1 - P(\texttt{passed}_{itm} = 1)}\right) = \beta_0 + \beta_1\,\text{Aggressive}_i + \beta_2\,\text{Polite}_i + \beta_3\,\text{Role}_i + \beta_4\,\text{JSON}_i + \beta_5\,\text{XML}_i + \gamma_m + u_t, \quad u_t \sim \mathcal{N}(0, \sigma^2_{\text{task}})
\label{eq:mixed-logit}
\end{equation}
where $i$ indexes the prompt variant, $t \in \{1, \dots, 50\}$ indexes the base task, and $m$ indexes the model within the stratum. In Equation~(\ref{eq:mixed-logit}), $\beta_1, \dots, \beta_5$ are the fixed effects of the five prompt contrasts relative to the neutral, no-role, plaintext baseline; $\gamma_m$ are model fixed effects that absorb baseline capability differences between models in the stratum; and $u_t$ is a task-specific random intercept with variance $\sigma^2_{\text{task}}$ that explicitly models the repeated evaluation of the same 50 base tasks. Models are estimated via variational Bayes as implemented in \texttt{statsmodels}' \texttt{BinomialBayesMixedGLM}. For each prompt factor $k \in \{1, \dots, 5\}$, we report:
\begin{enumerate}
\tightlist
\item The \textbf{Odds Ratio} $\text{OR}_k = \exp(\hat{\beta}_k)$ and its $95\%$ posterior interval $\exp(\hat{\beta}_k \pm 1.96\,\text{SD}(\hat{\beta}_k))$.
\item The \textbf{Average Marginal Effect ($\text{AME}_k$)} in percentage points, computed by evaluating the predicted probability difference $100 \times \bigl[\text{expit}(\mathbf{x}_{i(k=1)}^\top \boldsymbol{\beta} + \gamma_m + u_t) - \text{expit}(\mathbf{x}_{i(k=0)}^\top \boldsymbol{\beta} + \gamma_m + u_t)\bigr]$ for every observation in the stratum—averaging over the observed covariate distribution and the estimated task random-intercept distribution $\mathcal{N}(0, \hat{\sigma}^2_{\text{task}})$ across $4,000$ Monte Carlo posterior draws—along with its $2.5\text{th}$ and $97.5\text{th}$ percentile interval.
\end{enumerate}

\paragraph{Linear Mixed-Effects Models for Continuous Behavioral Outcomes.}
The same repeated-measures structure is applied to continuous behavioral outcomes. Because raw token counts and response latencies are right-skewed with long upper tails when models emit verbose explanations, we apply the natural logarithmic transformation $y = \log(1 + x)$ (\texttt{np.log1p}) to both \texttt{completion\_tokens} and \texttt{duration\_seconds} and fit linear mixed-effects models (\texttt{statsmodels} \texttt{MixedLM} via maximum likelihood) with the same fixed effects ($\beta_1, \dots, \beta_5$ and model fixed effects $\gamma_m$) and base-task random intercept $u_t$. Multiplying a log-linear coefficient by $100$ ($100 \times \hat{\beta}_k$) provides an accessible approximation of the percentage change in token volume or latency associated with that prompt factor. Binary result-only/code-only compliance (\texttt{result\_only} $\in \{0, 1\}$) is modeled with the same logistic mixed-effects specification as Equation~(\ref{eq:mixed-logit}) across all four strata (using \texttt{strict\_extract\_code(response\_text)} to recover extracted code in API Python, where compliance is $93.39\%$ overall).

\paragraph{Clustered Linear Probability Models (LPMs) as a Frequentist Robustness Check and Limitations.}
Alongside the mixed-effects models, we estimate linear probability models (OLS regressions of binary $y_{itm} \in \{0, 1\}$ on the five prompt contrasts and model fixed effects) using Liang--Zeger heteroskedasticity- and cluster-robust (``sandwich'') standard errors clustered on \texttt{base\_task} ($G = 50$ task clusters). We include clustered LPMs for two methodological reasons: first, they provide a purely frequentist check on the variational-Bayes mixed-logit AMEs without assuming normally distributed task random intercepts $u_t$ or variational posterior approximations; second, because OLS is a linear operator, LPM coefficients can be directly compared across nested specifications when adding behavioral control variables in Section~\ref{mediation-and-behavioral-controls} (whereas logistic regression coefficients suffer from non-collapsibility, meaning log-odds coefficients shift when covariates are added even without confounding because the latent error variance is rescaled). Conversely, LPMs have well-known limitations for binary data: predicted probabilities are not bounded by $[0, 1]$ and residuals are inherently heteroskedastic ($\text{Var}(y_{itm}) = p_{itm}(1 - p_{itm})$). Accordingly, mixed-effects logistic regression remains our primary inferential basis for all Pass@1 claims, and clustered LPMs are used strictly as a robustness check and additive decomposition tool. We require directional and magnitude agreement across the mixed-logit AME, the clustered LPM, and the raw contrast before treating any prompt effect as a headline finding.

\paragraph{Two-Way Factor-by-Factor and Factor-by-Model Interaction Specifications.}
Because the $3 \times 2 \times 3$ design crosses tone, role, and input format within the same prompt variants, we also extend the task-clustered LPM specification to test two sets of interaction hypotheses that cannot be evaluated in a main-effects-only model:
\begin{enumerate}
\tightlist
\item \emph{Two-way prompt-factor interactions ($\text{Tone} \times \text{Role}$, $\text{Tone} \times \text{Format}$, and $\text{Role} \times \text{Format}$):} Within each stratum, we augment the five main-effect indicators and model fixed effects with all eight two-way interaction terms: $\text{Aggressive} \times \text{Role}$, $\text{Polite} \times \text{Role}$ ($\text{Tone} \times \text{Role}$, $2\text{ df}$); $\text{Aggressive} \times \text{JSON}$, $\text{Aggressive} \times \text{XML}$, $\text{Polite} \times \text{JSON}$, $\text{Polite} \times \text{XML}$ ($\text{Tone} \times \text{Format}$, $4\text{ df}$); and $\text{Role} \times \text{JSON}$, $\text{Role} \times \text{XML}$ ($\text{Role} \times \text{Format}$, $2\text{ df}$). These terms test whether the effect of one prompt dimension depends on another (for example, whether structured JSON/XML wrappers mitigate the verbosity and compliance penalty of polite phrasing in code generation). For each stratum, we report point estimates in percentage points, task-clustered standard errors ($G = 50$ task clusters), $95\%$ confidence intervals, an omnibus joint Wald $\chi^2(8)$ test across all eight interaction parameters, and Benjamini--Hochberg FDR-adjusted $q$-values both within each 8-test stratum ($q_{\text{stratum}}$) and across all $8 \times 4 = 32$ interaction tests ($q_{\text{all32}}$).
\item \emph{Factor-by-model interactions ($\text{Factor} \times \text{Model}$):} To formally test RQ2's question of whether prompt-factor effects are consistent across models within each stratum, we estimate interacted regressions crossing the five prompt contrasts with the $M - 1$ model indicators ($5 \times 5 = 25$ interaction parameters in the 6-model local strata; $5 \times 1 = 5$ in the 2-model API strata) and compute task-clustered joint Wald $F$-tests for the overall $\text{Factor} \times \text{Model}$ interaction as well as the sub-families $\text{Tone} \times \text{Model}$, $\text{Role} \times \text{Model}$, and $\text{Format} \times \text{Model}$.
\end{enumerate}

\subsection{Multiple Comparisons}\label{multiple-comparisons}

When disaggregating prompt effects to individual models, the number of simultaneous statistical tests grows rapidly: testing 5 prompt contrasts across 6 local models and 2 task domains produces $6 \times 5 \times 2 = 60$ local model-specific baseline contrasts. If each test were evaluated at an unadjusted threshold of $\alpha = 0.05$, approximately $60 \times 0.05 = 3$ contrasts would appear ``statistically significant'' purely by chance under the null hypothesis. To prevent false-positive rules from being inferred from noisy per-model subtables, we apply the Benjamini--Hochberg (BH) step-up procedure to control the False Discovery Rate (FDR) at $q = 0.05$ across the full family of 60 local model-specific contrasts (as well as within the separate 30-test code and 30-test math sub-families). Specifically, given $M$ ranked $p$-values $p_{(1)} \le p_{(2)} \le \dots \le p_{(M)}$, the BH procedure identifies the largest rank $r^*$ satisfying $p_{(r)} \le \frac{r}{M} q$ and declares all contrasts with $r \le r^*$ significant after FDR correction.

We selected BH-FDR correction over Bonferroni family-wise error rate correction ($\alpha / 60 = 0.00083$) because Bonferroni is overly conservative when tests share correlated task structures and would mask genuine model-specific vulnerabilities (such as Llama 3.1's sensitivity to JSON wrappers in math). At the same time, we explicitly distinguish three tiers of evidence throughout the manuscript: (i) stratum-level primary mixed-effects contrasts (Table~\ref{tab:mixed-effects-pass}), (ii) BH-FDR-corrected model-specific contrasts, and (iii) uncorrected exploratory contrasts, which are reported for transparency but never promoted to general prompt-engineering rules.

\subsection{Behavioral-Adjustment Analysis}\label{mediation-and-behavioral-controls}

To address RQ4, we conduct a behavioral-adjustment analysis that compares each prompt factor's total correctness contrast ($\hat{\beta}_{\text{total}}$, estimated with model and base-task fixed effects) against its behavior-adjusted contrast ($\hat{\beta}_{\text{adjusted}}$) after sequentially and jointly adding three observable output behaviors to the regression: verbosity ($\log(1 + \texttt{completion\_tokens})$), format adherence (\texttt{result\_only} compliance), and latency ($\log(1 + \texttt{duration\_seconds})$). Using the linear probability decomposition, the difference $\hat{\beta}_{\text{total}} - \hat{\beta}_{\text{adjusted}}$ quantifies how much of an observed pass-rate shift co-varies with measurable changes in output length, cleanliness, and response time. The decomposition is estimated separately for all four strata across all three behavioral controls (for API Python, where \texttt{gemini-flash} is at $99.89\%$ compliance and all non-compliance variance originates from \texttt{deepseek-v4-flash} at $86.89\%$, we report both the tokens-plus-duration adjustment and the full three-control adjustment including \texttt{strict\_extract\_code} compliance).

A critical methodological limitation—which we emphasize throughout the paper—is that output length (in tokens), compliance, and duration are all \emph{post-treatment variables} realized after the prompt framing is applied. Without a separate experimental instrument that manipulates output length independently of prompt wording, conditioning on post-treatment variables cannot identify a formal causal mediation decomposition or prove internal cognitive mechanisms. Nevertheless, stated strictly as an associational decomposition, this analysis provides valuable diagnostic insight for software systems: if a prompt factor's pass-rate penalty disappears after controlling for token length and compliance (as occurs for polite tone in local code), the accuracy difference is empirically tied to verbosity and parser-interface friction; conversely, if a substantial penalty remains after controlling for all three behavioral metrics (as occurs for JSON wrappers in local math), surface verbosity and formatting compliance account for only part of the accuracy drop. Because each cloud API stratum contains only two commercial providers, the API behavioral contrasts describe an average of two models rather than a cross-section of model architectures; we therefore report them as descriptive complements to the local decomposition rather than as independent tests of a mediation mechanism.

\subsection{Dominant-Model Sensitivity}\label{dominant-model-sensitivity}

Finally, in multi-model panels where one model achieves a much higher baseline accuracy than its peers (\texttt{nemotron-3-nano-4b} at $22.8\%$ code and $91.2\%$ math in the local panel; \texttt{deepseek-v4-flash} at $82.4\%$ math in the cloud API panel), pooled regressions can be disproportionately driven by that single model's idiosyncratic prompt preferences. Our dominant-model sensitivity analysis therefore performs a two-stage leave-one-out check: first, we recompute pooled stratum pass rates after excluding the highest-performing model; second, and more importantly, we re-estimate both the mixed-effects logistic model (Equation~\ref{eq:mixed-logit}) and the clustered LPM for all five prompt contrasts on the reduced panel (Table~\ref{tab:dominant-effects}).

This leave-one-dominant-model check directly tests whether headline prompt effects generalize across the remaining models in the tier or depend on a single outlier architecture. Its main structural limitation applies to the cloud API stratum: because the API group evaluates two commercial models, excluding \texttt{deepseek-v4-flash} reduces the leave-one-out panel to a single model (\texttt{gemini-flash}). We report these single-model estimates explicitly because comparing the two-model pooled AME against the \texttt{gemini-flash}-only AME transparently reveals how much of each cloud wrapper or role effect is contributed by each model.

\subsection{Formal Comparison of Model and Prompt Effects}\label{formal-model-prompt-comparison}

To test RQ1 formally rather than by comparing descriptive ranges, we contrast the model and prompt sources of Pass@1 variation within each stratum using four complementary analyses (\texttt{run\_variance\_decomposition.py} in the replication package). Because every stratum is a fully crossed, balanced design (models $\times$ 50 base tasks $\times$ 18 prompt variants, one run per cell), the total sum of squares of the binary outcome partitions exactly into orthogonal components:
\begin{equation}
SS_{\text{total}} = SS_{M} + SS_{T} + SS_{P} + SS_{M \times T} + SS_{M \times P} + SS_{T \times P} + SS_{\text{resid}},
\label{eq:anova-decomposition}
\end{equation}
where $M$ denotes model, $T$ base task, and $P$ the 18-level prompt variant (further split into orthogonal Tone, Role, and Format main effects plus prompt-factor interactions).
\begin{enumerate}
\tightlist
\item \textbf{Variance decomposition (effect size):} We report each component's share of total variation, $\eta^2 = SS_{\text{component}} / SS_{\text{total}}$. Comparing $\eta^2_{M}$ with $\eta^2_{P}$ answers ``how much of the variation in Pass@1 is attributable to which model is used versus which prompt variant is used,'' independently of statistical significance.
\item \textbf{Repeated-measures $F$-tests (significance):} Treating base task as a random blocking factor, each fixed effect is tested against its own interaction with task ($F_M = MS_M / MS_{M \times T}$; $F_P = MS_P / MS_{T \times P}$; analogously for Tone, Role, and Format), and the Model $\times$ Prompt interaction is tested against the residual. This is the classical repeated-measures error structure and avoids pseudoreplication across the 50 tasks.
\item \textbf{Task-cluster bootstrap test of the difference:} To test directly whether model variation exceeds prompt variation, we resample the 50 base tasks with replacement ($B = 2{,}000$, fixed seed) and recompute $\eta^2_{M} - \eta^2_{P}$ and $\eta^2_{M} / \eta^2_{P}$; a $95\%$ percentile interval for the difference that excludes zero indicates that model variation significantly exceeds prompt variation.
\item \textbf{Robustness on the inferential model scales:} We additionally report (i) cluster-robust joint Wald tests of the model block versus the five-contrast prompt block in the task-clustered LPM and (ii) a latent-scale (logit) partition from a fixed-effects logistic regression, comparing the variance of the model and prompt components of the linear predictor against the logistic residual variance $\pi^2/3$.
\end{enumerate}
Because ANOVA on a binary outcome is a linear-probability approximation, the logit-scale partition serves as a check that conclusions do not depend on the outcome scale.

\subsection{Effect-Size Reporting and Magnitude Interpretation}\label{effect-size-reporting}

We report effect sizes explicitly for every headline claim and use $p$-values only to assess compatibility with the null hypothesis, never to rank the magnitude of effects. Five effect-size families are used, each matched to its outcome scale: (i) risk differences and Average Marginal Effects (AMEs) in percentage points for binary Pass@1 and compliance; (ii) odds ratios (ORs) as the multiplicative scale of the mixed-effects logit; (iii) $\eta^2$ variance shares for the model-versus-prompt decomposition; (iv) relative percentage changes ($100 \times \hat{\beta}$) for log-transformed token and duration outcomes; and (v) attenuation percentages, $100 \times (\hat{\beta}_{\text{total}} - \hat{\beta}_{\text{adjusted}}) / \hat{\beta}_{\text{total}}$, for the behavioral-adjustment decomposition.

Because percentage points, odds ratios, and log-scale percentage changes are not directly interchangeable, we interpret magnitude relative to three anchors reported in the same stratum: the between-model Pass@1 spread, the task-level $\eta^2$ share, and the baseline rate. This avoids treating a small but precisely estimated coefficient as practically important, or a large but imprecisely estimated coefficient as unimportant. For binary outcomes, the AME is the primary magnitude metric because it is invariant to the baseline rate, with the OR reported alongside it to make the multiplicative model scale transparent; for continuous outcomes, the relative percentage change is primary, with the untransformed token or second difference also reported where useful.

Concretely, the headline effect sizes are: (i) RQ1: $\eta^2_{\text{Model}} = 7.2$--$41.0\%$ versus $\eta^2_{\text{Prompt}} \le 0.41\%$ in three strata (ratio $19$--$123\times$), with the API Python exception ($0.06\%$ vs.\ $0.29\%$); (ii) RQ2: the eight supported Pass@1 contrasts span AMEs of $1.0$--$5.2$ pp, while 12 of 20 intervals include zero, and the largest single effects are $+5.2$ pp (API Python XML) and $-4.1$ pp (local math JSON); (iii) RQ3: polite tone changes local-code tokens by $+12.40\%$, duration by $+10.76\%$, and compliance by $-7.05$ pp, while JSON changes local-math tokens by $-19.80\%$ and duration by $-11.16\%$; and (iv) RQ4: behavioral adjustment attenuates the local-math JSON penalty by $32.0\%$ and the local-code polite-tone dip by $134.2\%$ (reversing its sign).

We therefore distinguish statistical detectability from practical magnitude throughout: several contrasts reach $p < 0.05$ with AMEs of only $1.0$--$2.2$ pp, whereas some larger AMEs (e.g., API math XML $-2.3$ pp, $p = 0.069$) remain interval-compatible with zero. Statements that an effect is ``larger'' or ``more robust'' in one setting are accompanied by the corresponding effect-size estimate and interval; where the evidence is a significance difference without a meaningful magnitude difference, we say so explicitly.

\subsection{Sample Size and Design-Based Power}\label{sample-size-and-power}

Fifty base tasks per stratum is a modest task sample, and a prompt contrast estimated from it can only be interpreted against the precision the design actually achieved. We therefore report a design-based power analysis rather than leaving the question implicit. Because every prompt contrast is a within-task comparison---each base task is evaluated under all 18 prompt variants under every model in the stratum---the sampling unit that governs precision is the base task, and the task-clustered linear-probability standard errors used in Section~\ref{clustered-linear-probability-robustness-check} are the appropriate variance estimates. For a two-sided test at $\alpha = 0.05$, the minimum detectable effect (MDE) at power $1 - \beta$ is
\begin{equation}
\mathrm{MDE}_{1-\beta} = \left(z_{1-\alpha/2} + z_{1-\beta}\right)\hat{\sigma}_{\text{cluster}},
\label{eq:mde}
\end{equation}
where $\hat{\sigma}_{\text{cluster}}$ is the task-clustered standard error of the contrast in percentage points. Table~\ref{tab:power-analysis} reports this quantity for every stratum, together with the task-level intraclass correlation of Pass@1, the realized design effect, and the number of base tasks that would be required at the same design to detect a 1~pp or 2~pp contrast at $80\%$ power.

Two results deserve emphasis. First, the within-task design is what makes 50 tasks sufficient: task identity is the dominant driver of Pass@1 variance (task-level ICC $0.12$--$0.49$), but because each contrast is estimated within task, this variance is blocked out rather than left in the error term. After task fixed effects, the clustered standard errors are at or below the i.i.d. benchmark in every stratum (variance ratio $0.41$--$1.23$, median $0.59$--$0.90$), so the 18 within-task repetitions act as blocking rather than as a variance penalty and the effective sample size equals or exceeds the raw row count. Second, the two model tiers differ sharply in resolution. The local strata detect contrasts of roughly $1.3$--$2.2$ pp (local code) and $1.9$--$3.1$ pp (local math) at $80\%$ power, and $1.5$--$2.6$ pp and $2.3$--$3.6$ pp at $90\%$ power; the API strata, restricted to two providers and therefore 36 rather than 108 observations per task, detect only $3.9$--$6.0$ pp (API Python) and $3.5$--$4.4$ pp (API math) at $80\%$ power.

The consequences for interpretation are explicit. (i) The eight supported Pass@1 contrasts have achieved power between $0.45$ and $0.96$ at their observed effect sizes (median $0.60$), so most of the effects we report as statistically supported are also effects the design resolves comfortably, while the lowest value belongs to the API Python JSON contrast, which we already treat as suggestive rather than definitive; the twelve null contrasts are accompanied by MDEs that bound how large an undetected effect could plausibly be. (ii) In the API strata, a null contrast of $1$--$2$ pp is not evidence of absence, because such an effect lies well below the detectable floor; this is why we describe API nulls as inconclusive rather than as equivalence. (iii) Extrapolating $\hat{\sigma}_{\text{cluster}} \propto K^{-1/2}$, resolving a 2~pp contrast at $80\%$ power would require approximately $21$--$63$ base tasks in local code, $48$--$122$ in local math, $193$--$449$ in API Python, and $149$--$241$ in API math; resolving $1$~pp would require $84$--$251$, $189$--$487$, $771$--$1{,}795$, and $595$--$961$ tasks respectively. The 50-task design is therefore adequate for the effect sizes this paper claims and explicitly underpowered for sub-percentage-point ones, which is the appropriate scope statement for a study of this size rather than a defect to be concealed.

\begin{table}[htbp]
\centering
\caption{Design-based power analysis: minimum detectable Pass@1 contrast at $80\%$ and $90\%$ power (two-sided $\alpha = 0.05$) for the task-clustered prompt contrasts in each stratum, with the task-level intraclass correlation, realized design effect, and base tasks that would be required to detect a 1~pp or 2~pp contrast at the same design.}\label{tab:power-analysis}
\footnotesize
\setlength{\tabcolsep}{4pt}
\begin{tabular}{@{}lcccccccc@{}}
\toprule
\textbf{Stratum} & \textbf{Observations} & \textbf{Task} & \textbf{Design} & \textbf{Cluster SE} & \textbf{MDE 80\%} & \textbf{MDE 90\%} & \textbf{Tasks for} & \textbf{Tasks for} \\
 & \textbf{(models $\times$ tasks)} & \textbf{ICC} & \textbf{effect} & \textbf{(pp)} & \textbf{(pp)} & \textbf{(pp)} & \textbf{1 pp} & \textbf{2 pp} \\
\midrule
Local code & $5{,}400$ ($6 \times 50$) & $0.19$ & $0.41$--$0.94$ & $0.46$--$0.80$ & $1.29$--$2.24$ & $1.50$--$2.59$ & $84$--$251$ & $21$--$63$ \\
Local math & $5{,}400$ ($6 \times 50$) & $0.12$ & $0.45$--$1.12$ & $0.69$--$1.11$ & $1.94$--$3.12$ & $2.25$--$3.61$ & $189$--$487$ & $48$--$122$ \\
API Python & $1{,}800$ ($2 \times 50$) & $0.49$ & $0.79$--$1.23$ & $1.40$--$2.14$ & $3.93$--$5.99$ & $4.54$--$6.93$ & $771$--$1{,}795$ & $193$--$449$ \\
API math & $1{,}800$ ($2 \times 50$) & $0.22$ & $0.42$--$0.74$ & $1.23$--$1.56$ & $3.45$--$4.38$ & $3.99$--$5.07$ & $595$--$961$ & $149$--$241$ \\
\bottomrule
\end{tabular}

\vspace{2pt}
\begin{minipage}{\linewidth}
\footnotesize
\textit{Notes:} Ranges span the five marginal prompt contrasts (aggressive tone, polite tone, persona role, JSON wrapper, XML wrapper) pooled across all 18 factorial cells within each stratum. Task ICC is the one-way random-effects intraclass correlation of Pass@1 across the 50 base tasks ($108$ observations per task in local strata, $36$ in API strata). Cluster SE is the task-clustered standard error of the linear probability model with model fixed effects (Section~\ref{clustered-linear-probability-robustness-check}). The design effect is $(\hat{\sigma}_{\text{cluster}} / \hat{\sigma}_{\text{i.i.d.}})^2$ from the same model: values below one indicate that within-task blocking reduces rather than inflates the variance. MDE is computed from Equation~\ref{eq:mde}. The task counts are extrapolations under $\hat{\sigma}_{\text{cluster}} \propto K^{-1/2}$, i.e., holding the per-task observation count and residual structure fixed while increasing the number of base tasks. Generated by \texttt{run\_power\_analysis.py} in the replication package.
\end{minipage}
\end{table}

\section{Results}\label{sec:results}

Results are organized sequentially by RQ1--RQ4 and stratified by task domain and model tier (local $n = 10{,}800$ vs.\ cloud API $n = 3{,}600$, interpreted as complementary regimes). Descriptive contrasts are paired throughout with mixed-effects odds ratios, AMEs, $95\%$ intervals, and task-clustered LPM checks (magnitude conventions in Section~\ref{effect-size-reporting}). Table~\ref{tab:key-findings} provides a compact navigation map of all subsections and headline findings.

\begin{table}[htbp]
\centering
\caption{Overview of empirical results organized by Research Question (RQ1--RQ4), mapping each subsection to its integrated descriptive and statistical findings.}\label{tab:key-findings}
\footnotesize
\setlength{\tabcolsep}{4pt}
\begin{tabular}{@{}p{0.15\linewidth}p{0.18\linewidth}p{0.31\linewidth}p{0.32\linewidth}@{}}
\toprule
\textbf{RQ \& Subsection} & \textbf{Stratum \& Scope} & \textbf{Descriptive \& Quantitative Finding} & \textbf{Accompanying Statistical Result} \\
\midrule
\textbf{RQ1: Model vs.\ Prompt Scale} \newline (Sec.~\ref{model-ability-dominates-correctness}) & All 4 strata ($N = 14,400$; 6 local + 2 API models) & Between-model Pass@1 spans $0.4\%$--$22.8\%$ (local code), $5.7\%$--$91.2\%$ (local math), $66.7\%$--$69.0\%$ (API Python), and $24.1\%$--$82.4\%$ (API math). & $\eta^2_{\text{Model}} = 7.2$--$41.0\%$ vs.\ $\eta^2_{\text{Prompt}} \le 0.41\%$ in local code, local math, API math (bootstrap difference CIs exclude 0); API Python exception ($\eta^2 = 0.06\%$ vs.\ $0.29\%$, CI includes 0). \\
\addlinespace[3pt]
\textbf{RQ2 (Part A): Code Correctness} \newline (Sec.~\ref{prompt-factor-effects-on-code-tasks}, \ref{mixed-effects-inferential-results}) & Local code ($n = 5,400$) \& API Python ($n = 1,800$) & \emph{Local code:} Tone/wrappers null ($\Delta \in [-0.5, +0.8]$ pp); role framing $-1.0$ pp ($8.6\%$ vs.\ $9.6\%$). \newline \emph{API Python:} JSON $+3.2$ pp ($68.3\%$ vs.\ $65.2\%$), XML $+4.8$ pp ($70.0\%$ vs.\ $65.2\%$). & \emph{Local code:} Role $\text{OR} = 0.81$ $[0.68, 0.96]$, $\text{AME} = -1.0$ pp ($p = 0.013$); LPM $-1.04$ pp ($p = 0.025$). \newline \emph{API Python:} JSON $\text{OR} = 1.32$, $\text{AME} = +3.4$ pp ($p = 0.022$); XML $\text{OR} = 1.55$, $\text{AME} = +5.2$ pp ($p < 0.001$; LPM $p = 0.006$). \\
\addlinespace[3pt]
\textbf{RQ2 (Part B): Math Correctness} \newline (Sec.~\ref{prompt-factor-effects-on-math-tasks}, \ref{mixed-effects-inferential-results}) & Local math ($n = 5,400$) \& API math ($n = 1,800$) & \emph{Local math:} JSON $-4.2$ pp, XML $-2.2$ pp, Role $+2.2$ pp, Polite $+1.7$ pp; Llama 3.1 JSON $-15.7$ pp, Qwen 2.5 Role $+10.2$ pp. \newline \emph{API math:} Role $-2.8$ pp ($51.9\%$ vs.\ $54.7\%$). & \emph{Local math:} JSON $\text{OR} = 0.62$, $\text{AME} = -4.1$ pp ($p < 0.001$); XML $\text{OR} = 0.78$, $\text{AME} = -2.2$ pp ($p = 0.001$); Role $\text{OR} = 1.28$ ($p < 0.001$); BH-FDR $q < 0.01$ for Llama JSON \& Qwen Role. \newline \emph{API math:} Role $\text{OR} = 0.74$, $\text{AME} = -2.8$ pp ($p = 0.006$; LPM $p = 0.037$). \\
\addlinespace[3pt]
\textbf{RQ3: Output Behavior} \newline (Sec.~\ref{output-behavior-tokens-duration-and-result-only-compliance}) & Tokens, latency, \& compliance in all 4 strata & \emph{Local code:} Polite $+12.40\%$ tokens, $+10.76\%$ latency, $-7.05$ pp compliance; XML $+6.98$ pp compliance. \emph{Local math:} Polite $+10.97\%$ tokens, $+7.51\%$ latency; JSON $-19.80\%$ tokens, $-11.16\%$ latency. \emph{API Python:} Aggressive $+14.20\%$ tokens; wrappers $-6.50\%$/$-5.24\%$ tokens and $+2.88$/$+2.22$ pp compliance ($93.39\%$ overall). \emph{API math:} XML $-17.52\%$ tokens, $+5.13$ pp visible-channel compliance. & MixedLM/logit (Table~\ref{tab:behavior-effects}): local code polite tokens $p = 8.6 \times 10^{-21}$, compliance $\text{AME} = -7.05$ pp ($p = 1.3 \times 10^{-43}$); API Python aggressive tokens $p = 3.7 \times 10^{-5}$, JSON/XML compliance $p \le 0.014$; API math XML tokens $p = 0.023$, compliance $p = 6.5 \times 10^{-7}$. \\
\addlinespace[3pt]
\textbf{RQ4: Behavioral Adjustment} \newline (Sec.~\ref{mediation-and-behavioral-control-results}) & Post-treatment controls (tokens, compliance, latency) in all 4 strata & \emph{Local code:} Polite $-0.50 \to +0.17$ pp. \emph{Local math:} JSON $-4.17 \to -2.84$ pp ($32.0\%$); Role $+2.19 \to +1.38$ pp ($36.7\%$). \emph{API Python:} JSON $+3.17 \to +3.94$ and XML $+4.83 \to +5.43$ pp (tokens+duration) vs.\ $+1.52$ ($52.0\%$) and $+3.53$ pp ($27.0\%$) with audited compliance. \emph{API math:} Role $-2.78 \to -2.59$ pp ($6.6\%$); XML $-2.33 \to -0.72$ pp ($69.0\%$, exploratory). & Linear probability decomposition: local code polite total $p = 0.544$; local math JSON total $p = 7.6 \times 10^{-5}$; API Python XML total $p = 0.012$, tokens+duration adjusted $p = 0.004$, and 3-control adjusted $p = 0.037$; API math role total $p = 0.073$ and adjusted $p = 0.047$. \\
\addlinespace[3pt]
\textbf{Sensitivity \& Validation} \newline (Sec.~\ref{clustered-linear-probability-robustness-check}--\ref{empirical-validation-suite}) & Clustered LPM, leave-one-out, live retries, $k = 3$ repeats & Removing Nemotron preserves JSON penalty ($-4.1 \to -5.2$ pp); $k = 3$ repeats ($14{,}398$ runs) replicate local Python politeness ($-2.56$ pp) and local math JSON ($-6.78$ pp) penalties; live retries gain $+4$--$6$ pp (API) and $+8$--$14$ pp (local). & Full directional agreement across primary, $k = 3$, live pipeline, shared-task, and 15-task non-math SE validations. \\
\bottomrule
\end{tabular}
\end{table}

\subsection{Data Integrity and Experimental Coverage}\label{data-integrity-and-run-coverage}

The final dataset contains 14,400 single-shot runs: 10,800 local-model runs and 3,600 API-model runs. The Consolidated Results Package includes 12 local JSONL files and four API JSONL files, with 900 records in each model-task file. Each base task appears with all 18 prompt variants for each model-task combination. We found no duplicate full run keys and no prompt-metadata join gaps. Empty responses were rare: 14 in the local results, all from Nemotron code outputs, and none in the final API results after transient empty or error responses were retried and successful attempts were retained for evaluation.

\subsection{RQ1: Model Choice Versus Prompt Framing}\label{model-ability-dominates-correctness}

We first describe the model-level pass rates and then test formally whether model variation exceeds prompt variation (Section~\ref{formal-model-prompt-comparison}). Table~\ref{tab:model-level-performance} reports the comprehensive model-level overview across all four task-regime strata—including Pass@1 pass rate, median response duration, and mean output length in tokens—while the corresponding visual comparison of model pass rates is provided in Appendix~\ref{app:model-pass-rate-figure} (Figure~\ref{fig:model_performance_by_task}). As shown in Table~\ref{tab:model-level-performance}, local code pass rates span from $0.4\%$ (\texttt{mistral-7b-instruct}) and $1.1\%$ (\texttt{deepseek-coder-6.7b}) to $22.8\%$ (\texttt{nemotron-3-nano-4b}); local math pass rates span from $5.7\%$ (\texttt{deepseek-coder-6.7b}) and $7.7\%$ (\texttt{mistral-7b-instruct}) to $91.2\%$ (\texttt{nemotron-3-nano-4b}); API Python pass rates range from $66.7\%$ (\texttt{deepseek-v4-flash}) to $69.0\%$ (\texttt{gemini-flash}); and API math pass rates range from $24.1\%$ (\texttt{gemini-flash}) to $82.4\%$ (\texttt{deepseek-v4-flash}).

\begin{longtable}[]{@{}
  >{\raggedright\arraybackslash}p{(\linewidth - 8\tabcolsep) * \real{0.1667}}
  >{\raggedright\arraybackslash}p{(\linewidth - 8\tabcolsep) * \real{0.1667}}
  >{\raggedleft\arraybackslash}p{(\linewidth - 8\tabcolsep) * \real{0.2222}}
  >{\raggedleft\arraybackslash}p{(\linewidth - 8\tabcolsep) * \real{0.2222}}
  >{\raggedleft\arraybackslash}p{(\linewidth - 8\tabcolsep) * \real{0.2222}}@{}}
\caption{Model-level pass rates and output behavior across the four task-regime strata ($900$ runs per row).}\label{tab:model-level-performance}\tabularnewline
\toprule\noalign{}
\begin{minipage}[b]{\linewidth}\raggedright
Task kind
\end{minipage} & \begin{minipage}[b]{\linewidth}\raggedright
Model
\end{minipage} & \begin{minipage}[b]{\linewidth}\raggedleft
Pass rate
\end{minipage} & \begin{minipage}[b]{\linewidth}\raggedleft
Median duration
\end{minipage} & \begin{minipage}[b]{\linewidth}\raggedleft
Mean output tokens
\end{minipage} \\
\midrule\noalign{}
\endfirsthead
\toprule\noalign{}
\begin{minipage}[b]{\linewidth}\raggedright
Task kind
\end{minipage} & \begin{minipage}[b]{\linewidth}\raggedright
Model
\end{minipage} & \begin{minipage}[b]{\linewidth}\raggedleft
Pass rate
\end{minipage} & \begin{minipage}[b]{\linewidth}\raggedleft
Median duration
\end{minipage} & \begin{minipage}[b]{\linewidth}\raggedleft
Mean output tokens
\end{minipage} \\
\midrule\noalign{}
\endhead
\bottomrule\noalign{}
\endlastfoot
local code & nemotron-3-nano-4b & 22.8\% & 17.78s & 1646 \\
local code & qwen2.5-7b & 13.2\% & 8.75s & 316 \\
local code & gemma2-9b & 11.3\% & 5.52s & 175 \\
local code & llama3.1-8b & 5.8\% & 11.16s & 394 \\
local code & deepseek-coder-6.7b & 1.1\% & 9.04s & 491 \\
local code & mistral-7b-instruct & 0.4\% & 9.12s & 366 \\
local math & nemotron-3-nano-4b & 91.2\% & 3.15s & 289 \\
local math & llama3.1-8b & 43.2\% & 5.09s & 274 \\
local math & qwen2.5-7b & 23.6\% & 0.68s & 6 \\
local math & gemma2-9b & 15.2\% & 0.80s & 6 \\
local math & mistral-7b-instruct & 7.7\% & 0.89s & 79 \\
local math & deepseek-coder-6.7b & 5.7\% & 5.35s & 287 \\
API Python & gemini-flash & 69.0\% & 10.30s & 359 \\
API Python & deepseek-v4-flash & 66.7\% & 16.12s & 1352 \\
API math & deepseek-v4-flash & 82.4\% & 9.41s & 1127 \\
API math & gemini-flash & 24.1\% & 3.40s & 15 \\
\end{longtable}

Beyond Pass@1 accuracy, Table~\ref{tab:model-level-performance} reveals substantial between-model differences in default output verbosity and response latency. In local math, for example, \texttt{qwen2.5-7b} and \texttt{gemma2-9b} emit ultra-compact final answers averaging $6$ output tokens with sub-second median durations ($0.68\text{ s}$ and $0.80\text{ s}$), whereas \texttt{nemotron-3-nano-4b}, \texttt{llama3.1-8b}, and \texttt{deepseek-coder-6.7b} generate $274$--$289$ mean output tokens with median durations of $3.15\text{ s}$--$5.35\text{ s}$. A similar pattern holds in cloud API math, where \texttt{gemini-flash} averages $15$ output tokens ($3.40\text{ s}$ median duration, $24.1\%$ pass rate) while \texttt{deepseek-v4-flash} averages $1,127$ output tokens ($9.41\text{ s}$ median duration, $82.4\%$ pass rate). Descriptively, the between-model pass-rate ranges are $22.3$ percentage points in local code, $85.6$ in local math, $2.3$ in API Python, and $58.3$ in API math, whereas the range across all 18 prompt variants is $5.3$, $8.7$, $10.0$, and $11.0$ percentage points, respectively.

\paragraph{Formal Test of Model Versus Prompt Effects.}
Descriptive ranges alone do not establish that model effects exceed prompt effects, so Table~\ref{tab:model-vs-prompt-decomposition} reports the variance decomposition and formal tests specified in Section~\ref{formal-model-prompt-comparison}.

\begin{table}[htbp]
\centering
\caption{Formal comparison of model and prompt effects on Pass@1. $\eta^2$ is the share of total Pass@1 variation (balanced ANOVA, Equation~\ref{eq:anova-decomposition}); bracketed values are $95\%$ task-cluster bootstrap intervals ($B = 2{,}000$). $F$-tests treat base task as a random block ($F_M$ against $M \times T$; $F_P$ against $T \times P$; $F_{M \times P}$ against the residual). Wald tests are task-clustered joint tests in the LPM. Latent shares are logit-scale variance shares of the model and prompt components relative to $\pi^2/3$.}\label{tab:model-vs-prompt-decomposition}
\footnotesize
\setlength{\tabcolsep}{3pt}
\begin{tabular}{@{}p{0.27\linewidth}p{0.17\linewidth}p{0.17\linewidth}p{0.17\linewidth}p{0.17\linewidth}@{}}
\toprule
\textbf{Quantity} & \textbf{Local code} & \textbf{Local math} & \textbf{API Python} & \textbf{API math} \\
\midrule
Model range / prompt-variant range (pp) & $22.3$ / $5.3$ & $85.6$ / $8.7$ & $2.3$ / $10.0$ & $58.3$ / $11.0$ \\
\addlinespace[2pt]
$\eta^2_{\text{Model}}$ (\%) & $7.22$ $[4.45, 11.27]$ & $40.97$ $[33.06, 50.59]$ & $0.06$ $[0.00, 1.08]$ & $34.17$ $[23.52, 47.00]$ \\
$\eta^2_{\text{Prompt}}$, all 18 variants (\%) & $0.27$ $[0.22, 0.67]$ & $0.27$ $[0.19, 0.55]$ & $0.29$ $[0.27, 1.42]$ & $0.41$ $[0.33, 1.13]$ \\
\quad of which Tone / Role / Format (\%) & $0.01$ / $0.03$ / $0.02$ & $0.02$ / $0.06$ / $0.14$ & $0.00$ / $0.01$ / $0.18$ & $0.01$ / $0.08$ / $0.04$ \\
$\eta^2_{\text{Model} \times \text{Prompt}}$ (\%) & $0.55$ & $0.84$ & $0.43$ & $0.43$ \\
$\eta^2_{\text{Task}}$ / $\eta^2_{\text{Model} \times \text{Task}}$ (\%) & $19.65$ / $29.77$ & $12.78$ / $21.94$ & $50.13$ / $7.10$ & $23.80$ / $13.99$ \\
\addlinespace[2pt]
$\eta^2_{\text{Model}} - \eta^2_{\text{Prompt}}$, 95\% CI (pp) & $[4.07, 10.83]$ & $[32.57, 50.27]$ & $[-1.31, 0.39]$ & $[22.90, 46.35]$ \\
$\eta^2_{\text{Model}} / \eta^2_{\text{Prompt}}$, median [95\% CI] & $19.3$ $[9.9, 33.4]$ & $122.9$ $[73.7, 226.2]$ & $0.2$ $[0.0, 1.7]$ & $52.8$ $[27.2, 109.2]$ \\
\addlinespace[2pt]
$F_M$ (df), $p$ & $11.89$ $(5, 245)$, $2.6 \times 10^{-10}$ & $91.49$ $(5, 245)$, $5.0 \times 10^{-54}$ & $0.43$ $(1, 49)$, $0.515$ & $119.70$ $(1, 49)$, $9.4 \times 10^{-15}$ \\
$F_P$ (df), $p$ & $1.93$ $(17, 833)$, $0.013$ & $3.52$ $(17, 833)$, $2.0 \times 10^{-6}$ & $0.65$ $(17, 833)$, $0.849$ & $1.48$ $(17, 833)$, $0.094$ \\
$F_{M \times P}$ (df), $p$ & $0.75$ $(85, 4165)$, $0.959$ & $2.12$ $(85, 4165)$, $1.6 \times 10^{-8}$ & $1.02$ $(17, 833)$, $0.429$ & $1.54$ $(17, 833)$, $0.073$ \\
\addlinespace[2pt]
Clustered LPM Wald: model block $\chi^2$(df), $p$ & $26.7$ $(5)$, $6.6 \times 10^{-5}$ & $546.1$ $(5)$, $9.0 \times 10^{-116}$ & $0.4$ $(1)$, $0.513$ & $119.3$ $(1)$, $9.0 \times 10^{-28}$ \\
Clustered LPM Wald: prompt block $\chi^2$(5), $p$ & $5.8$, $0.331$ & $30.6$, $1.1 \times 10^{-5}$ & $8.9$, $0.114$ & $7.1$, $0.212$ \\
\addlinespace[2pt]
Latent (logit) share: model / prompt (\%) & $40.36$ / $0.14$ & $47.35$ / $0.46$ & $0.09$ / $0.27$ & $35.57$ / $0.23$ \\
\bottomrule
\end{tabular}
\end{table}

Three conclusions follow. First, in three of the four strata (local code, local math, and API math), model variation significantly exceeds prompt variation: model identity explains $7.2\%$, $41.0\%$, and $34.2\%$ of total Pass@1 variation, compared with only $0.27\%$--$0.41\%$ for all 18 prompt variants combined; the bootstrap intervals for $\eta^2_{\text{Model}} - \eta^2_{\text{Prompt}}$ exclude zero in all three ($[4.07, 10.83]$, $[32.57, 50.27]$, and $[22.90, 46.35]$ pp), the model effect is $19$ to $123$ times larger than the prompt effect, and the same ordering holds on the logit scale ($35.6\%$--$47.4\%$ vs.\ $0.14\%$--$0.46\%$ latent variance shares). Second, API Python is an explicit exception: the two cloud models differ by only $2.3$ percentage points ($F_M = 0.43$, $p = 0.515$), and neither model nor prompt variant explains a meaningful share of variation ($\eta^2 = 0.06\%$ vs.\ $0.29\%$; difference interval $[-1.31, 0.39]$ includes zero). In this stratum, the supported structured-wrapper gains (Section~\ref{prompt-factor-effects-on-code-tasks}) are comparable to or larger than the between-model gap, so the ``model dominates'' conclusion does \emph{not} apply. Third, the base task and its interaction with the model account for a large share of variation in every stratum ($\eta^2_{\text{Task}} + \eta^2_{\text{Model} \times \text{Task}} = 34.7\%$--$57.2\%$), exceeding the model share in local code, API Python, and API math and approaching it in local math ($34.7\%$ vs.\ $41.0\%$), confirming that problem difficulty and model--problem fit rival or dominate both experimental factors. Even where prompt effects are statistically detectable ($F_P$ significant in local code and local math), they account for well under $1\%$ of Pass@1 variation.

This establishes the scale for the rest of the results: where models differ substantially in capability, prompt factors are secondary to model and task differences, but still meaningful at the level of specific contrasts (practical guidelines consolidated in Section~\ref{practical-prompt-policy-recommendation}, Table~\ref{tab:actionable-guidelines}).

\subsection{RQ2 (Part A): Prompt-Factor Correctness Effects on Code Tasks}\label{prompt-factor-effects-on-code-tasks}

Table~\ref{tab:code-effects} reports all ten code-generation contrasts in a single ledger: pooled raw Pass@1 rates for the reference and the factor level, the unadjusted percentage-point difference with its raw $p$-value, the mixed-effects logistic odds ratio and average marginal effect (AME), and the task-clustered linear probability model (LPM) estimate. Three regularities structure the discussion. First, no code contrast exceeds $5.2$ pp in absolute AME, and seven of the ten odds-ratio intervals include $1.0$. Second, the two strata differ in which factors matter: in local code only persona role framing produces an interval excluding zero ($\text{AME} = -1.0$ pp), whereas in API Python only the two structured wrappers do ($\text{AME} = +3.4$ and $+5.2$ pp). Third, no factor carries the same sign and significance across both strata, so code-generation prompt effects are stratum-specific rather than universal.

\begin{table}[htbp]
\centering
\caption{Consolidated prompt-factor effects on functional Pass@1 for the two code-generation strata (all ten contrasts). Contrasts are shown against their reference level within each stratum, with pooled raw pass rates, the unadjusted percentage-point difference, and three inferential summaries of the same contrast.}\label{tab:code-effects}
\footnotesize
\setlength{\tabcolsep}{3pt}
\begin{tabular}{@{}lcccccc@{}}
\toprule
\textbf{Contrast (vs.\ reference)} & \textbf{Pass@1 \%} & \textbf{$\Delta$ pp ($p$)} & \textbf{OR [95\% CI]} & \textbf{AME pp [95\% CI]} & \textbf{$p$} & \textbf{LPM pp (SE, $p$)} \\
\midrule
\multicolumn{7}{@{}l}{\emph{Local code} ($n = 5{,}400$; 6 models, 50 base tasks)} \\
Aggressive tone & $9.3 \to 9.3$ & $+0.0$ (1.000) & 0.98 [0.80, 1.20] & $-0.1$ [$-1.0$, $+0.9$] & 0.858 & $+0.00$ (0.80, 1.000) \\
Polite tone & $9.3 \to 8.8$ & $-0.5$ (0.601) & 0.89 [0.72, 1.09] & $-0.5$ [$-1.4$, $+0.4$] & 0.265 & $-0.50$ (0.79, 0.524) \\
Persona role & $9.6 \to 8.6$ & $-1.0$ (0.185) & 0.81 [0.68, 0.96] & $-1.0$ [$-1.7$, $-0.2$] & 0.013 & $-1.04$ (0.46, 0.025) \\
JSON wrapper & $8.8 \to 8.9$ & $+0.2$ (0.860) & 1.01 [0.83, 1.24] & $+0.1$ [$-0.8$, $+1.1$] & 0.900 & $+0.17$ (0.53, 0.753) \\
XML wrapper & $8.8 \to 9.6$ & $+0.8$ (0.387) & 1.15 [0.94, 1.41] & $+0.7$ [$-0.3$, $+1.6$] & 0.164 & $+0.83$ (0.78, 0.287) \\
\midrule
\multicolumn{7}{@{}l}{\emph{API Python} ($n = 1{,}800$; 2 models, 50 base tasks)} \\
Aggressive tone & $67.5 \to 68.0$ & $+0.5$ (0.853) & 1.05 [0.83, 1.34] & $+0.6$ [$-2.4$, $+3.5$] & 0.686 & $+0.50$ (2.14, 0.815) \\
Polite tone & $67.5 \to 68.0$ & $+0.5$ (0.853) & 1.05 [0.83, 1.34] & $+0.6$ [$-2.2$, $+3.4$] & 0.686 & $+0.50$ (1.94, 0.797) \\
Persona role & $68.3 \to 67.3$ & $-1.0$ (0.650) & 0.92 [0.76, 1.11] & $-1.1$ [$-3.4$, $+1.2$] & 0.381 & $-1.00$ (1.40, 0.475) \\
JSON wrapper & $65.2 \to 68.3$ & $+3.2$ (0.244) & 1.32 [1.04, 1.68] & $+3.4$ [$+0.4$, $+6.3$] & 0.022 & $+3.17$ (1.73, 0.068) \\
XML wrapper & $65.2 \to 70.0$ & $+4.8$ (0.074) & 1.55 [1.21, 1.97] & $+5.2$ [$+2.3$, $+8.2$] & $<0.001$ & $+4.83$ (1.77, 0.006) \\
\bottomrule
\end{tabular}

\vspace{2pt}
\begin{minipage}{\linewidth}
\footnotesize
\textit{Notes:} Pass@1 is pooled over models and base tasks within each stratum; the reference levels are neutral tone ($9.3\%$ local, $67.5\%$ API), no persona role ($9.6\%$ local, $68.3\%$ API), and plaintext input ($8.8\%$ local, $65.2\%$ API). $\Delta$ pp is the unadjusted raw difference between the factor level and its reference, with the unclustered raw $p$-value in parentheses. OR, AME, and $p$ are from the mixed-effects logistic model with random intercepts for model and task; AME is the average marginal effect in percentage points. The LPM column reports the linear probability model coefficient in percentage points with task-clustered standard error and $p$-value. Intervals excluding the null value (OR $= 1.0$ or AME $= 0$ pp) appear for three of the ten contrasts: local-code persona role, API-Python JSON, and API-Python XML.
\end{minipage}
\end{table}

Across the local code stratum ($n = 5,400$ runs across 6 models and 50 base tasks), the marginal Pass@1 rates of all tone, role, and format levels stay within the narrow band of $8.6\%$ to $9.6\%$, and the largest contrast between any two levels is $1.0$ pp:
\begin{itemize}
\tightlist
\item \textbf{Social tone (local code):} Neutral and aggressive prompts both achieve a raw pass rate of $9.3\%$ ($\Delta = +0.0$ pp, raw $p = 1.000$; mixed-logit $\text{OR} = 0.98$, $95\%$ interval $[0.80, 1.20]$, $\text{AME} = -0.1$ pp $[-1.0, +0.9]$, $p = 0.858$; clustered LPM $+0.00$ pp, $\text{SE} = 0.80$, $p = 1.000$). Polite prompts achieve $8.8\%$ vs.\ $9.3\%$ for neutral ($\Delta = -0.5$ pp, raw $p = 0.601$; mixed-logit $\text{OR} = 0.89$ $[0.72, 1.09]$, $\text{AME} = -0.5$ pp $[-1.4, +0.4]$, $p = 0.265$; clustered LPM $-0.50$ pp, $\text{SE} = 0.79$, $p = 0.524$). Both tone contrasts are thus bounded within $\pm 0.5$ pp and are compatible with zero on all three specifications.
\item \textbf{Persona role framing (local code):} Adding an expert persona role reduces the pooled pass rate from $9.6\%$ (no role) to $8.6\%$ (role), a raw difference of $\Delta = -1.0$ pp (unclustered raw $p = 0.185$). Once between-model and between-task variance are controlled in the repeated-measures regressions, this is the only local-code contrast whose $95\%$ interval excludes zero: mixed-logit $\text{OR} = 0.81$ $[0.68, 0.96]$, $\text{AME} = -1.0$ pp $[-1.7, -0.2]$, $p = 0.013$, with the clustered LPM agreeing at $-1.04$ pp ($\text{SE} = 0.46$, $p = 0.025$). The effect is nevertheless one percentage point in absolute size and does not survive Benjamini--Hochberg correction across the 30 local-code contrasts ($q_{\text{code}} \ge 0.825$; Section~\ref{clustered-linear-probability-robustness-check}).
\item \textbf{Structured input wrappers (local code):} Relative to plaintext ($8.8\%$), JSON yields $8.9\%$ ($\Delta = +0.2$ pp, raw $p = 0.860$; mixed-logit $\text{OR} = 1.01$ $[0.83, 1.24]$, $\text{AME} = +0.1$ pp $[-0.8, +1.1]$, $p = 0.900$; clustered LPM $+0.17$ pp, $\text{SE} = 0.53$, $p = 0.753$) and XML yields $9.6\%$ ($\Delta = +0.8$ pp, raw $p = 0.387$; mixed-logit $\text{OR} = 1.15$ $[0.94, 1.41]$, $\text{AME} = +0.7$ pp $[-0.3, +1.6]$, $p = 0.164$; clustered LPM $+0.83$ pp, $\text{SE} = 0.78$, $p = 0.287$). Both wrapper intervals include zero and both upper bounds stay below $2$ pp, so no wrapper effect on local-code correctness can be claimed; as shown next, these wrappers act on behavioural outcomes instead.
\end{itemize}

For local code tasks, prompt framing changes operational behavior by substantially larger relative effect sizes than it changes Pass@1 (examined in detail under RQ3 in Section~\ref{output-behavior-tokens-duration-and-result-only-compliance}): polite prompts inflate output tokens by $+12.40\%$ ($p = 8.6 \times 10^{-21}$) and duration by $+10.76\%$ ($p = 4.7 \times 10^{-22}$), and reduce code-only compliance by $\text{AME} = -7.05$ pp ($p = 1.3 \times 10^{-43}$), whereas the corresponding polite-tone Pass@1 AME is only $-0.5$ pp (Table~\ref{tab:mixed-effects-pass}). JSON and XML wrappers compress tokens by $-3.46\%$ and $-4.10\%$ ($p \le 0.009$) and improve code-only compliance by $\text{AME} = +4.27$ pp and $+6.98$ pp ($p < 10^{-16}$), while their Pass@1 AMEs remain small ($+0.1$ and $+0.7$ pp) and interval-compatible with zero.

In the cloud API Python stratum ($n = 1,800$ runs across \texttt{gemini-flash} and \texttt{deepseek-v4-flash}), tone and role framing shift pooled Pass@1 by at most $1.1$ pp and remain compatible with zero, whereas the two structured wrappers produce the largest positive correctness effects observed anywhere in the code experiments ($+3.4$ and $+5.2$ pp AME):
\begin{itemize}
\tightlist
\item \textbf{Structured input wrappers (API Python):} Compared with plaintext ($65.2\%$), the JSON wrapper increases the raw pass rate by $\Delta = +3.2$ pp to $68.3\%$ (unclustered raw $p = 0.244$) and is statistically supported in the primary mixed-effects logistic model ($\text{OR} = 1.32$ $[1.04, 1.68]$, $\text{AME} = +3.4$ pp $[+0.4, +6.3]$, $p = 0.022$) and marginal in the clustered LPM ($+3.17$ pp, $\text{SE} = 1.73$, $p = 0.068$). The XML wrapper increases the raw pass rate by $\Delta = +4.8$ pp to $70.0\%$ (unclustered raw $p = 0.074$) and is supported by both the mixed-effects logistic model ($\text{OR} = 1.55$ $[1.21, 1.97]$, $\text{AME} = +5.2$ pp $[+2.3, +8.2]$, $p < 0.001$) and the clustered LPM ($+4.83$ pp, $\text{SE} = 1.77$, $p = 0.006$), with the two estimates agreeing within $0.4$ pp. Both wrappers also reduce mean output length and request time descriptively—JSON by $-105$ tokens ($-6.50\%$ in log-tokens, $\text{SE} = 3.44\%$, $p = 0.059$) and $-1.55$ s ($-4.45\%$ in log-duration, $p = 0.333$), and XML by $-77$ tokens ($-5.24\%$ in log-tokens, $\text{SE} = 3.44\%$, $p = 0.128$) and $-1.12$ s ($-4.45\%$ in log-duration, $p = 0.333$).
\item \textbf{Social tone and persona role (API Python):} Both aggressive and polite tones match at $68.0\%$ vs.\ $67.5\%$ for neutral ($\Delta = +0.5$ pp, raw $p = 0.853$; mixed-logit $\text{OR} = 1.05$ $[0.83, 1.34]$, $\text{AME} = +0.6$ pp, $p = 0.686$; clustered LPM $+0.50$ pp, $p \ge 0.797$). Role framing yields $67.3\%$ vs.\ $68.3\%$ without role ($\Delta = -1.0$ pp, raw $p = 0.650$; mixed-logit $\text{OR} = 0.92$ $[0.76, 1.11]$, $\text{AME} = -1.1$ pp $[-3.4, +1.2]$, $p = 0.381$; clustered LPM $-1.00$ pp, $\text{SE} = 1.40$, $p = 0.475$).
\end{itemize}

Table~\ref{tab:code-effects} makes the asymmetry explicit: in API Python the XML wrapper raises Pass@1 by $+4.8$ pp raw and $+5.2$ pp in the mixed model $[+2.3, +8.2]$, and JSON by $+3.2$ pp raw and $+3.4$ pp $[+0.4, +6.3]$, whereas in local code the same wrappers move Pass@1 by only $+0.8$ pp raw ($+0.7$ pp $[-0.3, +1.6]$) and $+0.2$ pp raw ($+0.1$ pp $[-0.8, +1.1]$). Locally, wrappers instead act on behaviour: code-only compliance rises by $+4.27$ pp (JSON) and $+6.98$ pp (XML) while output tokens fall by $3.46\%$ and $4.10\%$. This comparison therefore has direct implications for LLM-based coding workflows: structured wrappers (especially XML) improve pass rates when they clarify input field boundaries for capable instruction-tuned coding models under a strict code-only contract, whereas in lower-capacity local code models they primarily improve format compliance and compactness rather than algorithmic pass rates.

\subsection{RQ2 (Part B): Prompt-Factor Correctness Effects on Math Tasks}\label{prompt-factor-effects-on-math-tasks}

In the local strata, mathematical reasoning exhibits larger prompt-framing correctness effect sizes than code generation (largest absolute AME $4.1$ pp vs.\ $1.0$ pp), with several factors reversing sign relative to code generation; in the API strata, the largest accuracy effect is instead $+5.2$ pp for XML in API Python, compared with $-2.8$ pp for role framing in API math. In the local math stratum ($n = 5,400$), structured wrappers significantly degrade Pass@1 accuracy, whereas role framing and polite tone produce modest statistically significant gains:
\begin{itemize}
\tightlist
\item \textbf{Structured input wrappers (local math):} Relative to plaintext ($33.2\%$), the JSON wrapper reduces the pooled pass rate by $\Delta = -4.2$ pp to $29.1\%$ (unclustered raw $p = 0.007$), a penalty confirmed with closely agreeing estimates by both the mixed-effects logistic model ($\text{OR} = 0.62$ $[0.52, 0.72]$, $\text{AME} = -4.1$ pp $[-5.4, -2.7]$, $p < 0.001$) and the clustered LPM ($-4.17$ pp, $\text{SE} = 1.11$, $p < 0.001$). The XML wrapper reduces the pooled pass rate by $\Delta = -2.2$ pp to $31.0\%$ (unclustered raw $p = 0.153$), which is likewise significant once task and model heterogeneity are controlled in the mixed-effects model ($\text{OR} = 0.78$ $[0.66, 0.91]$, $\text{AME} = -2.2$ pp $[-3.4, -0.8]$, $p = 0.001$) and clustered LPM ($-2.22$ pp, $\text{SE} = 1.02$, $p = 0.030$).
\item \textbf{Persona role and social tone (local math):} Role framing increases the pooled pass rate from $30.0\%$ to $32.2\%$ ($\Delta = +2.2$ pp, raw $p = 0.083$; mixed-logit $\text{OR} = 1.28$ $[1.13, 1.45]$, $\text{AME} = +2.1$ pp $[+1.0, +3.2]$, $p < 0.001$; clustered LPM $+2.19$ pp, $\text{SE} = 0.69$, $p = 0.002$). Polite tone increases the pass rate from $30.2\%$ (neutral) to $31.9\%$ ($\Delta = +1.7$ pp, raw $p = 0.280$; mixed-logit $\text{OR} = 1.20$ $[1.03, 1.40]$, $\text{AME} = +1.6$ pp $[+0.2, +3.0]$, $p = 0.019$; clustered LPM $+1.67$ pp, $\text{SE} = 0.79$, $p = 0.035$). Aggressive tone shows no significant effect ($31.2\%$ vs.\ $30.2\%$, $\Delta = +0.9$ pp, raw $p = 0.539$; mixed-logit $\text{OR} = 1.11$ $[0.95, 1.29]$, $\text{AME} = +0.9$ pp $[-0.5, +2.3]$, $p = 0.198$; clustered LPM $+0.94$ pp, $\text{SE} = 0.70$, $p = 0.180$).
\end{itemize}

Disaggregating local math by individual model reveals pronounced model-specific heterogeneity, which we evaluate alongside Benjamini--Hochberg (BH) FDR-adjusted $q$-values across the full 60-test local family ($q_{\text{all}}$) and the 30-test math sub-family ($q_{\text{math}}$):
\begin{itemize}
\tightlist
\item \textbf{Model-specific wrapper penalties:} \texttt{llama3.1-8b} drops by $-15.7$ pp under JSON (unadjusted $p = 0.00011$, $q_{\text{all}} = 0.0066$, $q_{\text{math}} = 0.0033$, significant under both FDR scopes) and by $-9.7$ pp under XML (unadjusted $p = 0.0177$, $q_{\text{all}} = 0.212$, $q_{\text{math}} = 0.106$). \texttt{deepseek-coder-6.7b} drops by $-5.7$ pp under XML (unadjusted $p = 0.0039$, $q_{\text{all}} = 0.078$, $q_{\text{math}} = 0.039$, significant within the 30-test math family) and by $-4.3$ pp under JSON (unadjusted $p = 0.0354$, $q_{\text{all}} = 0.304$, $q_{\text{math}} = 0.177$).
\item \textbf{Model-specific role and tone gains:} \texttt{qwen2.5-7b} exhibits a $+10.2$ pp gain from role framing (unadjusted $p = 0.00030$, $q_{\text{all}} = 0.0091$, $q_{\text{math}} = 0.0045$, significant under both FDR scopes). Across all 60 local model-specific contrasts, seven have unadjusted $p < 0.05$ (including \texttt{deepseek-coder-6.7b} polite math at $+4.7$ pp, $p = 0.0134$, $q_{\text{math}} = 0.101$, and aggressive code at $+2.0$ pp, $p = 0.0327$, $q_{\text{code}} = 0.825$), but only two survive 60-test BH correction (\texttt{llama3.1-8b} JSON math and \texttt{qwen2.5-7b} role math) and three survive 30-test math BH correction (adding \texttt{deepseek-coder-6.7b} XML math), whereas zero of the 30 local code contrasts survive BH correction ($q_{\text{code}} \ge 0.825$).
\end{itemize}

In the cloud API math stratum ($n = 1,800$), persona role framing reverses direction relative to local math and imposes a statistically significant accuracy penalty, while tone and wrappers are non-significant:
\begin{itemize}
\tightlist
\item \textbf{Persona role framing (API math):} Role framing reduces the pass rate by $\Delta = -2.8$ pp from $54.7\%$ (no role) to $51.9\%$ (role; unclustered raw $p = 0.238$) while increasing mean response length by $+83$ tokens. Controlling for the large baseline gap between \texttt{deepseek-v4-flash} ($82.4\%$) and \texttt{gemini-flash} ($24.1\%$), both the mixed-effects logistic model ($\text{OR} = 0.74$ $[0.60, 0.92]$, $\text{AME} = -2.8$ pp $[-4.8, -0.8]$, $p = 0.006$) and the clustered LPM ($-2.78$ pp, $\text{SE} = 1.33$, $p = 0.037$) confirm a significant negative role effect.
\item \textbf{Social tone and structured wrappers (API math):} Relative to neutral ($53.5\%$), aggressive tone yields $53.8\%$ ($\Delta = +0.3$ pp, raw $p = 0.908$; mixed-logit $\text{OR} = 1.05$ $[0.81, 1.38]$, $\text{AME} = +0.5$ pp, $p = 0.697$; clustered LPM $+0.33$ pp, $p = 0.818$) and polite tone yields $52.5\%$ ($\Delta = -1.0$ pp, raw $p = 0.729$; mixed-logit $\text{OR} = 0.91$ $[0.70, 1.18]$, $\text{AME} = -0.9$ pp, $p = 0.476$; clustered LPM $-1.00$ pp, $p = 0.417$). Relative to plaintext ($54.5\%$), JSON yields $53.2\%$ ($\Delta = -1.3$ pp, raw $p = 0.643$; mixed-logit $\text{OR} = 0.87$ $[0.67, 1.14]$, $\text{AME} = -1.2$ pp, $p = 0.325$; clustered LPM $-1.33$ pp, $p = 0.310$) and XML yields $52.2\%$ ($\Delta = -2.3$ pp, raw $p = 0.418$; mixed-logit $\text{OR} = 0.78$ $[0.60, 1.02]$, $\text{AME} = -2.3$ pp $[-4.9, +0.1]$, $p = 0.069$; clustered LPM $-2.33$ pp, $\text{SE} = 1.56$, $p = 0.136$).
\end{itemize}

Together, Sections~\ref{prompt-factor-effects-on-code-tasks} and~\ref{prompt-factor-effects-on-math-tasks} answer RQ2 by demonstrating that prompt-framing effects on Pass@1 are domain- and model-contingent rather than universal: structured wrappers aid cloud Python code generation ($+3.4$ to $+5.2$ pp AME) but degrade local mathematical reasoning ($-2.2$ to $-4.1$ pp AME), and persona roles help local math ($+2.1$ pp AME, driven by Qwen 2.5) but hurt local code ($-1.0$ pp AME) and cloud math ($-2.8$ pp AME).

\subsection{RQ2 (Synthesis): Multi-Stratum Mixed-Effects Inferential Summary}\label{mixed-effects-inferential-results}
Table~\ref{tab:mixed-effects-pass} consolidates the primary mixed-effects logistic regression estimates across all four strata (20 factor--stratum contrasts) so that the cross-domain sign reversals reported in Sections~\ref{prompt-factor-effects-on-code-tasks} and~\ref{prompt-factor-effects-on-math-tasks} can be compared side-by-side on a unified odds-ratio and percentage-point AME scale.

\begin{longtable}[]{@{}llrrr@{}}
\caption{Mixed-effects logistic results for pass/fail by stratum: odds ratios and average marginal effects (AME) in percentage points, with 95\% intervals and posterior p-values.}\label{tab:mixed-effects-pass}\tabularnewline
\toprule\noalign{}
Stratum & Contrast & Odds ratio [95\% int.] & AME pp [95\% int.] & p \\
\midrule\noalign{}
\endfirsthead
\toprule\noalign{}
Stratum & Contrast & Odds ratio [95\% int.] & AME pp [95\% int.] & p \\
\midrule\noalign{}
\endhead
\bottomrule\noalign{}
\endlastfoot
API Python & Aggressive tone & 1.05 [0.83, 1.34] & +0.6 [-2.4, +3.5] & 0.686 \\
API Python & Polite tone & 1.05 [0.83, 1.34] & +0.6 [-2.2, +3.4] & 0.686 \\
API Python & Role framing & 0.92 [0.76, 1.11] & -1.1 [-3.4, +1.2] & 0.381 \\
API Python & JSON wrapper & 1.32 [1.04, 1.68] & +3.4 [+0.4, +6.3] & 0.022 \\
API Python & XML wrapper & 1.55 [1.21, 1.97] & +5.2 [+2.3, +8.2] & \textless{}0.001 \\
API math & Aggressive tone & 1.05 [0.81, 1.38] & +0.5 [-2.0, +3.1] & 0.697 \\
API math & Polite tone & 0.91 [0.70, 1.18] & -0.9 [-3.3, +1.6] & 0.476 \\
API math & Role framing & 0.74 [0.60, 0.92] & -2.8 [-4.8, -0.8] & 0.006 \\
API math & JSON wrapper & 0.87 [0.67, 1.14] & -1.2 [-3.8, +1.4] & 0.325 \\
API math & XML wrapper & 0.78 [0.60, 1.02] & -2.3 [-4.9, +0.1] & 0.069 \\
Local code & Aggressive tone & 0.98 [0.80, 1.20] & -0.1 [-1.0, +0.9] & 0.858 \\
Local code & Polite tone & 0.89 [0.72, 1.09] & -0.5 [-1.4, +0.4] & 0.265 \\
Local code & Role framing & 0.81 [0.68, 0.96] & -1.0 [-1.7, -0.2] & 0.013 \\
Local code & JSON wrapper & 1.01 [0.83, 1.24] & +0.1 [-0.8, +1.1] & 0.900 \\
Local code & XML wrapper & 1.15 [0.94, 1.41] & +0.7 [-0.3, +1.6] & 0.164 \\
Local math & Aggressive tone & 1.11 [0.95, 1.29] & +0.9 [-0.5, +2.3] & 0.198 \\
Local math & Polite tone & 1.20 [1.03, 1.40] & +1.6 [+0.2, +3.0] & 0.019 \\
Local math & Role framing & 1.28 [1.13, 1.45] & +2.1 [+1.0, +3.2] & \textless{}0.001 \\
Local math & JSON wrapper & 0.62 [0.52, 0.72] & -4.1 [-5.4, -2.7] & \textless{}0.001 \\
Local math & XML wrapper & 0.78 [0.66, 0.91] & -2.2 [-3.4, -0.8] & 0.001 \\
\end{longtable}

The synthesis in Table~\ref{tab:mixed-effects-pass} reinforces two central conclusions. First, the direction and statistical significance of every headline contrast align between the mixed-effects models and the clustered LPMs reported in Sections~\ref{prompt-factor-effects-on-code-tasks}--\ref{prompt-factor-effects-on-math-tasks} and Section~\ref{clustered-linear-probability-robustness-check}: structured wrappers improve API Python pass rates for both JSON ($\text{OR} = 1.32$, $p = 0.022$) and XML ($\text{OR} = 1.55$, $p < 0.001$), whereas in local math both JSON ($\text{OR} = 0.62$, $p < 0.001$) and XML ($\text{OR} = 0.78$, $p = 0.001$) reduce correctness, while role framing ($\text{OR} = 1.28$, $p < 0.001$) and polite tone ($\text{OR} = 1.20$, $p = 0.019$) provide small local-math gains. Conversely, role framing reduces correctness in API math ($\text{OR} = 0.74$, $p = 0.006$) and local code ($\text{OR} = 0.81$, $p = 0.013$), and social tone has no supported effect in either code stratum.

\paragraph{Moderation of Effect Sizes and Null-Compatible Odds Ratio Intervals.}
Inspecting the $95\%$ intervals and Average Marginal Effects in Table~\ref{tab:mixed-effects-pass} underscores the need for a moderated interpretation of prompt-framing effects on Pass@1. First, \textbf{12 of the 20 stratum-level contrasts have $95\%$ odds-ratio intervals that include $1.0$} (and AME intervals that straddle $0.0$ pp): these include all aggressive-tone contrasts across all four strata, polite tone in API Python, API math, and local code, role framing in API Python, JSON and XML wrappers in API math, and JSON and XML wrappers in local code. For three-fifths of the tested factor--stratum pairs, the empirical data are therefore compatible with zero accuracy effect. Second, even among the \textbf{8 contrasts whose $95\%$ intervals exclude $\text{OR} = 1.0$}, the absolute magnitudes of the Average Marginal Effects are small to moderate, spanning $1.0$ to $5.2$ percentage points:
\begin{itemize}
\tightlist
\item \emph{Moderate effects ($2.8$ to $5.2$ pp):} API Python XML wrapper ($+5.2$ pp, $95\%$ interval $[+2.3, +8.2]$, OR $1.55$), Local math JSON wrapper ($-4.1$ pp $[-5.4, -2.7]$, OR $0.62$), API Python JSON wrapper ($+3.4$ pp $[+0.4, +6.3]$, OR $1.32$), and API math Role framing ($-2.8$ pp $[-4.8, -0.8]$, OR $0.74$).
\item \emph{Small effects ($1.0$ to $2.2$ pp):} Local math XML wrapper ($-2.2$ pp $[-3.4, -0.8]$, OR $0.78$), Local math Role framing ($+2.1$ pp $[+1.0, +3.2]$, OR $1.28$), Local math Polite tone ($+1.6$ pp $[+0.2, +3.0]$, OR $1.20$), and Local code Role framing ($-1.0$ pp $[-1.7, -0.2]$, OR $0.81$).
\end{itemize}
Given the 50-base-task sample size per domain, statistically significant shifts of $1.0$ to $2.2$ percentage points represent minor accuracy perturbations that are easily outweighed by model choice or task difficulty. Consequently, our conclusions do not claim that prompt framing broadly transforms Pass@1 accuracy; rather, correctness sensitivity is concentrated in specific format--domain interactions (structured wrappers helping cloud Python extraction by $+3.4$ to $+5.2$ pp while hurting local math accuracy by $-2.2$ to $-4.1$ pp; practical guidelines consolidated in Section~\ref{practical-prompt-policy-recommendation}, Table~\ref{tab:actionable-guidelines}). Because the mixed-effects logistic models are fitted using variational Bayes, substantive claims are based on agreement across mixed-logit AMEs, intervals, clustered LPMs, and dominant-model sensitivity checks rather than on any single $p$-value.

\subsection{RQ3: Prompt-Factor Effects on Output Behavior (Tokens, Latency, and Compliance)}\label{output-behavior-tokens-duration-and-result-only-compliance}

Addressing RQ3, prompt framing alters observable operational behavior—output token length, request latency, and code-only/result-only compliance—with larger relative effect sizes and more consistent signs across strata than it alters Pass@1 correctness. Table~\ref{tab:behavior-effects} reports all twenty factor-by-stratum behavioral contrasts for the \emph{four} strata on a common scale: log-linear percent changes for token volume and duration, and mixed-effects logistic average marginal effects in percentage points for compliance. Two audit-verified measurement facts are stated up front rather than left implicit. First, the cloud API strata are modeled alongside the local strata across all three behavioral outcomes. Second, our post-hoc compliance audit (\texttt{audit\_compliance\_metrics.py}, Section~\ref{construct-validity}) established that (a) recovering API Python code-only compliance via \texttt{strict\_extract\_code(response\_text)} yields $93.39\%$ overall compliance ($1{,}681 / 1{,}800$ runs; $99.94\%$ agreement with strict \texttt{ast.parse} acceptance at $93.44\%$), with all non-compliance variance originating from \texttt{deepseek-v4-flash} ($86.89\%$) while \texttt{gemini-flash} sits at $99.89\%$; and (b) for math tasks, non-reasoning models log true raw text, whereas reasoning models (\texttt{nemotron-3-nano-4b} locally and \texttt{deepseek-v4-flash} on API) emit hidden thinking tokens counted in \texttt{completion\_tokens} while returning only concise answers in \texttt{response\_text} ($<20$ characters on $97.4\%$ and $61.9\%$ of runs, respectively), so math result-only compliance measures visible-channel conciseness rather than total reasoning verbosity.
\begin{itemize}
\tightlist
\item \textbf{Local code ($n = 5,400$):} Polite prompts increase output length by $+12.40\%$ (MixedLM $\hat{\beta} = +0.1240$, $\text{SE} = 0.0133$, $p = 8.6 \times 10^{-21}$) and request duration by $+10.76\%$ ($\hat{\beta} = +0.1076$, $\text{SE} = 0.0111$, $p = 4.7 \times 10^{-22}$), while reducing code-only compliance by $-7.05$ pp ($\text{AME}$, $p = 1.3 \times 10^{-43}$). Role framing increases tokens by $+3.17\%$ ($\text{SE} = 1.08\%$, $p = 0.0034$) and duration by $+2.41\%$ ($\text{SE} = 0.91\%$, $p = 0.0081$) and reduces compliance by $-1.18$ pp ($p = 0.0048$). Aggressive tone has no effect on tokens ($+0.55\%$, $p = 0.678$) or duration ($+0.52\%$, $p = 0.639$) but reduces compliance by $-1.48$ pp ($p = 0.0048$). Structured wrappers compress responses and improve compliance: JSON $-3.46\%$ tokens ($\text{SE} = 1.33\%$, $p = 0.0090$), $-3.32\%$ duration ($p = 0.0028$), $+4.27$ pp compliance ($p = 4.1 \times 10^{-17}$); XML $-4.10\%$ tokens ($p = 0.0020$), $-3.57\%$ duration ($p = 0.0013$), $+6.98$ pp compliance ($p = 6.3 \times 10^{-41}$).
\item \textbf{Local math ($n = 5,400$):} Polite tone increases tokens by $+10.97\%$ (MixedLM $\text{SE} = 2.57\%$, $p = 2.0 \times 10^{-5}$) and duration by $+7.51\%$ ($\text{SE} = 1.47\%$, $p = 3.5 \times 10^{-7}$) and reduces result-only compliance by $-1.36$ pp ($\text{AME}$, $p = 0.0053$). Role framing increases tokens by $+8.86\%$ ($\text{SE} = 2.10\%$, $p = 2.4 \times 10^{-5}$) and duration by $+3.98\%$ ($\text{SE} = 1.20\%$, $p = 0.00094$) and reduces compliance by $-1.09$ pp ($p = 0.0073$). Aggressive tone marginally increases tokens ($+4.39\%$, $p = 0.087$), significantly increases duration ($+4.24\%$, $p = 0.0040$), and reduces compliance by $-1.47$ pp ($p = 0.0025$). The JSON wrapper compresses local math outputs sharply—$-19.80\%$ tokens ($\text{SE} = 2.57\%$, $p = 1.3 \times 10^{-14}$), $-11.16\%$ duration ($p = 3.6 \times 10^{-14}$)—with a non-significant $+0.72$ pp compliance shift ($p = 0.135$), whereas XML produces only minor, non-significant shifts ($-1.96\%$ tokens, $p = 0.445$; $-1.14\%$ duration, $p = 0.438$; $-0.36$ pp compliance, $p = 0.457$).
\item \textbf{API Python ($n = 1,800$):} Token effects are comparable in relative size to the local strata. Aggressive tone increases output tokens by $+14.20\%$ ($\text{SE} = 3.44\%$, $p = 3.7 \times 10^{-5}$)—the largest token inflation in the study—while polite tone increases them by $+6.49\%$ ($\text{SE} = 3.44\%$, $p = 0.059$) and role framing changes them by $-1.53\%$ ($\text{SE} = 2.81\%$, $p = 0.587$). Latency is largely insensitive: every factor moves duration by at most $6.1\%$ ($|\beta| \le 0.061$, all $p \ge 0.105$). Both wrappers reduce tokens descriptively—JSON by $-6.50\%$ ($p = 0.059$) and XML by $-5.24\%$ ($p = 0.128$)—and both reduce duration by $-4.45\%$ ($p = 0.333$). For audited code-only compliance ($93.39\%$ overall, driven by \texttt{deepseek-v4-flash} at $86.89\%$), structured wrappers improve compliance by $+2.88$ pp under JSON ($p = 0.0020$) and $+2.22$ pp under XML ($p = 0.0137$), role framing reduces compliance by $-3.61$ pp ($p = 2.0 \times 10^{-7}$), and polite and aggressive tones increase compliance on \texttt{deepseek-v4-flash} by $+3.86$ pp ($p < 0.001$) and $+2.88$ pp ($p = 0.0017$).
\item \textbf{API math ($n = 1,800$):} The XML wrapper compresses API math responses significantly in token volume ($-17.52\%$, $\text{SE} = 7.70\%$, $p = 0.023$) and marginally in duration ($-9.76\%$, $p = 0.097$); aggressive tone shows the opposite, non-significant tendency ($+12.78\%$ tokens, $p = 0.097$), and JSON reduces tokens by $-5.20\%$ ($p = 0.500$) and duration by $-10.71\%$ ($p = 0.069$). For visible-channel result-only compliance (noting that \texttt{deepseek-v4-flash} emits hidden thinking tokens in \texttt{completion\_tokens}), JSON raises visible conciseness by $+5.80$ pp ($p = 3.4 \times 10^{-8}$) and XML by $+5.13$ pp ($p = 6.5 \times 10^{-7}$), whereas role framing lowers it by $-3.83$ pp ($p = 2.4 \times 10^{-6}$) and polite tone raises it marginally by $+1.98$ pp ($p = 0.054$); aggressive tone is null ($+0.84$ pp, $p = 0.401$).
\end{itemize}

Figure~\ref{fig:behavior-effects} visualizes the ten local behavioral contrasts across Panel (a) (percentage shifts in output length in tokens and duration) and Panel (b) (percentage-point shifts in code-only/result-only compliance), illustrating how polite phrasing inflates verbosity and latency while JSON substantially compacts local math responses; the corresponding ten cloud API contrasts are reported numerically in Table~\ref{tab:behavior-effects} rather than plotted, because the six-model local panel supplies the architectural cross-section that a figure of this kind is meant to convey, whereas each API panel contains only two providers.

\begin{center}
\begin{minipage}{0.90\linewidth}
\centering

\begin{minipage}[b]{0.74\linewidth}
\centering
\includegraphics[width=\linewidth]{behavior_effects_tokens_duration.png}
\textbf{(a)} Token and duration effects, percent
\end{minipage}

\vspace{0.6em}

\begin{minipage}[b]{0.74\linewidth}
\centering
\includegraphics[width=\linewidth]{behavior_effects_compliance.png}
\textbf{(b)} Result-only/code-only compliance effects, percentage points
\end{minipage}

\captionof{figure}{Prompt-factor effects on output behavior for local code and local math. Panel (a): effects on output length (tokens) and duration, in percent. Panel (b): effects on result-only/code-only compliance, in percentage points. Bars are mixed-effects point estimates from the same models as Table~\ref{tab:behavior-effects} (log-linear percent change in Panel (a); average marginal effect in Panel (b)); whiskers are $95\%$ confidence intervals, and every bar is labelled with its point estimate.}
\label{fig:behavior-effects}
\end{minipage}
\end{center}

Together, the twenty contrasts in Table~\ref{tab:behavior-effects} answer RQ3 across all four strata: prompt framing reliably governs how models respond (verbosity, latency, and format adherence) even when Pass@1 correctness is unchanged, making behavioral telemetry essential for evaluating automated coding and reasoning pipelines (practical guidelines consolidated in Section~\ref{practical-prompt-policy-recommendation}, Table~\ref{tab:actionable-guidelines}). The dissociation is not a local-model artifact. In API Python, aggressive tone inflates output tokens by $+14.20\%$ while its Pass@1 AME is only $+0.6$ pp ($p = 0.686$); in API math, the XML wrapper compresses tokens by $-17.52\%$ and raises visible-channel result-only compliance by $+5.13$ pp while its Pass@1 AME remains negative ($-2.3$ pp, $p = 0.069$). In both cloud strata, therefore, the measured behavioral response is larger and better identified than the corresponding accuracy effect.

\begin{table}[htbp]
\centering
\caption{Prompt-factor effects on operational output behavior across all four strata: relative change in output tokens, relative change in request duration, and change in code-only/result-only compliance, each with its model-based $p$-value.}\label{tab:behavior-effects}
\footnotesize
\setlength{\tabcolsep}{4pt}
\begin{tabular}{@{}llccc@{}}
\toprule
\textbf{Stratum} & \textbf{Factor (vs.\ reference)} & \textbf{$\Delta$ tokens \% ($p$)} & \textbf{$\Delta$ duration \% ($p$)} & \textbf{$\Delta$ compliance pp ($p$)} \\
\midrule
\multicolumn{5}{@{}l}{\emph{Local code} ($n = 5{,}400$; 6 models, 50 base tasks)} \\
  & Aggressive tone & $+0.55$ (0.678) & $+0.52$ (0.639) & $-1.48$ (0.005) \\
  & Polite tone & $+12.40$ ($<$0.001) & $+10.76$ ($<$0.001) & $-7.05$ ($<$0.001) \\
  & Persona role & $+3.17$ (0.003) & $+2.41$ (0.008) & $-1.18$ (0.005) \\
  & JSON wrapper & $-3.46$ (0.009) & $-3.32$ (0.003) & $+4.27$ ($<$0.001) \\
  & XML wrapper & $-4.10$ (0.002) & $-3.57$ (0.001) & $+6.98$ ($<$0.001) \\
\midrule
\multicolumn{5}{@{}l}{\emph{Local math} ($n = 5{,}400$; 6 models, 50 base tasks)} \\
  & Aggressive tone & $+4.39$ (0.087) & $+4.24$ (0.004) & $-1.47$ (0.003) \\
  & Polite tone & $+10.97$ ($<$0.001) & $+7.51$ ($<$0.001) & $-1.36$ (0.005) \\
  & Persona role & $+8.86$ ($<$0.001) & $+3.98$ (0.001) & $-1.09$ (0.007) \\
  & JSON wrapper & $-19.80$ ($<$0.001) & $-11.16$ ($<$0.001) & $+0.72$ (0.135) \\
  & XML wrapper & $-1.96$ (0.445) & $-1.14$ (0.438) & $-0.36$ (0.457) \\
\midrule
\multicolumn{5}{@{}l}{\emph{API Python} ($n = 1{,}800$; 2 models, 50 base tasks)$^{a}$} \\
  & Aggressive tone & $+14.20$ ($<$0.001) & $+2.35$ (0.609) & $+2.88$ (0.002) \\
  & Polite tone & $+6.49$ (0.059) & $+1.06$ (0.817) & $+3.86$ ($<$0.001) \\
  & Persona role & $-1.53$ (0.587) & $-6.09$ (0.105) & $-3.61$ ($<$0.001) \\
  & JSON wrapper & $-6.50$ (0.059) & $-4.45$ (0.333) & $+2.88$ (0.002) \\
  & XML wrapper & $-5.24$ (0.128) & $-4.45$ (0.333) & $+2.22$ (0.014) \\
\midrule
\multicolumn{5}{@{}l}{\emph{API math} ($n = 1{,}800$; 2 models, 50 base tasks)$^{b}$} \\
  & Aggressive tone & $+12.78$ (0.097) & $+0.73$ (0.902) & $+0.84$ (0.401) \\
  & Polite tone & $+1.18$ (0.878) & $-1.99$ (0.736) & $+1.98$ (0.054) \\
  & Persona role & $+2.80$ (0.656) & $-1.06$ (0.825) & $-3.83$ ($<$0.001) \\
  & JSON wrapper & $-5.20$ (0.500) & $-10.71$ (0.069) & $+5.80$ ($<$0.001) \\
  & XML wrapper & $-17.52$ (0.023) & $-9.76$ (0.097) & $+5.13$ ($<$0.001) \\
\bottomrule
\end{tabular}

\vspace{2pt}
\begin{minipage}{\linewidth}
\footnotesize
\textit{Notes:} Token and duration contrasts are log-linear mixed-model coefficients expressed as relative percentage changes ($100 \times \hat{\beta}$), with random intercepts for model and task. Compliance contrasts are mixed-effects logistic average marginal effects in percentage points for code-only (local code, API Python) and result-only (local math, API math) compliance. $^{a}$Audited API Python code-only compliance is computed directly from \texttt{strict\_extract\_code(response\_text)} ($93.39\%$ overall across $1{,}800$ runs; $99.94\%$ agreement with strict \texttt{ast.parse} acceptance at $93.44\%$; non-compliance variance originates from \texttt{deepseek-v4-flash} at $86.89\%$ while \texttt{gemini-flash} is at $99.89\%$). $^{b}$For math, cross-checking against our full-raw $k = 3$ dataset confirms that non-reasoning models logged true raw text, whereas reasoning models (\texttt{deepseek-v4-flash}, \texttt{nemotron-3-nano-4b}) emit hidden thinking tokens counted in \texttt{completion\_tokens} while returning only concise answers in \texttt{response\_text}; math result-only compliance measures visible-channel conciseness rather than total reasoning verbosity.
\end{minipage}
\end{table}

\subsection{RQ4: Behavioral-Adjustment Analysis of Correctness Contrasts}\label{mediation-and-behavioral-control-results}

Addressing RQ4, our behavioral-adjustment analysis compares each prompt factor's unadjusted total correctness contrast ($\hat{\beta}_{\text{total}}$, estimated with model and base-task fixed effects) against its behavior-adjusted contrast ($\hat{\beta}_{\text{adjusted}}$) after controlling sequentially and jointly for log output tokens, code-only/result-only compliance, and log duration. Because these behavioral variables are realized after the prompt is applied, this decomposition measures post-treatment associational co-variation rather than causal mediation. The decomposition is estimated across all four strata:
\begin{itemize}
\tightlist
\item \textbf{Local code behavioral adjustment ($n = 5,400$):} In local code, the unadjusted polite-tone contrast is $\hat{\beta}_{\text{total}} = -0.50$ pp ($p = 0.544$). Controlling for log tokens alone shifts the polite contrast to $+0.33$ pp; controlling for duration alone shifts it to $+0.41$ pp; and controlling jointly for all three behavioral metrics shifts the adjusted coefficient to $\hat{\beta}_{\text{adjusted}} = +0.17$ pp ($134.2\%$ attenuation of the negative sign). Thus, the slight descriptive polite-tone dip in local code is entirely accounted for by polite prompts producing longer ($+12.40\%$), slower ($+10.76\%$), and less compliant ($-7.05$ pp) outputs. For local code role framing, the total penalty of $\hat{\beta}_{\text{total}} = -1.04$ pp ($p = 0.025$ clustered) is attenuated by $16.0\%$ to $\hat{\beta}_{\text{adjusted}} = -0.87$ pp after joint behavioral adjustment ($-0.83$ pp with tokens alone, $-1.06$ pp with compliance alone, $-0.83$ pp with duration alone). Aggressive tone ($-0.00$ pp total vs.\ $-0.01$ pp adjusted), JSON ($+0.17$ pp total vs.\ $+0.06$ pp adjusted), and XML ($+0.83$ pp total vs.\ $+0.82$ pp adjusted) remain near zero before and after adjustment.
\item \textbf{Local math behavioral adjustment ($n = 5,400$):} In local math, the total JSON wrapper penalty is $\hat{\beta}_{\text{total}} = -4.17$ pp ($p = 7.6 \times 10^{-5}$). Controlling for log tokens alone attenuates the JSON penalty by $25.3\%$ to $-3.11$ pp; controlling for duration alone attenuates it by $14.4\%$ to $-3.57$ pp; controlling for compliance alone leaves it at $-4.13$ pp ($0.9\%$ attenuation); and controlling jointly for all three behavioral covariates attenuates the penalty by $32.0\%$ to $\hat{\beta}_{\text{adjusted}} = -2.84$ pp. Thus, roughly one-third of the local math JSON accuracy drop co-varies with output token compaction ($-19.80\%$) and shorter generation time ($-11.16\%$), while a substantial direct negative association ($-2.84$ pp) remains. Similarly, the local math role-framing gain ($\hat{\beta}_{\text{total}} = +2.19$ pp, $p = 0.011$ unclustered, $p = 0.002$ clustered) is attenuated by $36.7\%$ to $\hat{\beta}_{\text{adjusted}} = +1.38$ pp after controlling for its increased token length ($+8.86\%$) and duration ($+3.98\%$). The local math XML penalty changes by only $8.4\%$ from $-2.22$ pp ($p = 0.035$) to $-2.04$ pp after joint adjustment, and polite tone shifts by $9.3\%$ from $+1.67$ pp to $+1.51$ pp.
\item \textbf{API Python behavioral adjustment ($n = 1{,}800$):} Controlling for log tokens and duration alone does not attenuate the two supported wrapper contrasts: JSON moves from $+3.17$ pp ($p = 0.101$) to $+3.94$ pp ($p = 0.036$, $-24.3\%$ attenuation) and XML moves from $+4.83$ pp ($p = 0.012$) to $+5.43$ pp ($p = 0.004$, $-12.4\%$ attenuation). When audited code-only compliance (\texttt{strict\_extract\_code}, $93.39\%$ overall) is also added in the full three-control specification, JSON attenuates by $52.0\%$ to $+1.52$ pp ($p = 0.370$; $+1.05$ pp with compliance alone) and XML attenuates by $27.0\%$ to $+3.53$ pp ($p = 0.037$; $+3.16$ pp with compliance alone), while tone and role contrasts remain null in all specifications ($p \ge 0.127$). This confirms that structured wrappers aid API Python both by reducing parser/compliance failures on \texttt{deepseek-v4-flash} (accounting for roughly half of the JSON gain and one-quarter of the XML gain) and through a direct $+3.53$ pp XML accuracy advantage that persists even after controlling for compliance, tokens, and latency.
\item \textbf{API math behavioral adjustment ($n = 1{,}800$; joint controls: tokens, visible-channel compliance, and duration):} The role-framing penalty is essentially unchanged by behavioral adjustment, moving from $-2.78$ pp ($p = 0.073$) to $-2.59$ pp ($p = 0.047$), an attenuation of $6.6\%$. The JSON contrast is null before and after adjustment ($-1.33 \to -1.37$ pp, $p = 0.393$), and the two tone contrasts remain null as well ($+0.33 \to -1.42$ pp, $p = 0.374$ for aggressive; $-1.00 \to -1.39$ pp, $p = 0.384$ for polite). Only the XML contrast changes materially, shrinking from $-2.33$ pp ($p = 0.218$) to $-0.72$ pp ($p = 0.651$), a $69.0\%$ attenuation that is consistent with its already exploratory status: once token compaction ($-17.52\%$), visible conciseness ($+5.13$ pp), and duration ($-9.76\%$) are held constant, little of the negative XML contrast remains, and neither estimate excludes zero.
\end{itemize}

Read together, the four strata show that behavioral adjustment has a consistent diagnostic role across both deployment tiers. In local code, controlling for observable output behavior explains the minor polite-tone dip; in local math, token compaction accounts for roughly one-third ($32.0\%$) of the JSON penalty while leaving a $-2.84$ pp direct format penalty; in API Python, audited code-only compliance accounts for $52.0\%$ of the JSON gain and $27.0\%$ of the XML gain while leaving a statistically supported $+3.53$ pp direct XML advantage ($p = 0.037$); and in API math, the $-2.59$ pp role penalty survives adjustment ($6.6\%$ attenuation). These API estimates rest on two commercial providers per stratum and are reported as descriptive complements to the six-model local decomposition.

Figure~\ref{fig:mediation-controls} displays the total and behavior-adjusted correctness contrasts across all five prompt factors for local code and local math, illustrating the complete attenuation of the polite-tone code contrast and the partial ($32.0\%$) attenuation of the local-math JSON penalty.

\begin{center}
\begin{minipage}{0.90\linewidth}
\centering
\includegraphics[width=0.85\linewidth]{mediation_controls.png}
\captionof{figure}{Prompt-factor correctness effects before and after adjustment for output length (tokens), result-only/code-only compliance, and duration. Panels show the two local strata, where a six-model cross-section makes the adjustment interpretable; the corresponding cloud API strata are reported in Section~\ref{mediation-and-behavioral-control-results}.}
\label{fig:mediation-controls}
\end{minipage}
\end{center}

Overall, these results answer RQ4 by showing that observable output behavior accounts for distinct shares of prompt sensitivity across domains: in local code, minor tone and role shifts align primarily with verbosity and parser-compliance friction, whereas in local math, response compaction accounts for roughly one-third of the JSON accuracy penalty while leaving a significant direct format penalty. Extending the same decomposition to the cloud API strata shows that audited compliance accounts for roughly half ($52.0\%$) of the API Python JSON gain and one-quarter ($27.0\%$) of the XML gain while leaving a significant $+3.53$ pp direct XML advantage ($p = 0.037$), and that the API math role penalty likewise survives adjustment, whereas the exploratory API math XML contrast does not (practical guidelines consolidated in Section~\ref{practical-prompt-policy-recommendation}, Table~\ref{tab:actionable-guidelines}).


\subsection{Cross-Cutting Robustness Check: Clustered Linear-Probability Models and Two-Way Interaction Analysis}\label{clustered-linear-probability-robustness-check}

As a frequentist robustness check on the variational-Bayes mixed-effects logistic results in Sections~\ref{prompt-factor-effects-on-code-tasks}--\ref{mixed-effects-inferential-results}, we fit pass/fail linear probability models separately by model group and task type with prompt-factor and model fixed effects and Liang--Zeger standard errors clustered by base task ($G = 50$ task clusters). The clustered LPM addresses the same repeated-task dependence as the mixed model without parametric random-effect distributional assumptions.

Across all four strata, the clustered LPM estimates closely match the mixed-effects Average Marginal Effects in both magnitude and statistical significance: in API Python, XML is $+4.83$ pp (clustered $\text{SE} = 1.77$, $p = 0.006$ vs.\ mixed-logit $\text{AME} = +5.2$ pp, $p < 0.001$) and JSON is $+3.17$ pp (clustered $\text{SE} = 1.73$, $p = 0.068$ vs.\ $\text{AME} = +3.4$ pp, $p = 0.022$); in API math, role framing is $-2.78$ pp (clustered $\text{SE} = 1.33$, $p = 0.037$ vs.\ $\text{AME} = -2.8$ pp, $p = 0.006$); in local code, role framing is $-1.04$ pp (clustered $\text{SE} = 0.46$, $p = 0.025$ vs.\ $\text{AME} = -1.0$ pp, $p = 0.013$); and in local math, JSON is $-4.17$ pp (clustered $\text{SE} = 1.11$, $p < 0.001$ vs.\ $\text{AME} = -4.1$ pp, $p < 0.001$), XML is $-2.22$ pp (clustered $\text{SE} = 1.02$, $p = 0.030$ vs.\ $\text{AME} = -2.2$ pp, $p = 0.001$), role framing is $+2.19$ pp (clustered $\text{SE} = 0.69$, $p = 0.002$ vs.\ $\text{AME} = +2.1$ pp, $p < 0.001$), and polite tone is $+1.67$ pp (clustered $\text{SE} = 0.79$, $p = 0.035$ vs.\ $\text{AME} = +1.6$ pp, $p = 0.019$). Thus, every headline finding is invariant to the choice between Bayesian mixed-logit and frequentist clustered-LPM estimation.

\paragraph{Two-Way Prompt-Factor Interaction Results ($\text{Tone} \times \text{Role}$, $\text{Tone} \times \text{Format}$, and $\text{Role} \times \text{Format}$).}
To test whether the effect of one prompt factor depends on the level of another within the $3 \times 2 \times 3$ design, Table~\ref{tab:two-way-interactions} reports all eight two-way interaction terms across each of the four strata, estimated with task-clustered standard errors ($G = 50$), $95\%$ confidence intervals, omnibus joint Wald $\chi^2(8)$ tests, and Benjamini--Hochberg FDR-adjusted $q$-values within each 8-test stratum ($q_{\text{stratum}}$) and across all 32 tests ($q_{\text{all32}}$).

\begin{table}[htbp]
\centering
\caption{Two-way prompt-factor interaction estimates ($\text{Tone} \times \text{Role}$, $\text{Tone} \times \text{Format}$, and $\text{Role} \times \text{Format}$) on Pass@1 (percentage points) from task-clustered linear probability models ($G = 50$ base-task clusters per stratum), with $95\%$ confidence intervals, raw $p$-values, within-stratum BH $q_{\text{stratum}}$, and global 32-test BH $q_{\text{all32}}$.}\label{tab:two-way-interactions}
\footnotesize
\setlength{\tabcolsep}{4.5pt}
\begin{tabular}{@{}llrrrrrr@{}}
\toprule
\textbf{Stratum \& Joint Test} & \textbf{Two-Way Interaction Term} & \textbf{Coef.\ (pp)} & \textbf{Cluster SE} & \textbf{95\% CI (pp)} & \textbf{Raw $p$} & \textbf{$q_{\text{stratum}}$} & \textbf{$q_{\text{all32}}$} \\
\midrule
\textbf{Local code ($n = 5,400$)} & $\text{Aggressive} \times \text{Role}$ & $-0.44$ & $1.36$ & $[-3.10,\; +2.22]$ & $0.743$ & $0.880$ & $0.942$ \\
$\chi^2(8) = 14.79$, $F = 1.85$, & $\text{Polite} \times \text{Role}$ & $+1.00$ & $1.44$ & $[-1.82,\; +3.82]$ & $0.487$ & $0.880$ & $0.888$ \\
$p_{\text{joint}} = 0.063$ & $\text{Aggressive} \times \text{JSON}$ & $+0.33$ & $1.35$ & $[-2.31,\; +2.98]$ & $0.805$ & $0.880$ & $0.942$ \\
 & $\text{Aggressive} \times \text{XML}$ & $+1.67$ & $1.04$ & $[-0.37,\; +3.70]$ & $0.109$ & $0.290$ & $0.616$ \\
 & $\text{Polite} \times \text{JSON}$ & $+4.17$ & $1.57$ & $[+1.09,\; +7.24]$ & $0.008$ & $0.048^*$ & $0.193$ \\
 & $\text{Polite} \times \text{XML}$ & $+3.33$ & $1.33$ & $[+0.73,\; +5.94]$ & $0.012$ & $0.048^*$ & $0.193$ \\
 & $\text{Role} \times \text{JSON}$ & $-0.56$ & $1.17$ & $[-2.85,\; +1.74]$ & $0.635$ & $0.880$ & $0.942$ \\
 & $\text{Role} \times \text{XML}$ & $+0.11$ & $0.74$ & $[-1.33,\; +1.56]$ & $0.880$ & $0.880$ & $0.942$ \\
\midrule
\textbf{Local math ($n = 5,400$)} & $\text{Aggressive} \times \text{Role}$ & $-0.11$ & $1.54$ & $[-3.12,\; +2.90]$ & $0.942$ & $0.942$ & $0.973$ \\
$\chi^2(8) = 22.16$, $F = 2.77$, & $\text{Polite} \times \text{Role}$ & $-2.00$ & $1.20$ & $[-4.35,\; +0.35]$ & $0.095$ & $0.444$ & $0.616$ \\
$p_{\text{joint}} = 0.005$ & $\text{Aggressive} \times \text{JSON}$ & $+1.83$ & $1.51$ & $[-1.13,\; +4.80]$ & $0.226$ & $0.444$ & $0.888$ \\
 & $\text{Aggressive} \times \text{XML}$ & $+1.00$ & $1.42$ & $[-1.79,\; +3.79]$ & $0.483$ & $0.551$ & $0.888$ \\
 & $\text{Polite} \times \text{JSON}$ & $-1.83$ & $1.67$ & $[-5.12,\; +1.45]$ & $0.274$ & $0.444$ & $0.888$ \\
 & $\text{Polite} \times \text{XML}$ & $-2.67$ & $1.69$ & $[-5.99,\; +0.65]$ & $0.115$ & $0.444$ & $0.616$ \\
 & $\text{Role} \times \text{JSON}$ & $+1.67$ & $1.58$ & $[-1.42,\; +4.76]$ & $0.290$ & $0.444$ & $0.888$ \\
 & $\text{Role} \times \text{XML}$ & $+1.56$ & $1.61$ & $[-1.59,\; +4.70]$ & $0.333$ & $0.444$ & $0.888$ \\
\midrule
\textbf{API Python ($n = 1,800$)} & $\text{Aggressive} \times \text{Role}$ & $+3.67$ & $4.15$ & $[-4.46,\; +11.80]$ & $0.377$ & $1.000$ & $0.888$ \\
$\chi^2(8) = 3.62$, $F = 0.45$, & $\text{Polite} \times \text{Role}$ & $+0.33$ & $3.04$ & $[-5.63,\; +6.30]$ & $0.913$ & $1.000$ & $0.973$ \\
$p_{\text{joint}} = 0.890$ & $\text{Aggressive} \times \text{JSON}$ & $-2.50$ & $4.38$ & $[-11.08,\; +6.08]$ & $0.568$ & $1.000$ & $0.942$ \\
 & $\text{Aggressive} \times \text{XML}$ & $+1.00$ & $5.12$ & $[-9.03,\; +11.03]$ & $0.845$ & $1.000$ & $0.942$ \\
 & $\text{Polite} \times \text{JSON}$ & $-0.00$ & $5.27$ & $[-10.33,\; +10.33]$ & $1.000$ & $1.000$ & $1.000$ \\
 & $\text{Polite} \times \text{XML}$ & $+3.00$ & $4.22$ & $[-5.27,\; +11.27]$ & $0.477$ & $1.000$ & $0.888$ \\
 & $\text{Role} \times \text{JSON}$ & $+1.00$ & $2.96$ & $[-4.81,\; +6.81]$ & $0.736$ & $1.000$ & $0.942$ \\
 & $\text{Role} \times \text{XML}$ & $+1.00$ & $3.92$ & $[-6.69,\; +8.69]$ & $0.799$ & $1.000$ & $0.942$ \\
\midrule
\textbf{API math ($n = 1,800$)} & $\text{Aggressive} \times \text{Role}$ & $-2.00$ & $2.81$ & $[-7.51,\; +3.51]$ & $0.477$ & $0.764$ & $0.888$ \\
$\chi^2(8) = 7.63$, $F = 0.95$, & $\text{Polite} \times \text{Role}$ & $+2.67$ & $2.96$ & $[-3.14,\; +8.47]$ & $0.368$ & $0.764$ & $0.888$ \\
$p_{\text{joint}} = 0.471$ & $\text{Aggressive} \times \text{JSON}$ & $+3.00$ & $4.16$ & $[-5.15,\; +11.15]$ & $0.471$ & $0.764$ & $0.888$ \\
 & $\text{Aggressive} \times \text{XML}$ & $-0.50$ & $3.55$ & $[-7.46,\; +6.46]$ & $0.888$ & $0.888$ & $0.942$ \\
 & $\text{Polite} \times \text{JSON}$ & $+8.00$ & $3.89$ & $[+0.37,\; +15.63]$ & $0.040$ & $0.318$ & $0.424$ \\
 & $\text{Polite} \times \text{XML}$ & $+1.00$ & $3.79$ & $[-6.43,\; +8.43]$ & $0.792$ & $0.888$ & $0.942$ \\
 & $\text{Role} \times \text{JSON}$ & $-2.67$ & $2.88$ & $[-8.32,\; +2.98]$ & $0.355$ & $0.764$ & $0.888$ \\
 & $\text{Role} \times \text{XML}$ & $-0.67$ & $3.61$ & $[-7.74,\; +6.40]$ & $0.853$ & $0.888$ & $0.942$ \\
\bottomrule
\end{tabular}
\end{table}

Three findings emerge from Table~\ref{tab:two-way-interactions}. First, across all 32 two-way factor-by-factor interaction parameters, \textbf{29 have $95\%$ confidence intervals that include zero, and zero of the 32 interactions survive global Benjamini--Hochberg correction across the 32-test family} ($\min q_{\text{all32}} = 0.193$). Moreover, the omnibus joint Wald test of all eight interaction terms is completely null in API Python ($\chi^2(8) = 3.62$, $p = 0.890$) and API math ($\chi^2(8) = 7.63$, $p = 0.471$), and all eight individual interaction terms in local math have unadjusted $p \ge 0.095$ ($q_{\text{stratum}} \ge 0.444$). In particular, all eight $\text{Role} \times \text{Format}$ terms across the four strata are small and statistically null ($p \ge 0.290$), indicating that the effect of persona role framing does not depend on whether the prompt is serialized in plaintext, JSON, or XML.

Second, within the local code stratum ($\chi^2(8) = 14.79$, $p_{\text{joint}} = 0.063$), two theoretically interpretable $\text{Tone} \times \text{Format}$ terms reach significance under within-stratum BH correction ($q_{\text{stratum}} = 0.048$): $\text{Polite} \times \text{JSON}$ ($+4.17$ pp, $\text{SE} = 1.57$, $95\%$ CI $[+1.09, +7.24]$, $p = 0.008$) and $\text{Polite} \times \text{XML}$ ($+3.33$ pp, $\text{SE} = 1.33$, $95\%$ CI $[+0.73, +5.94]$, $p = 0.012$). Inspecting the underlying simple-effect cell means explains why: under unstructured plaintext, polite wording reduces local code Pass@1 by $-3.00$ pp from $10.00\%$ (plaintext $\times$ neutral, $n = 600$) to $7.00\%$ (plaintext $\times$ polite, $n = 600$) because polite conversational preambles interfere with code extraction; however, when the prompt is enclosed in a JSON or XML wrapper, the structural delimiters suppress this polite-wording accuracy dip ($9.83\%$ vs.\ $8.67\%$ under JSON; $9.50\%$ vs.\ $9.17\%$ under XML). Because these two terms do not survive global 32-test FDR correction ($q_{\text{all32}} = 0.193$), we treat them as suggestive evidence that structured wrappers buffer against polite-tone verbosity rather than as standalone confirmatory rules.

\paragraph{Factor-by-Model Interaction Tests ($\text{Factor} \times \text{Model}$ Heterogeneity for RQ2).}
To formally test RQ2's second clause—whether prompt-factor effects are homogeneous or heterogeneous across models within each stratum—we also compute task-clustered omnibus Wald $F$-tests for $\text{Factor} \times \text{Model}$ interactions:
\begin{itemize}
\tightlist
\item \textbf{Local math ($n = 5,400$, 6 models):} The overall $\text{Factor} \times \text{Model}$ interaction is significant ($\chi^2(25) = 91.68$, $F = 3.67$, $p = 2.5 \times 10^{-9}$), driven specifically by $\text{Format} \times \text{Model}$ heterogeneity ($F(10) = 4.07$, $p = 1.6 \times 10^{-5}$) whose practical magnitude is large and model-specific: \texttt{llama3.1-8b} loses $-15.7$ pp under JSON and $-9.7$ pp under XML, and \texttt{deepseek-coder-6.7b} loses $-5.7$ pp under XML, whereas \texttt{qwen2.5-7b} gains $+10.2$ pp from role framing. By contrast, $\text{Tone} \times \text{Model}$ is marginal ($F(10) = 1.79$, $p = 0.057$) and $\text{Role} \times \text{Model}$ is overall non-significant ($F(5) = 1.61$, $p = 0.152$), reflecting that single $+10.2$ pp outlier alongside near-zero shifts on the other five models.
\item \textbf{API math ($n = 1,800$, 2 models):} The overall $\text{Factor} \times \text{Model}$ test is non-significant ($\chi^2(5) = 9.03$, $F = 1.81$, $p = 0.107$), with zero heterogeneity for $\text{Role} \times \text{Model}$ ($F(1) = 0.06$, $p = 0.797$) and $\text{Tone} \times \text{Model}$ ($F(2) = 0.07$, $p = 0.920$)—confirming that the $-2.8$ pp API math role-framing penalty is consistent across both cloud models—while $\text{Format} \times \text{Model}$ shows moderate divergence ($F(2) = 3.41$, $p = 0.032$).
\item \textbf{API Python ($n = 1,800$, 2 models) and Local code ($n = 5,400$, 6 models):} In API Python, neither the omnibus $\text{Factor} \times \text{Model}$ test ($\chi^2(5) = 5.80$, $F = 1.16$, $p = 0.326$) nor any sub-family ($\text{Tone} \times \text{Model}$ $p = 0.119$; $\text{Role} \times \text{Model}$ $p = 0.489$; $\text{Format} \times \text{Model}$ $p = 0.295$) rejects homogeneity. Similarly, in local code, none of the three individual factor-by-model sub-families is statistically significant ($\text{Tone} \times \text{Model}$ $F(10) = 0.86$, $p = 0.575$; $\text{Role} \times \text{Model}$ $F(5) = 0.93$, $p = 0.461$; $\text{Format} \times \text{Model}$ $F(10) = 0.91$, $p = 0.521$).
\end{itemize}
These formal interaction tests directly substantiate our RQ2 conclusion: within code generation, prompt-factor effects are broadly consistent in direction and magnitude across models, whereas in mathematical reasoning, structured-wrapper sensitivity varies substantially by model architecture—$\text{Format} \times \text{Model}$ is supported locally ($p = 1.6 \times 10^{-5}$) and in API math ($p = 0.032$), with model-specific wrapper penalties as large as $-15.7$ pp (Llama 3.1 JSON).

\subsection{Cross-Cutting Sensitivity: Dominant-Model Leave-One-Out Analysis}\label{dominant-model-sensitivity-and-aggregation-risk}

A single high-performing model can substantially influence pooled averages. In the local setting, removing Nemotron reduces the pooled code pass rate from 9.1\% to 6.4\% and the pooled math pass rate from 31.1\% to 19.1\%. In the API setting, removing DeepSeek V4 Flash has little effect on Python code generation, changing the pass rate from 67.8\% to 69.0\%, but it reduces the API math pass rate from 53.3\% to 24.1\%. These shifts show why pooled averages are reported together with model-stratified and dominant-model sensitivity results.

Pass-rate averages alone do not show whether the prompt effects themselves are driven by the dominant model. We therefore re-estimated the main tone, role, and format contrasts with and without the dominant model in each group, using both the clustered linear probability model and the mixed-effects logistic model. Table~\ref{tab:dominant-effects} reports the percentage-point effects in both panels. For the API group, the leave-one-out panel contains only Gemini Flash, so those columns should be read as single-model estimates.

\begin{longtable}[]{@{}llrrrr@{}}
\caption{Prompt-factor effects with and without the dominant model: clustered-LPM percentage-point contrasts and mixed-effects average marginal effects (AME pp).}\label{tab:dominant-effects}\tabularnewline
\toprule\noalign{}
Stratum & Contrast & LPM full & LPM w/o dom. & AME full & AME w/o dom. \\
\midrule\noalign{}
\endfirsthead
\toprule\noalign{}
Stratum & Contrast & LPM full & LPM w/o dom. & AME full & AME w/o dom. \\
\midrule\noalign{}
\endhead
\bottomrule\noalign{}
\endlastfoot
API Python & JSON wrapper & +3.2 & +1.3 & +3.4 & +1.6 \\
API Python & XML wrapper & +4.8 & +2.7 & +5.2 & +3.1 \\
API math & Role framing & -2.8 & -2.4 & -2.8 & -4.1 \\
API math & JSON wrapper & -1.3 & -3.7 & -1.2 & -6.1 \\
API math & XML wrapper & -2.3 & -0.0 & -2.3 & -0.2 \\
Local code & Role framing & -1.0 & -0.9 & -1.0 & -0.9 \\
Local math & Polite tone & +1.7 & +2.3 & +1.6 & +2.3 \\
Local math & Role framing & +2.2 & +2.4 & +2.1 & +2.3 \\
Local math & JSON wrapper & -4.2 & -5.3 & -4.1 & -5.2 \\
Local math & XML wrapper & -2.2 & -2.5 & -2.2 & -2.4 \\
\end{longtable}

The re-estimation shows that the local prompt effects are not simply artifacts of Nemotron's high pass rate. Without Nemotron, the local math JSON penalty becomes slightly larger, changing from -4.1 to -5.2 percentage points in the mixed-effects AME estimate. The local math gains for polite tone and role framing remain similar in direction and magnitude, and the local-code role penalty also remains stable at $-1.0$ percentage points. Thus, Nemotron changes the pooled pass rate substantially, but it does not explain away the local prompt-effect estimates.

The API group shows a different pattern, which is expected in a two-model panel. The API Python wrapper effects shrink after removing DeepSeek V4 Flash: the XML AME decreases from +5.2 to +3.1 percentage points, and the JSON AME decreases from +3.4 to +1.6 percentage points. These effects should therefore be interpreted as two-model averages with a larger contribution from DeepSeek V4 Flash. In contrast, the API math role-framing penalty remains negative and becomes larger in the Gemini-only panel, changing from -2.8 to -4.1 percentage points. The API math XML effect does not remain after removing DeepSeek V4 Flash, which supports treating it as exploratory.

This sensitivity analysis supports the main methodological point of the paper: prompt-policy conclusions should not be based on pooled averages alone. They need to be interpreted by model group, task type, and dominant-model dependence.

Given Nemotron's high local math pass rate of 91.2\%, we also audited the corresponding repository file. The file contains 900 distinct runs, covering 50 base tasks across 18 prompt variants, with no duplicate task IDs. Recomputing pass/fail scores from \texttt{expected\_answer} and \texttt{extracted\_values} reproduced 821 out of 900 passes with zero mismatches against the stored \texttt{passed} field, and a scan found no expected-answer leakage. At the task level, 42 of the 50 tasks were solved under all 18 prompt conditions, 6 were solved under some conditions, and 2 were not solved under any condition. This supports treating the 91.2\% value as a reproducible result under the stated evaluator, while still motivating the no-Nemotron sensitivity check.

\subsection{Qualitative Diagnostics: Extraction and Execution Failure Modes}\label{error-and-failure-modes}

The experiment shows some of the possible failure modes. Regarding code-based tasks, a few of these failures include wrong algorithms, mistakes in syntax, missing functions, or even outputs incorporating prose or markdown that the extractor cannot cleanly process. The use of both polite and role prompts may lead to parsing issues since they can lead to more explanatory text. Formats such as JSON and XML may make outputs more compact, but they do not guarantee correctness.

The types of failures in math tasks include wrong answers, missing final-answer markers, too brief or vague answers, and extraction ambiguity. Using JSON reduces the duration and length of local math responses while decreasing the pass rate, suggesting that using compact formatting can lower accuracy in certain reasoning tasks even if the requested output is final-result-only. This does not mean that showing reasoning is always required; rather, the results show that less verbose responses are associated with lower local math accuracy in this study.

Representative cases demonstrate the contrast between task-related failures and parser or interface failures. In one Nemotron code run, the response returned by the model was empty, leading to the evaluation of \texttt{No code extracted}. In a different incident with Qwen 2.5, Python code was executed; however, the model used helper names that did not satisfy the requirements of the callable interface, due to which the code execution failed with name-resolution errors. One more instance occurred when a local math run was conducted in Qwen 2.5 using the JSON-format prompt. The result was that the answer \texttt{6} was given when the expected answer was \texttt{12}, showing the compact nature of the answer, which was easily extractable yet erroneous. These examples support the broader interpretation that correctness, extraction compatibility, and output behavior should be reported separately.

\subsection{End-to-End Evaluation: Simulated Offline Replay Versus Live Pipeline With Real Retries}\label{end-to-end-replay}

The harness measures model outputs, whereas production systems consume them through parsers, validators, and retry logic. We evaluate this connection in two clearly separated ways: (i) an \emph{offline replay} of all $14{,}400$ stored responses (where single-attempt rates are measured and retry gains are \emph{simulated} under an independent-redraw upper bound), and (ii) a \emph{live end-to-end Python pipeline} ($450$ workflows across three models) executing \emph{real} error-feedback retry calls.

\paragraph{Offline Replay (Simulated Retry Upper Bound, Table~\ref{tab:pipeline-replay}).}
Replaying all $14{,}400$ stored responses through a concrete downstream consumer---strict extraction (fenced code block or bare code body for Python; numeric parse for math), syntactic validation by \texttt{ast.parse}, and held-out test execution---shows that parser-level failures are a minority of code failures: strict parsing fails on $2.4\%$ of local-code runs ($2.7\%$ of local-code failures) and on $6.6\%$ of API Python runs ($20.4\%$ of API failures; $8.5\%$ under plaintext vs.\ $5.3\%$ JSON and $5.8\%$ XML). Under an idealized independent-redraw simulation ($q$ equal to each model--task's neutral-plaintext pass rate), two retries would lift API Python from $67.8\%$ to $96.7\%$ ($1.43\times$ calls) and local code from $9.1\%$ to $25.8\%$ ($2.73\times$ calls). However, this simulation is strictly an \emph{upper bound} because failures on the same task are correlated rather than independent.

\begin{table}[htbp]
\centering
\caption{End-to-end replay of all $14{,}400$ stored responses through a strict downstream parser, validator, and reformat-retry policy: measured single-attempt success, parser-level failure exposure, and simulated success under one and two retries with the associated call and latency cost.}\label{tab:pipeline-replay}
\footnotesize
\setlength{\tabcolsep}{3pt}
\begin{tabular}{@{}lcccccccc@{}}
\toprule
\textbf{Stratum} & \textbf{Pass@1} & \textbf{Strict-parse} & \textbf{Parser share} & \textbf{Success} & \textbf{Success} & \textbf{Success} & \textbf{Calls per} & \textbf{Added} \\
 & \textbf{(1 attempt)} & \textbf{failures} & \textbf{of failures} & \textbf{R = 0} & \textbf{R = 1} & \textbf{R = 2} & \textbf{1,000 (R = 2)} & \textbf{latency (R = 2)} \\
\midrule
Local code ($n = 5{,}400$) & $9.1\%$ & $2.4\%$ & $2.7\%$ & $9.1\%$ & $17.9\%$ & $25.8\%$ & $2{,}730$ & $+173\%$ \\
Local math ($n = 5{,}400$) & $31.1\%$ & $0.4\%$ & $0.6\%$ & $31.1\%$ & $52.5\%$ & $67.2\%$ & $2{,}165$ & $+116\%$ \\
API Python ($n = 1{,}800$) & $67.8\%$ & $6.6\%$ & $20.4\%$ & $67.8\%$ & $89.7\%$ & $96.7\%$ & $1{,}425$ & $+42\%$ \\
API math ($n = 1{,}800$) & $53.3\%$ & $0.0\%$ & $0.0\%$ & $53.3\%$ & $80.8\%$ & $92.1\%$ & $1{,}659$ & $+66\%$ \\
\bottomrule
\end{tabular}

\vspace{2pt}
\begin{minipage}{\linewidth}
\footnotesize
\sloppy
\textit{Notes:} Pass@1 is the recorded harness pass rate. Strict-parse failures are responses from which the replay consumer could not extract a syntactically valid artifact; the parser share divides those by all failures. $R$ denotes the number of reformat retries; $R = 0$ is measured from the stored data, while $R = 1$ and $R = 2$ apply $1 - (1 - p_0)(1 - q)^R$ with $q$ equal to the model--task neutral-plaintext pass rate, i.e.\ the independent-redraw assumption. Calls per 1,000 and added latency are expected values for the $R = 2$ policy, with latency computed from each response's recorded duration. For math tasks, non-reasoning models store the full raw text in \texttt{response\_text} while reasoning models (\texttt{nemotron-3-nano-4b}, \texttt{deepseek-v4-flash}) emit hidden thinking tokens in a separate channel and return concise visible answers, so math strict-parse failures reflect unparseable visible numeric expressions. Generated by \texttt{run\_pipeline\_replay.py} in the replication package.
\end{minipage}
\end{table}

\paragraph{Live End-to-End Pipeline (Measured Real Retries, Table~\ref{tab:live-pipeline-measured}).}
To measure actual retry recovery, tokens, and wall-clock latency without simulation assumptions, we executed a live 4-stage Python pipeline (model generation $\to$ strict parser $\to$ \texttt{ast.parse} $\to$ isolated 7-assertion subprocess execution $\to$ up to two \emph{real} follow-up LLM calls fed the exact failure stage and error message) across all $50$ Python tasks $\times$ $3$ representative policies (Neutral+Plaintext, Neutral+XML, Polite+Plaintext) for \texttt{gemini-flash} (API, $T = 0$), \texttt{qwen2.5:7b}, and \texttt{llama3.1:8b} (local Ollama, $T = 0.2$; $450$ live workflows in total; Attempt~0 reused from the $k = 3$ runs where available). Table~\ref{tab:live-pipeline-measured} reports the \emph{measured} outcomes alongside the clearly separated simulated upper bounds in Table~\ref{tab:pipeline-replay}. Three empirical findings emerge:
\begin{itemize}
\tightlist
\item \textbf{Real $T = 0$ API retries recover far less than the simulated upper bound:} On \texttt{gemini-flash}, measured success rises from $68.0\%$--$70.0\%$ ($R = 0$) to $74.0\%$--$76.0\%$ ($R = 1$) and remains flat at $74.0\%$--$76.0\%$ ($R = 2$)---a measured gain of only $+4.0$ to $+6.0$ pp (at $1{,}540$--$1{,}580$ calls per $1{,}000$ tasks), compared with the simulated $+28.9$ pp upper bound ($96.7\%$). Persistent failures are predominantly algorithmic blind spots rather than independent sampling draws.
\item \textbf{Local $T = 0.2$ retries yield larger absolute recovery but at steep workflow cost:} On \texttt{llama3.1:8b} and \texttt{qwen2.5:7b}, two real retries raise success by $+8.0$ to $+14.0$ pp ($10$--$22\% \to 20$--$32\%$), requiring $2{,}480$--$2{,}740$ calls per $1{,}000$ tasks. On \texttt{llama3.1:8b}, Polite+Plaintext is strictly dominated by Neutral+Plaintext in end-to-end workflow efficiency ($20.0\%$ vs.\ $32.0\%$ $R = 2$ success, $1{,}074$ vs.\ $829$ completion tokens, $44.6$\,s vs.\ $36.7$\,s wall-clock LLM time).
\item \textbf{LLM inference dominates end-to-end latency ($>95\%$):} Downstream strict parsing, AST validation, and 7-test subprocess execution add only $0.07$--$0.30$\,s per workflow on API tasks and $0.46$--$1.50$\,s locally, compared with $26.9$--$44.6$\,s of cumulative LLM inference time.
\end{itemize}

\begin{table}[htbp]
\centering
\caption{Measured live end-to-end Python pipeline with up to two \emph{real} error-feedback retries ($n = 450$ workflows: $50$ tasks $\times$ $3$ policies $\times$ $3$ models). Clearly separated from the simulated offline replay upper bounds in Table~\ref{tab:pipeline-replay}.}\label{tab:live-pipeline-measured}
\footnotesize
\setlength{\tabcolsep}{3.5pt}
\begin{tabular}{@{}llccccccccc@{}}
\toprule
\textbf{Model (Regime)} & \textbf{Policy} & \textbf{$R=0$} & \textbf{$R=1$} & \textbf{$R=2$} & \textbf{Gain} & \textbf{Calls / 1k} & \textbf{Total Tok} & \textbf{LLM (s)} & \textbf{Consumer (ms)} \\
\midrule
\texttt{gemini-flash} (API) & Neutral+Plaintext & $68.0\%$ & $74.0\%$ & $74.0\%$ & $+6.0$ pp & $1{,}580$ & $552.3$ & $36.93$ & $302.8$ \\
\texttt{gemini-flash} (API) & Neutral+XML & $70.0\%$ & $76.0\%$ & $76.0\%$ & $+6.0$ pp & $1{,}540$ & $564.0$ & $30.39$ & $285.0$ \\
\texttt{gemini-flash} (API) & Polite+Plaintext & $70.0\%$ & $74.0\%$ & $74.0\%$ & $+4.0$ pp & $1{,}560$ & $571.2$ & $34.44$ & $65.6$ \\
\midrule
\texttt{llama3.1:8b} (Local) & Neutral+Plaintext & $18.0\%$ & $28.0\%$ & $32.0\%$ & $+14.0$ pp & $2{,}540$ & $828.8$ & $36.71$ & $1{,}257.9$ \\
\texttt{llama3.1:8b} (Local) & Neutral+XML & $10.0\%$ & $18.0\%$ & $22.0\%$ & $+12.0$ pp & $2{,}720$ & $866.4$ & $36.47$ & $694.7$ \\
\texttt{llama3.1:8b} (Local) & Polite+Plaintext & $12.0\%$ & $14.0\%$ & $20.0\%$ & $+8.0$ pp & $2{,}740$ & $1{,}074.3$ & $44.57$ & $821.8$ \\
\midrule
\texttt{qwen2.5:7b} (Local) & Neutral+Plaintext & $22.0\%$ & $30.0\%$ & $32.0\%$ & $+10.0$ pp & $2{,}480$ & $611.2$ & $26.86$ & $1{,}501.3$ \\
\texttt{qwen2.5:7b} (Local) & Neutral+XML & $18.0\%$ & $30.0\%$ & $30.0\%$ & $+12.0$ pp & $2{,}520$ & $705.4$ & $28.13$ & $461.0$ \\
\texttt{qwen2.5:7b} (Local) & Polite+Plaintext & $22.0\%$ & $30.0\%$ & $30.0\%$ & $+8.0$ pp & $2{,}480$ & $675.8$ & $26.99$ & $1{,}243.0$ \\
\bottomrule
\end{tabular}
\end{table}

\subsection{Empirical Validation: Repeated Generations ($k = 3$), Neutral Controls, Shared-Task Check, and Non-Math SE Benchmark}\label{empirical-validation-suite}

To directly address run-to-run nondeterminism, prompt-length confounding, cross-regime dataset confounding, and external validity beyond math-derived tasks, we executed four targeted follow-up experiments using full raw-response logging.

\paragraph{1. Independent Repeated Generations ($k = 3$ across All 8 Models, $14{,}398$ Runs).}
We repeated all $50$ base tasks across the $6$ core one-factor-at-a-time conditions (Neutral+Plaintext baseline plus JSON, XML, Role, Polite, and Aggressive) for $k = 3$ independent runs per cell across all six local Ollama models ($T = 0.2$, Kaggle T4 GPU; every contrast is estimated within model and session so hardware is held fixed across conditions) and both cloud API models ($T = 0$; \texttt{gemini-3.1-flash-lite-preview} and \texttt{deepseek/deepseek-v4-flash} via AI/ML API after Pioneer AI's host retirement), yielding $14{,}398 / 14{,}400$ valid runs ($99.99\%$; Table~\ref{tab:k3-repeated-variability}). Within-cell variability is sharply model-dependent: \texttt{gemini-flash} ($T = 0$) exhibits $0.00\%$ cell-level flip rate ($100.00\%$ exact agreement across all $600$ cells), whereas \texttt{deepseek-v4-flash} ($T = 0$) exhibits $23.67\%$ flips in Python and $15.00\%$ in math, confirming that hosted APIs can be substantially nondeterministic even at temperature $0$ (moreover, switching third-party API hosts from retired Pioneer AI to AI/ML API lowered \texttt{deepseek-v4-flash}'s Python pass rate from $66.7\%$ to $39.67\%$, illustrating how provider-side routing and quantization shift hosted model performance). Local models ($T = 0.2$) show $2.33\%$--$20.00\%$ flips in code and $2.67\%$--$28.33\%$ in math. Crucially, aggregate run-to-run SD across repetitions is small ($0.00$--$2.08$ pp), and repeating $k = 3$ tightens the $80\%$-power MDE for \emph{single-cell pairwise contrasts against the neutral-plaintext baseline} ($n = 600$ at $k = 1 \to n = 1{,}800$ at $k = 3$, distinct from the 18-variant marginal main-effect MDEs of $1.29$--$3.12$ pp in Table~\ref{tab:power-analysis}) by $15$--$40\%$: from $3.45$--$5.02$ pp ($k = 1$) down to $3.05$--$3.86$ pp in local code and $3.00$--$3.94$ pp in local math. Headline local conclusions replicate cleanly: under $k = 3$, Local Python Polite vs.\ Neutral is $-2.56$ pp ($95\%$ CI $[-4.85, -0.27]$) and Local Math JSON vs.\ Plaintext is $-6.78$ pp ($95\%$ CI $[-9.53, -4.03]$), whereas all other pairwise contrasts (including API Python XML at $+1.67$ pp $[-4.40, +7.73]$ and JSON at $+0.33$ pp $[-4.56, +5.23]$ under the noisier AI/ML API host) have confidence intervals crossing zero.

\begin{table}[htbp]
\centering
\caption{Independent $k = 3$ repeated-generation variability across all 8 study models ($50$ tasks $\times$ $6$ core conditions $\times$ $3$ repeats $= 14{,}398 / 14{,}400$ valid runs).}\label{tab:k3-repeated-variability}
\footnotesize
\setlength{\tabcolsep}{4pt}
\begin{tabular}{@{}llccccc@{}}
\toprule
\textbf{Stratum} & \textbf{Model} & \textbf{Runs} & \textbf{Pass \%} & \textbf{Rep SD (pp)} & \textbf{Flip Rate \%} & \textbf{Exact Agree \%} \\
\midrule
API Python & \texttt{deepseek-v4-flash} (AI/ML API) & $900$ & $39.67$ & $2.08$ & $23.67$ & $76.33$ \\
API Python & \texttt{gemini-flash} & $900$ & $68.33$ & $0.00$ & $0.00$ & $100.00$ \\
API math & \texttt{deepseek-v4-flash} (AI/ML API) & $900$ & $80.11$ & $0.51$ & $15.00$ & $85.00$ \\
API math & \texttt{gemini-flash} & $900$ & $23.67$ & $0.00$ & $0.00$ & $100.00$ \\
\midrule
Local Python & \texttt{deepseek-coder-6.7b} & $899$ & $3.00$ & $0.58$ & $5.00$ & $95.00$ \\
Local Python & \texttt{gemma2-9b} & $900$ & $13.89$ & $0.51$ & $7.00$ & $93.00$ \\
Local Python & \texttt{llama3.1-8b} & $900$ & $8.78$ & $0.69$ & $12.67$ & $87.33$ \\
Local Python & \texttt{mistral-7b-instruct} & $900$ & $1.00$ & $0.88$ & $2.33$ & $97.67$ \\
Local Python & \texttt{nemotron-3-nano-4b} & $900$ & $17.11$ & $0.38$ & $20.00$ & $80.00$ \\
Local Python & \texttt{qwen2.5-7b} & $900$ & $18.33$ & $1.33$ & $12.00$ & $88.00$ \\
\midrule
Local math & \texttt{deepseek-coder-6.7b} & $899$ & $4.78$ & $0.52$ & $10.67$ & $89.33$ \\
Local math & \texttt{gemma2-9b} & $900$ & $17.56$ & $0.19$ & $2.67$ & $97.33$ \\
Local math & \texttt{llama3.1-8b} & $900$ & $48.22$ & $0.51$ & $28.33$ & $71.67$ \\
Local math & \texttt{mistral-7b-instruct} & $900$ & $9.33$ & $1.00$ & $10.00$ & $90.00$ \\
Local math & \texttt{nemotron-3-nano-4b} & $900$ & $90.56$ & $0.19$ & $4.33$ & $95.67$ \\
Local math & \texttt{qwen2.5-7b} & $900$ & $24.00$ & $0.33$ & $4.00$ & $96.00$ \\
\bottomrule
\end{tabular}
\end{table}

\paragraph{2. Neutral Paraphrase and Length-Matched Controls ($180$ Runs).}
Evaluating $30$ Python tasks across two semantically neutral controls---a length-matched short neutral paraphrase and a longer neutral technical prompt matched in token length to Polite/Aggressive ($\approx +58$ prompt tokens)---yielded $73.3\%$ vs.\ $66.7\%$ on \texttt{gemini-flash}, $20.0\%$ vs.\ $20.0\%$ on \texttt{llama3.1:8b}, and $26.7\%$ vs.\ $30.0\%$ on \texttt{qwen2.5:7b}. Neutral rewording and length variation alone therefore induce a $0.0$--$6.7$ pp noise band on small task subsets, reinforcing why $1$--$2$ pp tone contrasts must be interpreted conservatively.

\paragraph{3. Shared-Task Cross-Regime Comparison ($7$ Models on $20$ Identical TOANMATH Tasks).}
To de-confound model capability from task-dataset difficulty, we evaluated \texttt{gemini-flash} on $20$ local TOANMATH Python tasks ($\times 3$ policies $= 60$ runs) and compared against all six local models on the exact same $20$ tasks. \texttt{gemini-flash} achieved $70.0\%$ mean Pass@1 ($75.0\%$ Neutral, $70.0\%$ XML, $65.0\%$ Polite) versus $0.0\%$--$23.3\%$ across the six local models---a $47$--$70$ pp capability gap on identical tasks, confirming that model-tier differences are not an artifact of dataset calibration.

\paragraph{4. Independent Non-Math SE Coding Benchmark ($15$ Tasks, $135$ Runs).}
To test external validity beyond math-derived tasks, we constructed and unit-test-verified $15$ standalone non-math software-engineering Python tasks ($75$ assertions: \texttt{.env}/CLI/URL/log/Markdown parsing, semver comparison, interval merging, LRU cache, bracket balancing, JSON flattening, RLE compression, Unix path normalization, topological task scheduling, sliding-window rate limiting, and record deduplication) evaluated across \texttt{gemini-flash}, \texttt{llama3.1:8b}, and \texttt{qwen2.5:7b} ($\times 3$ policies). XML matched or outperformed Plaintext on every model (\texttt{gemini-flash} $93.3\%$ vs.\ $93.3\%$; \texttt{llama3.1:8b} $53.3\%$ vs.\ $40.0\%$; \texttt{qwen2.5:7b} $53.3\%$ vs.\ $46.7\%$), while Polite phrasing inflated tokens and latency on all three models (up to $+24.5\%$ tokens and $+38.2\%$ latency on \texttt{gemini-flash} with $0.0$ pp accuracy gain), replicating our core operational-overhead finding outside mathematics.

\section{Discussion}\label{sec:discussion}
\subsection{Main Interpretation}\label{main-interpretation}

We observed that prompt framing is not a universal accuracy tool, but it is a real engineering variable: model ability dominates correctness whenever models differ in capability, while prompt choices change output form, parseability, latency, cost, and sometimes correctness. Prompts must therefore be evaluated as part of the overall LLM system, not adopted as universal rules.

Prompt effects are also better understood as interface effects between the model, parser, test harness, and evaluator than as changes in model competence. Polite code prompts increase explanatory language, which makes extraction harder, while structured wrappers reduce verbosity. In local math tasks, the compactness induced by some prompt formats is associated with lower answer accuracy. The practical lesson is that prompts, parsers, validators, and task definitions form one measurement system.

\subsection{Implications for LLM-Based Coding Systems}\label{implications-for-llm-based-coding-systems}

For LLM-based coding tools, prompts are part of the model--evaluator interface. Reliable systems need more than prompt design: they need explicit output contracts, parsers that tolerate common valid forms, and automated validation. Workflows should separate any reasoning or explanation from the final artifact, produce a clean code artifact at the end, and avoid relying on social framing cues such as politeness or role labels for correctness.

The API Python results show the value of structured wrappers in some coding workflows: raw pass-rate contrasts improve under JSON and XML while tokens and duration decrease, and the mixed-effects model supports both effects. This suggests that structured wrappers can improve compatibility with automated evaluation when they align with the task and model. The local math result shows that this pattern does not transfer uniformly across task settings.

The same principle applies to agentic workflows, where model outputs are consumed by downstream systems. A response that provides human-friendly explanations may be unusable to automated tools due to extra commentary or formatting inconsistencies, while machine-friendly output may be too terse for human debugging. Prompt policy must therefore be directed at the actual output consumer: the instruction should state the required artifact, the evaluator should accept common valid forms, and systems serving both human and automated consumers should separate these phases rather than asking one prompt to satisfy incompatible goals.

The end-to-end evaluations in Section~\ref{end-to-end-replay} turn that principle into measured workflow economics. While an idealized independent-redraw simulation suggests two retries could lift API Python from $67.8\%$ to $96.7\%$ (Table~\ref{tab:pipeline-replay}), our live $450$-workflow pipeline with real error-feedback retries (Table~\ref{tab:live-pipeline-measured}) shows that actual $T = 0$ API recovery is only $+4$--$6$ pp ($68$--$70\% \to 74$--$76\%$, with zero additional recovery on the second retry), whereas local $T = 0.2$ models gain $+8$--$14$ pp ($10$--$22\% \to 20$--$32\%$) at $2.5$--$2.7\times$ the calls and $27$--$45$\,s of LLM inference time (with parser/AST/test execution contributing $<1.5$\,s). Because parser-level failures account for $2.7\%$ of local-code failures and $20.4\%$ of API Python failures, retry and parser budgets should be sized to the dominant failure mode: resilient extraction and structured wrappers help when model competence is high, whereas low-accuracy local deployments benefit far more from upgrading the base model than from repeated retries.

\subsection{Why Code and Math Differ}\label{why-code-and-math-differ}

Programming and mathematical evaluations impose different requirements for output results. Code evaluation favors artifacts that are executable and compatible with the parser. Extra explanations become barriers by interfering with extraction. Structured wrappers help to signify the nature of the required output and reduce irrelevant text.

Mathematical evaluation, on the other hand, requires a correct answer derived from the reasoning process. The output contract asks for just providing the answer; therefore, extraction is easier, although very compact framing may coincide with less successful answer derivation. In this experiment, JSON framing lowers the number of local math tokens and reduces generation duration, while the pass rate drops. This does not prove that visible reasoning is required; it only shows that format constraints interact with task types in ways that affect downstream accuracy.

\subsection{Reconsidering Prompt-Engineering Folklore}\label{reconsidering-prompt-engineering-folklore}

The findings give a mixed view of prompt-engineering folklore. Politeness and role framing do not provide universal correctness gains. Polite prompts increase verbosity and duration and do not help local code, although they show a small positive effect in local math under the mixed-effects analysis. Role framing is also mixed: it improves local math in the pooled analysis ($+2.1$ pp AME), but reduces local code ($-1.0$ pp) and API math ($-2.8$ pp). JSON/XML wrappers are useful in API Python, especially XML, while JSON is associated with lower accuracy in local math. The result is therefore not that any one prompt style is always good or bad; it is that these choices depend on the model, task, and evaluator.

Therefore, the implications of this result are highly conditional. The effect of a prompt modification varies with the model's default settings, task type (coding vs. math), and evaluation criteria (pass rate vs. behavior). The conditional nature of these results also explains the discrepancy in the advice found in existing literature: different prompt versions bundle multiple interventions at once (politeness lengthens prompts, roles add persona text, and JSON alters both structure and separators), and different models have different default verbosity levels and wrapper sensitivities. Consequently, each bundled element can pull the model's behavior in a different direction. Factorial research isolates these components, though it cannot guarantee that any single component will be universally effective. This is why prompt practices should always be verified locally, described clearly, and measured in terms of both correctness and behavior.

\subsection{Consolidated Practical Guidelines and Cost--Benefit Proportionality}\label{practical-prompt-policy-recommendation}

To consolidate all stakeholder recommendations in one place without repeating them across sections, Table~\ref{tab:actionable-guidelines} maps our empirical findings to concrete starting practices, explicitly distinguishing \textbf{directly measured evidence} from \textbf{proposed practices / hypotheses} requiring local validation. Because our API stratum includes two models and our primary coding benchmarks are math-derived (supplemented by $15$ non-math SE tasks), these guidelines should be treated as \emph{evidence-grounded starting defaults} rather than universal prompting laws.

\begin{table}[htbp]
\centering
\caption{Consolidated stakeholder guidelines grounded in measured evidence, explicitly separating directly supported findings from proposed practices / hypotheses requiring local validation.}\label{tab:actionable-guidelines}
\footnotesize
\setlength{\tabcolsep}{3.5pt}
\begin{tabular}{@{}p{0.15\linewidth}p{0.26\linewidth}p{0.31\linewidth}p{0.23\linewidth}@{}}
\toprule
\textbf{Stakeholder} & \textbf{Directly Measured Evidence} & \textbf{Evidence-Grounded Starting Practice} & \textbf{Status / Scope Caveat} \\
\midrule
\textbf{1. Practitioners \& Developers} & Model identity explains $7.2$--$41.0\%$ of Pass@1 variance vs.\ $\le 0.41\%$ for prompts in 3 of 4 strata; shared-task gap is $47$--$70$ pp across 7 models. & Invest first in base-model selection; treat zero-shot framing as secondary for correctness. & \textbf{Supported} across all 4 strata and shared-task check. \\
\addlinespace[3pt]
 & Polite tone yields no code Pass@1 gain (and $-2.56$ pp in $k=3$ local code) while inflating tokens ($+12.4\%$ primary; up to $+24.5\%$ on SE benchmark) and latency ($+10.8\%$; $+38.2\%$ SE). Aggressive tone is null. & Start with neutral, direct imperative instructions; omit polite preambles and aggressive pressure cues in automated templates. & \textbf{Supported} across primary, $k=3$, and non-math SE benchmarks. \\
\addlinespace[3pt]
 & XML/JSON improve API Python Pass@1 ($+5.2$/$+3.4$ pp), API Python compliance ($+2.22$/$+2.88$ pp), and local code compliance ($+6.98$/$+4.27$ pp); XML matches or beats Plaintext on all 3 SE-benchmark models. Conversely, JSON reduces local math Pass@1 ($-4.1$ pp primary; $-6.78$ pp in $k=3$). & Use XML (or JSON) input delimiters as a starting default for automated code tasks, but keep local math prompts in plaintext unless validated. & \textbf{Supported} in studied models; verify locally before transferring across new model families (only 2 API models tested). \\
\midrule
\textbf{2. Tool, Agent \& Harness Builders} & Strict parser acceptance ($97.6\%$ local, $93.4\%$ API) exceeds $<15\%$ extra-text compliance ($72.0\%$ local); live retries recover $+4$--$6$ pp on API ($T=0$) and $+8$--$14$ pp locally at $2.5$--$2.7\times$ calls. & Pair explicit output contracts with resilient multi-stage extractors (fence regex $\to$ \texttt{ast.parse} $\to$ tests) and cap retries at 1--2 attempts. & \textbf{Supported} by compliance audit and 450-workflow live pipeline. \\
\addlinespace[3pt]
 & Local math JSON accuracy drops co-occur with $-19.8\%$ token compaction (associational, not proven causal). & Test two-field schemas (\texttt{<scratchpad>} + \texttt{<final\_answer>}) when structured output is required for multi-step math. & \textbf{Proposed practice / hypothesis} (not directly tested here). \\
\midrule
\textbf{3. Empirical SE \& AI Researchers} & Wrapper effects reverse sign across code and math; $k=3$ flip rates reach $23.7\%$ even on $T=0$ APIs; neutral paraphrases shift 30-task pass rates by $0$--$6.7$ pp. & Report exact templates, model snapshots, $k \ge 3$ run-to-run variability, leave-one-model-out checks, and behavioral telemetry alongside Pass@1. & \textbf{Supported} methodological recommendation. \\
\bottomrule
\end{tabular}
\end{table}

\paragraph{Balancing Expected Effect Sizes Against Benchmarking Cost (Proportionality Criteria).}
Because 12 of the 20 mixed-effects correctness contrasts in Table~\ref{tab:mixed-effects-pass} are compatible with zero effect and the 8 supported accuracy shifts are modest ($1.0$--$5.2$ pp), exhaustive factorial prompt benchmarking is disproportionate for many projects. Teams can scale evaluation effort across three tiers:
\begin{enumerate}
\tightlist
\item \textbf{Tier 1 --- Low-Volume or Human-in-the-Loop Workflows ($<10^3$ calls/month):} Adopt the neutral starting defaults in Table~\ref{tab:actionable-guidelines} and invest budget in base-model selection rather than prompt-matrix sweeps.
\item \textbf{Tier 2 --- High-Volume Automated Services ($>10^4$ calls/month):} Run a lightweight 2--3 variant behavioral smoke test (e.g., Neutral Plaintext vs.\ XML on 30--50 representative tasks) to capture $10$--$20\%$ token and latency savings.
\item \textbf{Tier 3 --- Unattended Execution Pipelines and Published Benchmarks:} Conduct multi-prompt validation with $k \ge 3$ repeats and version-control prompt templates, model identifiers, and parser rules together.
\end{enumerate}

\section{Threats to Validity}\label{threats-to-validity}

\subsection{Internal Validity}\label{internal-validity}

The crossed design maintains internal validity because every base task appears under every prompt variant. This reduces confounding between task difficulty and prompt condition. However, prompt variants also differ in length and framing. In practice, polite, aggressive, role, and structured-wrapper prompts can add different token counts. In this context, some effects may reflect prompt length or instruction density rather than tone or format alone.

Issues with hardware variability, API variability, and timeouts should also be taken into consideration. Empty responses are retained in order to preserve the factorial design, but their causes are not fully observable from the stored outputs. Table~\ref{tab:execution-provenance} and Section~\ref{execution-environment-and-run-protocol} report the complete local hardware platform, runtime parameters, run dates, and retry policies directly, with full runner logs archived in the replication package.

The crossed design limits run-order and environmental confounds, but does not eliminate them. Every base task appears under every prompt variant to avoid comparing easy tasks in one condition with hard tasks in another. However, local runtime load and hosted API serving conditions may still vary during different calls. Duration results should therefore be read as observed system latency in the executed environment, not as a hardware-independent property of a prompt. Correctness and compliance are less sensitive to timing noise, but timeouts and empty responses show that runtime behavior remains part of the evaluated system.

\subsection{Construct Validity}\label{construct-validity}

Pass@1 and test pass rate do not reflect every aspect of code quality. Although a program has passed the available tests, it still may be brittle, inefficient, or poorly styled. Conversely, a response may contain a correct solution but still fail in terms of extraction. Such limitations are not only flaws; they reflect the engineering realities of an automated LLM pipeline.

Math correctness depends on the deterministic four-step extraction, \LaTeX{}/symbolic canonicalization, and AST evaluation pipeline ($|\hat{y}_i - y_i| \le 10^{-6}$, Table~\ref{tab:outcome-operationalization}). Similarly, code-only compliance ($L_{\text{extra}} / L_{\text{resp}} < 0.15$) and math result-only compliance ($0 < L_{\text{resp}} < 50$ characters) are deterministic operational indicators of visible-channel conciseness. Validating the $15\%$ code-only rule directly against strict \texttt{ast.parse} acceptance across all $7{,}200$ primary Python runs shows $99.94\%$ agreement on API Python ($93.39\%$ code-only vs.\ $93.44\%$ parser-accepted after correcting the missing top-level \texttt{extracted\_code} key) and $99.43\%$ precision / $73.39\%$ recall on local Python ($25.96\%$ of local outputs contain $\ge 15\%$ surrounding prose yet still yield a valid fenced Python block, with a much lower Pass@1 of $2.14\%$ vs.\ $11.88\%$ when $<15\%$). Thus, $<15\%$ extra text measures zero-chatter adherence rather than raw parser survivability, and both metrics are reported separately. For math, cross-checking the primary logs against our full-raw $k = 3$ runs confirms that non-reasoning models logged genuine raw text, whereas reasoning models (\texttt{nemotron-3-nano-4b}, \texttt{deepseek-v4-flash}) emit hidden thinking tokens in a separate channel counted inside \texttt{completion\_tokens}; math compliance therefore reflects visible-channel brevity rather than total reasoning verbosity.

Moreover, task definitions and prompt template operationalization influence construct validity. Coding tasks focus on functional accuracy according to the provided tests, which does not necessarily guarantee maintainability, security, style, or generalization. For math tasks, the focus is on the extracted final answer, not the quality or faithfulness of an explanation. Similarly, tone (neutral, polite, aggressive) and role framing are operationalized through deterministic, rule-based lexical templates (Table~\ref{tab:prompt-templates-exact}) rather than free-form crowdsourced paraphrases or an external sentiment classifier. While our four-author blind audit confirmed $100\%$ inter-rater agreement on tone classification and semantic-constraint equivalence across all six domain--tone templates—and our 180-run neutral paraphrase and length-matched experiment (Section~\ref{empirical-validation-suite}) directly bounds the effect of neutral rewording and $+58$-token prompt-length expansion ($0.0$--$6.7$ pp on 30-task subsets)—each tone level still represents one canonical realization of politeness or urgency (with a natural length difference of roughly 30--64 input tokens over the neutral baseline). These choices are appropriate for an automated prompt-framing experiment, but they should not be confused with full measures of programming expertise, mathematical reasoning, or the entire pragmatic spectrum of human politeness. The paper therefore emphasizes observable outcomes rather than broad claims about model ability.

\subsection{Statistical Conclusion Validity}\label{statistical-conclusion-validity}

The experiment includes repeated observations by base task and model. The main inferential analysis therefore uses mixed-effects logistic models with a random intercept per base task, which model this dependence structure explicitly, supported by clustered linear-probability robustness checks. Simple pairwise tests and descriptive percentage-point contrasts are retained only as the interpretability layer; they do not fully model the dependence structure on their own. Furthermore, many pairwise comparisons are exploratory, and uncorrected p-values should not be overinterpreted. Model-specific contrasts require false-discovery-rate correction before they can support strong inferential claims.

The analysis uses different levels of evidence. Descriptive contrasts are useful for direction and effect size. Local model-specific contrasts are corrected with the Benjamini-Hochberg method because of the large number of comparisons made. Mixed-effects logistic models provide the main inference for pass/fail, with clustered linear-probability models as a fixed-effect robustness check. Results such as API Python XML ($\text{AME} = +5.2$ pp $[+2.3, +8.2]$) and local math JSON ($\text{AME} = -4.1$ pp $[-5.4, -2.7]$) are more robust because they combine non-trivial effect sizes with directional stability across the mixed model, the clustered check, and the dominant-model sensitivity panels; marginal findings such as API Python JSON ($+3.4$ pp $[+0.4, +6.3]$), which loses significance when DeepSeek V4 Flash is removed, should be read as suggestive rather than definitive.

One remaining limitation is estimation method and statistical power. First, the mixed-effects logistic models were fitted using variational Bayes rather than full maximum-likelihood or Markov-chain Monte Carlo estimation, which might result in somewhat understated posterior uncertainty; however, because the mixed-effects results show close agreement with the clustered-LPM robustness check, the estimation method does not alter any substantive conclusions. Second, the 50-task design has a finite resolution, which Section~\ref{sample-size-and-power} reports explicitly: at $80\%$ power the task-clustered contrasts detect $1.3$--$2.2$ pp in local code, $1.9$--$3.1$ pp in local math, $3.9$--$6.0$ pp in API Python, and $3.5$--$4.4$ pp in API math, with $90\%$-power thresholds about $16\%$ wider (Table~\ref{tab:power-analysis}). Contrasts below those floors---including any sub-percentage-point prompt effect and most API contrasts of $1$--$2$ pp---cannot be distinguished from zero here, so the corresponding null results are reported as inconclusive rather than as evidence of equivalence. Third, while the crossed $3 \times 2 \times 3$ design ensures orthogonal marginal balance for main effects ($1,800$--$2,700$ observations per factor level in local strata and $600$--$900$ in API strata), estimating two-way factor-by-factor interactions partitions observations into smaller sub-cells ($600$--$900$ locally and $200$--$300$ in API strata, with clustered standard errors of $0.74$--$1.69$ pp locally and $2.81$--$5.27$ pp in API strata, Table~\ref{tab:two-way-interactions}). Consequently, while our data rule out large synergistic or antagonistic two-way interactions ($>3$--$4$ pp locally or $>8$--$10$ pp in API strata) and identify significant $\text{Format} \times \text{Model}$ heterogeneity in math, detecting subtle higher-order interactions ($<2$ pp) would require a larger base-task sample or multi-run replications per cell; the task-count extrapolation in Table~\ref{tab:power-analysis} indicates that resolving a $2$ pp contrast at $80\%$ power would need roughly $21$--$63$ tasks locally and $149$--$449$ tasks in the API strata.

\subsection{Single-Run Primary Matrix and Empirical $k = 3$ Run-to-Run Variability}\label{single-run-noise}

In the primary $14{,}400$-run factorial matrix, each (prompt, task, model) cell was executed once: those $14{,}400$ runs represent $8$ models $\times$ $2$ task kinds $\times$ $50$ base tasks $\times$ $18$ treatment variants, \emph{not} $14{,}400$ independent replications. Because every prompt contrast is estimated within base task, the task-clustered standard errors of the primary matrix ($0.46$--$0.80$ pp local code, $0.69$--$1.11$ pp local math, $1.40$--$2.14$ pp API Python, $1.23$--$1.56$ pp API math) already lie at $0.43$--$0.83\times$ the single-run Bernoulli noise floor.

To directly measure within-cell nondeterminism rather than relying on theoretical bounds alone, we executed $k = 3$ independent repeated generations across all $8$ study models, all $50$ base tasks, and the $6$ core prompt conditions ($14{,}398 / 14{,}400$ valid runs; Section~\ref{empirical-validation-suite} and Table~\ref{tab:k3-repeated-variability}). This replication establishes three empirical facts: (i) hosted APIs at temperature $0$ are not uniformly deterministic (\texttt{gemini-flash} shows $0.00\%$ cell flips across $600$ cells, whereas \texttt{deepseek-v4-flash} via AI/ML API shows $23.67\%$ flips in Python and $15.00\%$ in math, alongside a provider-host shift in Python baseline from $66.7\%$ on Pioneer AI to $39.67\%$ on AI/ML API); (ii) local models at $T = 0.2$ show $2.33\%$--$28.33\%$ cell-level flips, yet aggregate run-to-run SD across repetitions remains tiny ($0.19$--$1.33$ pp locally; $0.00$--$2.08$ pp on APIs); and (iii) repeating $k = 3$ tightens the $80\%$-power MDE for \emph{single-cell pairwise contrasts against the neutral-plaintext baseline} ($n = 600 \to 1{,}800$ per contrast, distinct from the 18-variant marginal main-effect MDEs of $1.29$--$3.12$ pp in Table~\ref{tab:power-analysis}) by $15$--$40\%$ (from $3.45$--$5.02$ pp at $k = 1$ to $3.05$--$3.86$ pp in local code and $3.00$--$3.94$ pp in local math) while replicating the core negative effects of local code politeness ($-2.56$ pp, $95\%$ CI $[-4.85, -0.27]$) and local math JSON wrapping ($-6.78$ pp, $[-9.53, -4.03]$). Sub-threshold contrasts ($1$--$2$ pp) remain inconclusive and are treated as exploratory rather than evidence of equivalence.

\subsection{External Validity}\label{external-validity}

In the given set of tasks, there are 50 code as well as 50 mathematics tasks per model group. The tasks are domain-specific, in the sense that local tasks are derived from the national exam in Vietnam while the API tasks are from MathNet competitions. Since our coding tasks are constructed as self-contained algorithmic functions derived from mathematical problems, the study evaluates controlled single-function synthesis rather than repository-level software development \citep{jimenez2024swebench}. Real-world software engineering involves multi-file module dependencies, third-party APIs, ambiguous user requirements, stateful I/O, and iterative debugging workflows where prompt-framing sensitivity may differ substantially from standalone algorithmic synthesis. Furthermore, within the same algorithmic problem family, our complexity-tier disaggregation (Section~\ref{tasks}) shows that prompt-framing effects do not scale linearly with task difficulty: framing sensitivity is concentrated in solvable easy-to-medium tasks where the model already possesses sufficient domain competence and where output compliance or response compaction determines success, whereas on hard tasks that exceed a model's parameter capacity, surface prompt framing yields negligible accuracy gains. Consequently, our findings generalize most directly to automated single-shot code-generation and quantitative-evaluation pipelines on tasks near a model's competence boundary, and should not be extrapolated uncritically to complex multi-file software repositories or tasks far beyond a model's capability ceiling.

The size of the task sample additionally bounds the granularity of generalization we claim. The 50 tasks per benchmark span the ranges documented in Table~\ref{tab:task-dataset-characterization} (problem statements of $2$--$211$ words, reference or worked solutions of $2$--$60$ lines, and four topic families per benchmark), and Section~\ref{sample-size-and-power} shows that this sample resolves prompt contrasts of roughly $1.3$--$3.1$ pp at $80\%$ power in the local strata. A replication on 50 freshly drawn problems from the same family would therefore attach a $95\%$ sampling interval of about $\pm 0.9$ to $\pm 1.6$ pp to a local-code contrast and $\pm 1.4$ to $\pm 2.2$ pp to a local-math contrast. We consequently interpret our estimates as evidence about the sign, direction, and order of magnitude of prompt-framing effects on this problem population---and about the cross-regime reversals---rather than as task-level predictions for arbitrary new problems of the same family, particularly harder variants where the complexity-tier analysis shows that framing sensitivity itself changes.

The evaluated models comprise six local open-weight models in the single-GPU 4B--9B parameter tier (spanning hybrid Mamba-2/Transformer, code-specialized, and general GQA/SWA Transformer architectures) alongside two commercial cloud API models. While this design tests across-family heterogeneity within a fixed 16\,GiB local hardware envelope and contrasts it with hosted cloud endpoints, it does not cover intermediate open-weight parameter scales (e.g., 14B, 32B, or 70B models requiring multi-GPU or high-VRAM server hardware) or full-precision (FP16/BF16) local deployments. Moreover, commercial API model behavior may evolve across provider updates, making the exact model identifiers and execution timestamps reported in Table~\ref{tab:execution-provenance} essential for reproducibility.

This research does not claim that the selected prompt factors exhaust the prompt-design space. There may be many other prompt-related factors like few-shot examples, chain-of-thought instructions, explicit schemas, constrained decoding, tool-use instructions, or multi-turn repair prompts. All of these factors may interact with tone, role, and wrapper format in different ways.

The external validity of the experiment is also restricted due to the one-shot evaluation setting. Real coding assistants often operate over multiple turns, inspect test failures, edit files, or call tools. Such feedback loops may reduce some prompt-framing effects since the system can repair malformed outputs, or they may introduce new effects because the prompt also controls how the model responds to tool feedback. Our research can be understood best as evidence about initial generation under controlled prompt variation rather than as a full model of long-running software agents.

\subsection{Evaluation Validity, Oracle Audit, and Replication Package Verification}\label{evaluation-validity-and-audit-evidence}

We conducted three independent post-hoc audits to verify data, oracle, and artifact integrity:
\begin{itemize}
\tightlist
\item \textbf{Dominant-Model Audit (Nemotron Math):} Auditing the \texttt{nemotron-3-nano-4b} local math file confirmed complete row counts, task and variant balance, zero duplicates, and independent pass recomputation matching the stored $91.2\%$ pass rate (replicated at $90.56\%$ across $k = 3$ full-raw repeats). Comparing primary and $k = 3$ full-raw logs confirms that Nemotron and DeepSeek V4 Flash emit hidden thinking tokens inside \texttt{completion\_tokens} while returning concise answers in \texttt{response\_text}.
\item \textbf{Independent 200-Task Specification and Test-Oracle Audit:} Executing and statically inspecting all $200$ tasks ($100$ Python tasks / $700$ assertions + $100$ Math tasks; \texttt{audit\_task\_oracles.py}) confirmed that $100/100$ reference Python solutions pass all $7$ assertions with zero duplicate tests, and $100/100$ math expected answers parse cleanly. API Code ($50/50$) and API Math ($50/50$) are $100\%$ clean; Local Math is $49/50$ clean (task~48 has a decimal answer, \texttt{961.6}, without an explicit rounding note); and Local Code surfaces three transparent oracle nuances: (i) all $50$ local code tasks communicate the exact target function name and signature via the two in-prompt example assertions (\texttt{test\_list}) rather than repeating the identifier inside the prose problem statement (\texttt{FUNC\_NAME\_NOT\_IN\_TEXT = 50/50}), (ii) $3$ tasks translate fixed-constant exam items with zero parameters (tasks~13, 36, 46; task~46 constant-return), and (iii) $16$ tasks use exact float equality (\texttt{==}) on rounded floats rather than tolerance bounds, requiring strict adherence to the prompt's rounding instructions.
\item \textbf{Recursive Replication Package Audit:} Recursively scanning the public replication repository ($60$ directories, $1{,}364$ files) confirmed $0$ legacy thesis directories or files, $0$ secret/\texttt{.env} files, $0$ redundant intermediate backup/scratch trees, and full alignment with the Data Availability Statement.
\end{itemize}

\section{Conclusion}\label{conclusion}

This study investigated how three selected zero-shot prompt-framing dimensions (tone, role framing, and plaintext/JSON/XML input wrappers) affect accuracy and operational behavior across Python code generation and mathematical reasoning, combining a $14{,}400$-run primary factorial matrix ($8$ models $\times$ $50$ base tasks $\times$ $18$ variants) with $14{,}398$ independent $k = 3$ repeated runs, $450$ live end-to-end pipeline workflows with real retries, $180$ neutral paraphrase/length controls, a $20$-task shared-regime comparison, and a $15$-task non-math SE coding benchmark.

Across these evaluations, a consistent hierarchy emerges: wherever candidate models differ in capability, model selection dominates Pass@1 correctness ($7.2\%$--$41.0\%$ of variance vs.\ $\le 0.41\%$ for prompts; $47$--$70$ pp gap on shared tasks), whereas zero-shot prompt framing produces modest, selective Pass@1 shifts ($1.0$--$5.2$ pp across 8 supported contrasts; 12 of 20 intervals compatible with zero) alongside systematic operational shifts in tokens, latency, and parser compliance. Polite phrasing adds token and latency overhead without improving code accuracy—replicated across primary runs, $k = 3$ repetitions ($-2.56$ pp in local code), and the non-math SE benchmark—while structured wrappers improve code-interface compatibility yet degrade local mathematical reasoning under JSON ($-4.1$ pp primary; $-6.78$ pp in $k = 3$). In live end-to-end execution, real error-feedback retries recover $+4$--$6$ pp on API Python ($T = 0$) and $+8$--$14$ pp locally, substantially below idealized independent-retry simulations. Rather than prescribing universal prompting laws, we consolidate evidence-grounded starting defaults, hypotheses requiring local validation, and three cost--benefit proportionality tiers in Section~\ref{practical-prompt-policy-recommendation} (Table~\ref{tab:actionable-guidelines}). Future work should extend this framework to repository-level software engineering \citep{jimenez2024swebench}, intermediate open-weight parameter scales (14B--70B), human-rated tone equivalence panels, and multi-turn agentic tool workflows.


\section{Data Availability Statement}\label{data-availability-statement}

The complete, self-contained replication package—comprising the 18-variant factorial prompt templates, local and cloud API evaluation runners, the consolidated $14{,}400$-run primary JSONL dataset, the $14{,}398$-run $k = 3$ repeated-generation dataset, the $450$-workflow live end-to-end retry pipeline logs, the $180$ neutral paraphrase/length-matched controls, the $20$-task shared-regime comparison, the $15$-task non-math SE coding benchmark, statistical analysis scripts (mixed-effects, clustered LPM, two-way interactions, BH-FDR, behavioral adjustment, variance decomposition, power analysis, and dominant-model sensitivity), and the 200-task oracle, compliance, and recursive repository audit logs, with all legacy thesis draft directories removed—is publicly available at: \url{https://github.com/tycoi2005/prompt-framing-experiment}.

\section{Funding Statement}\label{funding-statement}

This work received no specific grant from any funding agency, commercial, or not-for-profit sectors.

\section{Competing Interests}\label{competing-interests}

The authors declare no competing interests.

\section{AI Tools Declaration}\label{ai-tools-declaration}

AI tools were used for editing support, code assistance, and drafting assistance. The experimental outputs were generated by the evaluated models and were not manually altered for the reported analyses.

\section{Ethics Statement}\label{ethics-statement}

This study does not involve human participants. It analyzes model
outputs on coding and mathematical reasoning tasks. Public replication
should respect task-data licensing and API logging policies, including
releasing reconstruction scripts or task identifiers instead of copied
task text when redistribution rights are unclear.

\appendix
\section{Supplementary Visualizations}\label{app:model-pass-rate-figure}

Figure~\ref{fig:model_performance_by_task} provides a visual bar-chart comparison of the model-level Pass@1 pass rates across the four task-regime strata reported in Table~\ref{tab:model-level-performance} (Section~\ref{model-ability-dominates-correctness}).

\begin{center}
\begin{minipage}{0.90\linewidth}
\centering
\includegraphics[width=0.85\linewidth]{model_performance_by_task.png}
\captionof{figure}{Model pass rates by task kind across the local and cloud API strata (complementing Table~\ref{tab:model-level-performance}). The formal comparison of model and prompt effects is reported in Table~\ref{tab:model-vs-prompt-decomposition}.}
\label{fig:model_performance_by_task}
\end{minipage}
\end{center}


\printcredits

%\bibliographystyle{cas-model2-names}
\bibliographystyle{elsarticle-num}
\bibliography{literature}

\end{document}
